% The Riemann Problem in Special-Relativistic Magnetohydrodynamics
% Lecture notes.  Build:  latexmk -pdf rmhd_riemann_course.tex
\documentclass[11pt,a4paper,oneside,openany]{book}

\usepackage[T1]{fontenc}
\usepackage[utf8]{inputenc}
\usepackage{lmodern}
\usepackage{amsmath,amssymb,amsthm,mathtools,bm}
\usepackage{stmaryrd}
\usepackage[a4paper,top=2.6cm,bottom=2.8cm,left=2.7cm,right=2.7cm,headheight=14pt]{geometry}
\usepackage{microtype}
\usepackage{booktabs,array,tabularx,multirow}
\usepackage{enumitem}
\usepackage{xcolor}
\usepackage{tikz}
\usetikzlibrary{arrows.meta,calc,decorations.pathreplacing,patterns}
\usepackage{pgfplots}
\pgfplotsset{compat=1.18}
\usepgfplotslibrary{polar}
\usepackage[most]{tcolorbox}
\usepackage{caption}
\captionsetup{font=small,labelfont=bf}
\usepackage[authoryear,round]{natbib}
\usepackage{url}
\usepackage[colorlinks=true,linkcolor=blue!55!black,citecolor=green!35!black,
            urlcolor=blue!55!black]{hyperref}
% citation commands used in the text
\newcommand{\textcite}{\citet}
\newcommand{\parencite}{\citep}
% links used in the bibliography
\newcommand{\arxiv}[1]{\href{https://arxiv.org/abs/#1}{arXiv:#1}}
\newcommand{\doi}[1]{\href{https://doi.org/#1}{doi:#1}}
\usepackage[capitalise,nameinlink]{cleveref}

% ---------------------------------------------------------------- theorem-like
\theoremstyle{definition}
\newtheorem{definition}{Definition}[chapter]
\newtheorem{example}{Example}[chapter]
\newtheorem{exercise}{Exercise}[chapter]
\theoremstyle{plain}
\newtheorem{proposition}{Proposition}[chapter]
\newtheorem{lemma}{Lemma}[chapter]
\theoremstyle{remark}
\newtheorem{remark}{Remark}[chapter]
\crefname{exercise}{Exercise}{Exercises}
\crefname{proposition}{Proposition}{Propositions}
\crefname{definition}{Definition}{Definitions}
\crefname{remark}{Remark}{Remarks}
\crefname{example}{Example}{Examples}

\newenvironment{solution}{\par\smallskip\noindent\small\textit{Solution.}\ }{\par\medskip}

% boxes
\newtcolorbox{codebox}[1][]{colback=gray!6,colframe=gray!55,boxrule=0.5pt,arc=1.5pt,
  left=4pt,right=4pt,top=3pt,bottom=3pt,fonttitle=\bfseries\small,
  title={In the code#1},breakable}
\newtcolorbox{keybox}[1][]{colback=blue!3,colframe=blue!45!black,boxrule=0.5pt,arc=1.5pt,
  left=4pt,right=4pt,top=3pt,bottom=3pt,fonttitle=\bfseries\small,title={#1},breakable}

% ---------------------------------------------------------------- notation
\newcommand{\jump}[1]{\llbracket #1\rrbracket}
\newcommand{\avg}[1]{(\!(#1)\!)}
\newcommand{\dd}{\mathrm{d}}
\newcommand{\vv}{\bm{v}}
\newcommand{\BB}{\bm{B}}
\newcommand{\EE}{\bm{E}}
\newcommand{\SSS}{\bm{S}}
\newcommand{\UU}{\bm{U}}
\newcommand{\FF}{\bm{F}}
\newcommand{\VV}{\bm{V}}
\newcommand{\calB}{\mathcal{B}}
\newcommand{\calA}{\mathcal{A}}
\newcommand{\calN}{\mathcal{N}}
\newcommand{\Ptot}{P}
\newcommand{\vB}{\eta}
\newcommand{\sgn}{\operatorname{sign}}
\newcommand{\code}[1]{\texttt{\detokenize{#1}}}
\newcommand{\file}[1]{\texttt{\detokenize{#1}}}
% flag for anything not verified against a source or the code
\newcommand{\unverified}[1]{\footnote{\textbf{Not independently verified:} #1}}

\hypersetup{pdftitle={The Riemann Problem in Special-Relativistic Magnetohydrodynamics},
            pdfauthor={Lecture notes}}

\begin{document}

% ======================================================================== title
\begin{titlepage}
\centering
\vspace*{3cm}
{\Huge\bfseries The Riemann Problem in\\[0.3em] Special-Relativistic Magnetohydrodynamics\par}
\vspace{1.2cm}
{\Large Lecture notes for a graduate course\par}
\vspace{0.6cm}
{\large with derivations, exercises, and notes on an exact solver in practice\par}
\vspace{2.5cm}
\begin{minipage}{0.82\textwidth}
\small
These notes develop the theory of the one-dimensional Riemann problem of ideal
special-relativistic magnetohydrodynamics (SRMHD) from the conservation laws up to
the exact and approximate solvers used in multidimensional codes. Derivations are
given step by step. The algebraic identities marked as such were checked
symbolically or numerically with the script \file{sympy_checks.py} that accompanies
these notes (\cref{app:verification}). Statements about the reference
implementation (the exact solver in \file{rmhd_final} and the finite-volume code in
\file{HLLD}) are restricted to what its source and its documentation record.
\end{minipage}
\vfill
{\small Version of 15 October 2026\par}
\end{titlepage}

\frontmatter
\tableofcontents

% ======================================================================== conventions
\chapter*{Conventions and notation}
\addcontentsline{toc}{chapter}{Conventions and notation}

\paragraph{Units and metric.} We use units with $c=1$ and the Minkowski metric
$g_{\mu\nu}=\operatorname{diag}(-1,1,1,1)$ in Cartesian coordinates $(t,x,y,z)$, so
covariant derivatives reduce to partial derivatives. Greek indices run over
$0,\dots,3$, Latin indices over $1,2,3$; repeated indices are summed. The magnetic field
is measured in Heaviside--Lorentz-type units in which the factor $\sqrt{4\pi}$ has been
absorbed into $\BB$: the magnetic pressure is $b^2/2$ and the Maxwell stress has no
$4\pi$. (A Gaussian field is $\sqrt{4\pi}$ times ours.) This is the convention of the
codes discussed here.

\paragraph{Fluid variables.} $\rho$ is the rest-mass density, $\epsilon$ the specific
internal energy, $p$ the gas pressure, $h=1+\epsilon+p/\rho$ the specific enthalpy of
the gas, $\vv$ the three-velocity, $W=(1-v^2)^{-1/2}$ the Lorentz factor and
$u^\mu=W(1,\vv)$ the four-velocity. The laboratory magnetic field is $\BB$, the comoving
field four-vector is $b^\mu$, and
\begin{equation*}
\vB\equiv\vv\cdot\BB,\qquad
b^2\equiv b^\mu b_\mu,\qquad
w\equiv\rho h+b^2,\qquad
\Ptot\equiv p+\tfrac12 b^2,\qquad
\sigma\equiv \frac{b^2}{\rho h}.
\end{equation*}
$w$ is the \emph{total} enthalpy density, $\Ptot$ the \emph{total} pressure and
$\sigma$ the magnetisation. Two warnings about other people's notation: the exact solver
of \textcite{GiacomazzoRezzolla2006} writes $h$ for the \emph{total} specific enthalpy
$h_g+b^2/\rho$ and $p$ for the total pressure $p_g+b^2/2$ (our $\Ptot$), with $h_g,p_g$ for
our $h,p$, and \textcite{Anton2010} write $\mathcal{E}$ for our
$w$. Where we quote their equations we translate them.

\paragraph{Speeds.} $c_s$ is the sound speed, $c_a^2=b^2/w$ the square of the Alfvén speed
in the fluid frame, and $\zeta^2=c_s^2+c_a^2-c_s^2c_a^2$. Characteristic speeds in the
laboratory are written $\lambda$, with subscripts e (entropy), A (Alfvén), s (slow), f
(fast). Shock speeds are $V_s$, with Lorentz factor $W_s$.

\paragraph{Jumps.} Across a discontinuity, ``a'' labels the state \emph{ahead} of the
front (into which it moves) and ``b'' the state \emph{behind}. Following
\textcite{GiacomazzoRezzolla2006} we write
\begin{equation*}
\jump{Q}\equiv Q_a-Q_b,\qquad \avg{Q}\equiv Q_a+Q_b .
\end{equation*}

\paragraph{Planar symmetry.} In the Riemann problem all quantities depend on $t$ and on
the coordinate $x$ normal to the initial discontinuity. ``Normal'' quantities carry the
index $x$ (or $n$), ``tangential'' ones the indices $y,z$ (or $t$, as in $B_t$).

\paragraph{Code references.} File names refer to the two repositories of the reference
implementation: \file{HLLD/} (finite-volume code; e.g.\ \file{src/physics/hlld.py}) and
\file{rmhd_final/} (exact solver; e.g.\ \file{rmhd/batched/planar5_b.py}). Its
documentation is cited by file name: \file{HLLD/docs/GUIDE.md} (``the guide'') and
\file{HLLD/docs/what_we_solve.md} (``the solver record'').

\mainmatter

% ======================================================================== chapter 1
\chapter{Introduction}\label{ch:intro}

\section{Hyperbolic conservation laws and the Riemann problem}

Ideal relativistic magnetohydrodynamics, like ideal gas dynamics, is a system of
conservation laws. In one space dimension it has the form
\begin{equation}\label{eq:cl1d}
\partial_t\UU+\partial_x\FF(\UU)=0,
\end{equation}
where $\UU\in\mathbb{R}^m$ collects densities of conserved quantities and $\FF(\UU)$
their fluxes. The system is \emph{hyperbolic} if the Jacobian
$\mathsf{A}(\UU)=\partial\FF/\partial\UU$ has real eigenvalues and a complete set of
eigenvectors for every admissible $\UU$; the eigenvalues are the speeds at which small
disturbances propagate. Solutions of \eqref{eq:cl1d} develop discontinuities in finite
time even from smooth data, so they are understood in the weak (integral) sense,
\begin{equation}\label{eq:integralform}
\int_{x_1}^{x_2}\UU(x,t_2)\,\dd x=\int_{x_1}^{x_2}\UU(x,t_1)\,\dd x
-\int_{t_1}^{t_2}\bigl[\FF(\UU(x_2,t))-\FF(\UU(x_1,t))\bigr]\dd t ,
\end{equation}
supplemented by admissibility conditions that select the physically relevant weak
solution (\cref{ch:discontinuities}).

\begin{definition}[Riemann problem]
The Riemann problem for \eqref{eq:cl1d} is the initial-value problem with piecewise
constant data
\begin{equation}\label{eq:rpdata}
\UU(x,0)=\begin{cases}\UU_L,& x<0,\\ \UU_R,& x>0.\end{cases}
\end{equation}
\end{definition}

Both \eqref{eq:cl1d} and \eqref{eq:rpdata} are invariant under the scaling
$x\to\alpha x$, $t\to\alpha t$ with $\alpha>0$. If the (admissible) solution is unique,
it must therefore be invariant too, $\UU(x,t)=\UU(\alpha x,\alpha t)$, i.e.\ it depends on
$x$ and $t$ only through the similarity variable
\begin{equation}
\xi\equiv x/t .
\end{equation}
The solution is a fan of \emph{waves} emanating from the origin and separating constant
states. A wave is either a discontinuity moving at constant speed (shock, contact,
rotational discontinuity) or a centred \emph{rarefaction}, a smooth self-similar fan.
For an $m$-component strictly hyperbolic system with genuinely nonlinear or linearly
degenerate fields, the classical theory \parencite{Lax1957} gives, for data close
enough together, exactly one wave per characteristic family and $m-1$ intermediate
constant states. Ideal MHD violates the hypotheses of that theorem at some states, which
is the source of most of the difficulty treated in these notes.

\section{Godunov's method}\label{sec:godunov}

Divide the line into cells $C_i=[x_{i-1/2},x_{i+1/2}]$ of width $\Delta x$ and store cell
averages $\UU_i^n=\Delta x^{-1}\int_{C_i}\UU(x,t^n)\,\dd x$. Applying
\eqref{eq:integralform} to $C_i\times[t^n,t^{n+1}]$ gives the \emph{exact} relation
\begin{equation}\label{eq:fvupdate}
\UU_i^{n+1}=\UU_i^{n}-\frac{\Delta t}{\Delta x}\bigl(\hat\FF_{i+1/2}-\hat\FF_{i-1/2}\bigr),
\qquad
\hat\FF_{i+1/2}=\frac{1}{\Delta t}\int_{t^n}^{t^{n+1}}\FF\bigl(\UU(x_{i+1/2},t)\bigr)\dd t .
\end{equation}
Any numerical method of this form is conservative: the sum of $\UU_i\Delta x$ changes
only through the boundary fluxes. \textcite{Godunov1959} closed \eqref{eq:fvupdate} by
replacing the unknown $\UU(x,t^n)$ with the piecewise-constant function equal to
$\UU_i^n$ in $C_i$. Every interface then carries a Riemann problem with data
$(\UU_i^n,\UU_{i+1}^n)$. As long as waves from neighbouring interfaces do not reach the
interface during the step, i.e.\ under the Courant--Friedrichs--Lewy condition
\begin{equation}
\frac{\Delta t}{\Delta x}\,\max|\lambda|\le 1,
\end{equation}
the solution \emph{at} $x_{i+1/2}$ is the constant $\UU^\ast_{i+1/2}=
\UU_{\mathrm{RP}}(\xi=0;\UU_i^n,\UU_{i+1}^n)$, and the flux integral is exact:
\begin{equation}\label{eq:godunovflux}
\hat\FF_{i+1/2}=\FF\bigl(\UU_{\mathrm{RP}}(0;\UU_i^n,\UU_{i+1}^n)\bigr).
\end{equation}
(The stricter bound $1/2$ guarantees that waves from neighbouring interfaces do not
interact inside a cell; it is needed if the new averages are computed by integrating the
juxtaposed Riemann solutions over each cell, but not for \eqref{eq:godunovflux}, which only
uses the value on the interface.)

\begin{figure}[t]
\centering
\begin{tikzpicture}[x=1.15cm,y=1.0cm,>=Stealth]
  % x-t plane with local Riemann fans
  \draw[->] (-0.3,2.4) -- (10.4,2.4) node[right] {$x$};
  \draw[->] (-0.3,2.4) -- (-0.3,4.4) node[above] {$t$};
  \draw[gray!60] (-0.3,4.0) -- (10.2,4.0) node[right,black,font=\small] {$t^{n+1}$};
  \node[left,font=\small] at (-0.3,2.4) {$t^n$};
  \foreach \i in {1,...,4}{
    \begin{scope}[shift={(2*\i,2.4)}]
      \draw[red!70!black] (0,0) -- (-0.75,1.6);
      \draw[red!70!black,dashed] (0,0) -- (0.15,1.6);
      \draw[red!70!black] (0,0) -- (0.6,1.6);
      \draw[very thick,blue!60!black] (0,0) -- (0,1.6);
    \end{scope}
  }
  % piecewise-constant data
  \foreach \i/\h in {0/1.6,1/1.25,2/0.7,3/0.85,4/0.45}{
    \draw[thick,blue!60!black] (2*\i,\h) -- (2*\i+2,\h);
    \node[font=\small] at (2*\i+1,\h+0.25) {$\UU_{\i}^n$};
  }
  \foreach \i in {0,...,5}{\draw[gray] (2*\i,0) -- (2*\i,2.2);}
  \draw[->] (-0.3,0) -- (10.4,0) node[right] {$x$};
\end{tikzpicture}
\caption{Godunov's method. Bottom: cell averages define piecewise-constant data at $t^n$.
Top: each interface carries a Riemann problem; the solution on the interface ($\xi=0$,
thick line) is constant in time and gives the flux \eqref{eq:godunovflux}.}
\label{fig:godunov}
\end{figure}

Godunov's scheme is first-order accurate, conservative, and \emph{upwind}: it uses
characteristic information to decide which side of the interface each piece of
information comes from. With the exact Riemann solution it inherits the admissibility
of the exact waves. Its accuracy is raised by \emph{reconstructing} a piecewise
polynomial from the averages and solving Riemann problems between the reconstructed
interface values (\cref{ch:multid}). If a conservative scheme converges as the grid is
refined, its limit is a weak solution \parencite[see the Lax--Wendroff theorem
in][]{LeVeque2002}; which weak solution is a separate question, and one that will
return in \cref{ch:nonunique}.

\subsection*{Why the Riemann problem is the building block}

Three properties make the Riemann problem the natural ``atom'' of a shock-capturing
scheme. (i)~It contains, exactly, the nonlinear wave structure that a scheme has to
respect at a discontinuity, including the correct jump conditions. (ii)~Its solution is
self-similar, so the flux at $\xi=0$ is a function of two states only, $\hat\FF=
\hat\FF(\UU_L,\UU_R)$, which can be evaluated face by face. (iii)~Every approximate
solver used in practice---HLL, HLLC, HLLD, Roe-type---is a model of the exact Riemann
fan that keeps some of its waves and discards others (\cref{ch:approx}); the exact
solution is the end point against which they are measured \parencite{Toro2009}.

\section{Why relativistic MHD}

Relativistic magnetised outflows and plasmas are central to high-energy astrophysics:
jets from active galactic nuclei and X-ray binaries, pulsar winds and their nebulae,
gamma-ray-burst outflows, and accretion flows and magnetospheres of compact objects.
Numerical models of these systems solve the equations of relativistic MHD with
Godunov-type schemes; reviews of the methods are given by \textcite{MartiMuller2015} and
in the textbook of \textcite{RezzollaZanotti2013}. Three features make the relativistic
magnetised case harder than Newtonian MHD:
\begin{itemize}[nosep]
\item the characteristic speeds are bounded by $c$ and crowd together as the flow or the
  field becomes relativistic, and they are roots of a quartic rather than of a
  biquadratic;
\item the conserved and primitive variables are related implicitly, so every flux
  evaluation needs a nonlinear inversion (\cref{sec:c2p});
\item in strongly magnetised regions the gas pressure is a small difference of large
  numbers, $p=\Ptot-b^2/2$, so small relative errors in the total pressure produce
  large relative errors (or negative values) in $p$.
\end{itemize}

\section{The exact solver in practice, and a roadmap}

The exact solution of the SRMHD Riemann problem was given by
\textcite{GiacomazzoRezzolla2006}, building on the special case of purely tangential
fields by \textcite{Romero2005} and on the relativistic hydrodynamic solvers of
\textcite{MartiMuller1994}, \textcite{Pons2000}, \textcite{RezzollaZanotti2001} and
\textcite{RezzollaZanottiPons2003}. Approximate solvers of the HLL family
\parencite{MignoneBodo2006,MignoneUglianoBodo2009} and linearised solvers
\parencite{Komissarov1999,Balsara2001,Anton2010} are what multidimensional codes use.
These notes are written for a reader who works on an exact SRMHD solver that is used as
the flux function of a two-dimensional constrained-transport code. That implementation
is described where it illustrates the theory (\cref{sec:student-impl,sec:failure-classes}),
always as recorded in its own documentation.

\paragraph{Roadmap.}
\Cref{ch:equations} derives the equations of ideal SRMHD and their conservative form.
\Cref{ch:characteristics} derives the characteristic speeds, their ordering and their
degeneracies. \Cref{ch:discontinuities,ch:rarefactions} treat the two kinds of
elementary waves, discontinuities and rarefactions. \Cref{ch:hydro} solves the
relativistic hydrodynamic Riemann problem completely, as a warm-up. \Cref{ch:srmhd}
assembles the seven-wave SRMHD problem and describes exact solution methods, including
the reference implementation. \Cref{ch:nonunique} discusses compound waves and
non-uniqueness. \Cref{ch:approx} derives the approximate solvers HLL, HLLC and HLLD and
reviews linearised solvers. \Cref{ch:multid} places the Riemann solver in a
multidimensional finite-volume code with constrained transport and lists the standard
tests. The appendices collect the algebra of $b^\mu$, the eigenstructure, a notation
table and a record of what was verified and how.

\section*{Exercises}
\addcontentsline{toc}{section}{Exercises}

\begin{exercise}\label{ex:selfsimilar}
Show that if $\UU(x,t)$ solves \eqref{eq:cl1d} with data \eqref{eq:rpdata}, then so does
$\UU_\alpha(x,t)=\UU(\alpha x,\alpha t)$ for every $\alpha>0$. Conclude that a unique
solution is a function of $x/t$.
\end{exercise}
\begin{solution}
$\partial_t\UU_\alpha=\alpha(\partial_t\UU)(\alpha x,\alpha t)$ and
$\partial_x\FF(\UU_\alpha)=\alpha(\partial_x\FF(\UU))(\alpha x,\alpha t)$, so
\eqref{eq:cl1d} holds; the data are scale invariant. Uniqueness gives
$\UU=\UU_\alpha$ for all $\alpha$; choose $\alpha=1/t$.
\end{solution}

\begin{exercise}
For linear advection $\partial_t u+a\partial_x u=0$ with $a>0$, show that Godunov's
method is the first-order upwind scheme $u_i^{n+1}=u_i^n-a\frac{\Delta t}{\Delta x}
(u_i^n-u_{i-1}^n)$, and that it is stable for $a\Delta t/\Delta x\le1$.
\end{exercise}

\begin{exercise}
Write the integral form \eqref{eq:integralform} on a rectangle that contains a single
discontinuity moving at constant speed $V_s$, and let the rectangle shrink onto the
discontinuity. Derive the Rankine--Hugoniot condition $\FF_a-\FF_b=V_s(\UU_a-\UU_b)$.
(This is used throughout \cref{ch:discontinuities}.)
\end{exercise}

% ======================================================================== chapter 2
\chapter{The equations of ideal SRMHD}\label{ch:equations}

\section{Fluid}

The rest-mass current and the energy--momentum tensor of a perfect fluid are
\begin{equation}
J^\mu=\rho u^\mu,\qquad
T^{\mu\nu}_{\mathrm{fl}}=\rho h\,u^\mu u^\nu+p\,g^{\mu\nu},
\end{equation}
with $u^\mu u_\mu=-1$ and $h=1+\epsilon+p/\rho$. In the fluid rest frame
$T^{00}_{\mathrm{fl}}=e\equiv\rho(1+\epsilon)$ is the energy density and
$T^{ij}_{\mathrm{fl}}=p\,\delta^{ij}$.

\section{Electromagnetic field and the ideal-MHD condition}\label{sec:bmu}

Let $F^{\mu\nu}$ be the Faraday tensor and ${}^\ast\!F^{\mu\nu}$ its dual. In the
laboratory frame we choose the components
\begin{equation}\label{eq:dualF}
{}^\ast\!F^{0i}=B^i,\qquad {}^\ast\!F^{ij}=-\varepsilon^{ijk}E_k ,
\end{equation}
with $\varepsilon^{ijk}$ the three-dimensional Levi-Civita symbol. With this choice the
homogeneous Maxwell equations $\partial_\mu{}^\ast\!F^{\mu\nu}=0$ read
$\nabla\cdot\BB=0$ ($\nu=0$) and $\partial_t\BB+\nabla\times\EE=0$ ($\nu=j$).

The comoving (fluid-frame) electric and magnetic fields are the four-vectors
\begin{equation}
e^\mu=F^{\mu\nu}u_\nu,\qquad b^\mu={}^\ast\!F^{\mu\nu}u_\nu .
\end{equation}
Both are orthogonal to $u^\mu$ by antisymmetry. A perfect conductor cannot sustain an
electric field in its rest frame; the \emph{ideal-MHD condition} is $e^\mu=0$, which in
the laboratory frame reads
\begin{equation}\label{eq:idealmhd}
\EE=-\vv\times\BB .
\end{equation}

\begin{proposition}[Comoving field]\label{prop:bmu}
Under \eqref{eq:idealmhd},
\begin{align}
b^0&=W(\vv\cdot\BB)=W\vB, \label{eq:b0}\\
b^i&=\frac{B^i}{W}+b^0v^i, \label{eq:bi}\\
b^2&=b^\mu b_\mu=\frac{B^2}{W^2}+(\vv\cdot\BB)^2, \label{eq:b2}
\end{align}
$u_\mu b^\mu=0$, and ${}^\ast\!F^{\mu\nu}=u^\mu b^\nu-u^\nu b^\mu$.
\end{proposition}

\begin{proof}
With $u_\nu=W(-1,\vv)$, \eqref{eq:dualF} gives
$b^0={}^\ast\!F^{0i}u_i=W\,\BB\cdot\vv$ and
$b^i={}^\ast\!F^{i0}u_0+{}^\ast\!F^{ij}u_j=WB^i-W\varepsilon^{ijk}v_jE_k
=W\,[\BB-\vv\times\EE]^i$.
With $\EE=-\vv\times\BB$,
$\vv\times\EE=-\vv\times(\vv\times\BB)=v^2\BB-\vv(\vv\cdot\BB)$, so
$b^i=W[(1-v^2)B^i+v^i\vB]=B^i/W+W\vB\,v^i$, which is \eqref{eq:bi}. Then
$u_\mu b^\mu=W(-b^0+\vv\cdot\bm{b})=W(-W\vB+\vB/W+W\vB v^2)=0$. Finally
$b^2=-(b^0)^2+|\bm b|^2=-W^2\vB^2+B^2/W^2+2\vB^2+W^2\vB^2v^2=B^2/W^2+\vB^2$. For the
last statement compute $u^0b^i-u^ib^0=W(B^i/W+W\vB v^i)-Wv^iW\vB=B^i={}^\ast\!F^{0i}$;
the spatial components follow in the same way.
\end{proof}

Inverting \eqref{eq:bi}, the laboratory field is $B^i=u^0b^i-u^ib^0=W(b^i-b^0v^i)$.
In the fluid rest frame $b^\mu=(0,\BB')$: $b^\mu$ \emph{is} the rest-frame magnetic
field, and $b^2$ its invariant strength. The \emph{magnetic pressure} is $b^2/2$.

\paragraph{Electromagnetic stress--energy.} In our units
$T^{\mu\nu}_{\mathrm{em}}=F^{\mu\alpha}F^{\nu}{}_{\alpha}-\frac14g^{\mu\nu}F_{\alpha\beta}F^{\alpha\beta}$.
Under the ideal-MHD condition,
\begin{equation}\label{eq:Tem}
T^{\mu\nu}_{\mathrm{em}}=b^2u^\mu u^\nu+\tfrac12b^2g^{\mu\nu}-b^\mu b^\nu .
\end{equation}
\emph{Proof.} In the fluid rest frame the field is purely magnetic, $\BB'=\bm b$, and the
Maxwell tensor has $T^{00}=b^2/2$, $T^{0i}=0$, $T^{ij}=\frac12b^2\delta^{ij}-b^ib^j$.
The right-hand side of \eqref{eq:Tem} evaluated with $u^\mu=(1,\bm 0)$ and
$b^\mu=(0,\bm b)$ gives the same components. Both sides are tensors, so they agree in
every frame. \qed

\section{Total energy--momentum tensor and the conservation laws}

Adding fluid and field,
\begin{equation}\label{eq:Tmunu}
\boxed{\;T^{\mu\nu}=(\rho h+b^2)\,u^\mu u^\nu+\Bigl(p+\frac{b^2}{2}\Bigr)g^{\mu\nu}-b^\mu b^\nu
=w\,u^\mu u^\nu+\Ptot\,g^{\mu\nu}-b^\mu b^\nu\;}
\end{equation}
The equations of ideal SRMHD are conservation of rest mass, of energy--momentum, and
the homogeneous Maxwell equations:
\begin{equation}\label{eq:covariant}
\nabla_\mu(\rho u^\mu)=0,\qquad
\nabla_\mu T^{\mu\nu}=0,\qquad
\nabla_\mu{}^\ast\!F^{\mu\nu}=\nabla_\mu(u^\mu b^\nu-u^\nu b^\mu)=0 .
\end{equation}
The inhomogeneous Maxwell equations only determine the current and charge densities,
which ideal MHD does not need. Together with an equation of state the system
\eqref{eq:covariant} is closed.

\section{Conservative form}\label{sec:consform}

In flat spacetime and Cartesian coordinates, \eqref{eq:covariant} is a system of
conservation laws $\partial_t\UU+\partial_i\FF^i=0$. Reading off the time components:

\paragraph{Mass.} $D\equiv\rho u^0=\rho W$ with flux $F^i(D)=\rho Wv^i$.

\paragraph{Momentum.} $S_j\equiv T^{0j}=wW^2v_j-b^0b_j$. Substituting
\eqref{eq:b0}--\eqref{eq:b2},
\begin{equation}
wW^2v_j-b^0b_j=\Bigl(\rho h+\frac{B^2}{W^2}+\vB^2\Bigr)W^2v_j-W\vB\Bigl(\frac{B_j}{W}+W\vB v_j\Bigr)
=(\rho hW^2+B^2)v_j-\vB B_j .
\end{equation}

\paragraph{Energy.} $T^{00}=wW^2-\Ptot-(b^0)^2$. Since $b^2W^2-(b^0)^2=B^2$ by
\eqref{eq:b0}--\eqref{eq:b2},
$T^{00}=\rho hW^2+B^2-p-\frac12b^2$. One evolves either $E\equiv T^{00}$ or
$\tau\equiv E-D$, which removes the rest-mass energy and is better conditioned in the
Newtonian regime.

\paragraph{Magnetic field.} The $\nu=j$ components of
$\partial_\mu(u^\mu b^\nu-u^\nu b^\mu)=0$ give
$\partial_tB^j+\partial_i(u^ib^j-u^jb^i)=0$, and
$u^ib^j-u^jb^i=Wv^i(B^j/W+b^0v^j)-Wv^j(B^i/W+b^0v^i)=v^iB^j-v^jB^i$. The $\nu=0$
component is the constraint $\partial_iB^i=0$.

\begin{keybox}[Conserved variables and $x$-fluxes]
\begin{equation}\label{eq:UF}
\UU=\begin{pmatrix} D\\ S_j\\ \tau\\ B^k\end{pmatrix}
=\begin{pmatrix}\rho W\\ (\rho hW^2+B^2)v_j-(\vv\cdot\BB)B_j\\
\rho hW^2+B^2-p-\tfrac12\bigl[B^2/W^2+(\vv\cdot\BB)^2\bigr]-\rho W\\ B^k\end{pmatrix},
\end{equation}
\begin{equation}\label{eq:Fx}
\FF^x=\begin{pmatrix} Dv_x\\ wW^2v_xv_j+\Ptot\delta_{xj}-b^xb_j\\ S_x-Dv_x\\
v_xB^k-v^kB^x\end{pmatrix}.
\end{equation}
Equivalent forms: $\SSS=\rho hW^2\vv+\EE\times\BB$,
$\tau=\rho hW^2-p-D+\frac12(|\EE|^2+B^2)$ with $|\EE|^2=|\vv\times\BB|^2=v^2B^2-\vB^2$, and
$T^{xj}=S_jv_x+\Ptot\delta_{xj}-B^xb_j/W$.
\end{keybox}

The equivalent forms follow from $\EE\times\BB=-(\vv\times\BB)\times\BB=\vv B^2-\BB\vB$,
from $B^2-\frac12b^2=\frac12(B^2+v^2B^2-\vB^2)$, and from expanding $b^xb_j$ with
\eqref{eq:bi}. All of \eqref{eq:UF} was checked symbolically against the components of
\eqref{eq:Tmunu} (\cref{app:verification}).

\begin{codebox}[: conserved variables and fluxes]
\file{HLLD/src/physics/hlld.py} builds $\UU$ and $\FF$ exactly as in \eqref{eq:UF}:
\code{primitive_to_conserved} uses $\tau=\rho hW^2-p-D+\tfrac12(B^2+B^2v^2-\vB^2)$, and
\code{compute_srmhd_fluxes} uses $\tau=(\rho h+b^2)W^2-(p+b^2/2)-(b^0)^2-D$, the momentum
flux $wWv_xu_j+\Ptot\delta_{xj}-b^xb_j$ and the induction flux $v_xB^k-v^kB^x$, for any
sweep direction (\code{idir}). The HLLD star-state algebra is written for the total
energy, so \code{energy_form} converts $\tau\to E=\tau+D$ and $F_\tau\to F_E=F_\tau+F_D$.
The exact solver (\file{rmhd_final/rmhd/batched/jump.py}, \code{cons_flux}) also uses
$E=T^{00}$; its state vector is $(\rho,\Ptot,v_x,v_y,v_z,B_y,B_z)$ with $B_n=B_x$ held
separately, and it carries the \emph{total} pressure $\Ptot$, not $p$.
\end{codebox}

\paragraph{One dimension.} In planar symmetry $\partial_y=\partial_z=0$, the
constraint gives $\partial_xB^x=0$ and the $x$-component of the induction equation gives
$\partial_tB^x=0$: the normal field is a constant parameter of the Riemann problem. The
one-dimensional system has seven unknowns,
$\UU=(D,S_x,S_y,S_z,\tau,B^y,B^z)$.

\section{Equation of state and sound speed}

We use the ideal gas $p=(\Gamma-1)\rho\epsilon$, so
\begin{equation}
h=1+\frac{\Gamma}{\Gamma-1}\frac{p}{\rho}.
\end{equation}
The relativistic sound speed is $c_s^2=(\partial p/\partial e)_s$ with $e=\rho(1+\epsilon)$
\parencite{Anton2010}. Along an isentrope the first law gives $\dd\epsilon=(p/\rho^2)\dd\rho$,
hence $\dd e=(1+\epsilon)\dd\rho+\rho\,\dd\epsilon=h\,\dd\rho$ and
\begin{equation}\label{eq:cs}
c_s^2=\frac{1}{h}\Bigl(\frac{\partial p}{\partial\rho}\Bigr)_s=\frac{\Gamma p}{\rho h}
=\frac{\Gamma(\Gamma-1)p}{(\Gamma-1)\rho+\Gamma p}<\Gamma-1 ,
\end{equation}
the bound being approached as $p/\rho\to\infty$; causality ($c_s<1$) holds for
$\Gamma\le2$.

\begin{codebox}[: equation of state]
\file{HLLD/src/physics/eos.py} implements a hybrid EOS, $p=K\rho^\Gamma+
(\Gamma_{\rm th}-1)\rho(\epsilon-\epsilon_{\rm cold})$, with
$c_s^2=(\Gamma p_{\rm cold}+\Gamma_{\rm th}p_{\rm th})/(\rho h)$. All production runs and
scripts construct it with $K=0$ and $\Gamma_{\rm th}=\Gamma$, which is the ideal gas
above (three legacy 1D configs, \file{configs/blastwave.yaml} and
\file{configs/static_with*_Bx.yaml}, use a cold polytrope with $K=1$). The exact solver forms $p=\Ptot-b^2/2$ inside every routine and uses
$h=1+\Gamma p/[(\Gamma-1)\rho]$ and $c_s^2=\Gamma p/(\rho h)$
(\file{rmhd/batched/wave_speeds_b.py}).
\end{codebox}

\section{Primitive recovery}\label{sec:c2p}

The fluxes are functions of the primitive variables $(\rho,\vv,p,\BB)$, but the scheme
evolves $\UU$. Recovering the primitives is a nonlinear problem with no closed-form
solution. Its structure is easiest to see with the variable $Z\equiv\rho hW^2$. From
\eqref{eq:UF},
\begin{equation}
\SSS\cdot\BB=(\rho hW^2+B^2)\vB-\vB B^2=Z\,\vB ,
\end{equation}
so $\vB=(\SSS\cdot\BB)/Z$; squaring $\SSS$ then gives
\begin{equation}\label{eq:c2p-v2}
v^2=\frac{S^2Z^2+(\SSS\cdot\BB)^2(2Z+B^2)}{Z^2(Z+B^2)^2},
\end{equation}
and the energy equation becomes one equation in $Z$ \parencite{MignoneBodo2006},
\begin{equation}\label{eq:c2p-f}
f(Z)\equiv Z-p(Z)+\Bigl(1-\frac{1}{2W^2}\Bigr)B^2-\frac{(\SSS\cdot\BB)^2}{2Z^2}-E=0,
\qquad
p(Z)=\frac{\Gamma-1}{\Gamma}\,\frac{Z-DW}{W^2},
\end{equation}
where $W=W(Z)$ through \eqref{eq:c2p-v2} and $E=\tau+D$. The identities
\eqref{eq:c2p-v2}--\eqref{eq:c2p-f} were checked in \cref{app:verification}.
\textcite{Noble2006} compared one- and two-variable schemes of this kind for general
relativistic MHD; \textcite{Kastaun2021} reformulated the inversion as a root-finding
problem in $\mu=1/(hW)$ (in the rescaled variables $\bm r=\SSS/D$, $\bm b=\BB/\sqrt D$) with a bracketing interval that is guaranteed to contain
the unique solution, which makes it robust across the physical parameter space.

\begin{codebox}[: primitive recovery]
\file{HLLD/src/physics/c2p.py} implements the scheme of \textcite{Kastaun2021},
vectorised over cells: a bracketed root-find for $\mu=1/(hW)$. According to the
guide (Part~2.6) the bracketed secant was found to stall on hot, fast gas
(e.g.\ $(\rho,p,v)=(0.1,1,0.99)$ left a relative error $\sim10^{-6}$), so cells that it
leaves unconverged are re-solved by bisection; cells where recovery still fails are
counted (\code{c2p_bad}); only cells whose recovered density falls below the atmosphere
value are floored.
\end{codebox}

\section{Newtonian limit}\label{sec:newtonian}

Introduce a bookkeeping parameter $\delta$ with $\vv=O(\delta)$, $\BB=O(\delta)$ (Alfvén
speeds comparable with flow speeds) and $p,\rho\epsilon=O(\delta^2)\rho$. Expanding
\eqref{eq:UF} (checked symbolically, \cref{app:verification}):
\begin{align}
D&=\rho+O(\delta^2), & S_x&=\rho v_x+O(\delta^3),\\
\tau&=\rho\epsilon+\tfrac12\rho v^2+\tfrac12B^2+O(\delta^4), &
T^{xx}&=\rho v_x^2+p+\tfrac12B^2-B_x^2+O(\delta^4),
\end{align}
and the energy flux $S_x-Dv_x=\bigl(\tfrac12\rho v^2+\rho\epsilon+p+B^2\bigr)v_x-(\vv\cdot\BB)B_x
+O(\delta^5)$. These are the conserved variables and fluxes of Newtonian ideal MHD in
Heaviside--Lorentz units; the induction equation is unchanged. The relativistic
corrections enter through $W$, through the enthalpy $\rho h$ as inertia, and through the
electric-field terms $\EE\times\BB$ and $|\EE|^2/2$ (displacement current), which make the
momentum of the field $\vv B^2-\BB\vB$ part of $\SSS$.

\section*{Exercises}
\addcontentsline{toc}{section}{Exercises}

\begin{exercise}
Verify $u_\mu b^\mu=0$ and $b^2=B^2/W^2+\vB^2$ directly from \eqref{eq:b0}--\eqref{eq:bi},
and show that $b^2W^2-(b^0)^2=B^2$.
\end{exercise}
\begin{solution}
$b^2W^2-(b^0)^2=(B^2/W^2+\vB^2)W^2-W^2\vB^2=B^2$.
\end{solution}

\begin{exercise}
Show that $\SSS=\rho hW^2\vv+\EE\times\BB$ and $\tau=\rho hW^2-p-D+\frac12(|\EE|^2+B^2)$, and
interpret each term.
\end{exercise}

\begin{exercise}
Show that $T^{ij}=S_jv^i+\Ptot\delta^{ij}-B^ib^j/W$. (Expand $b^ib^j$ with
\eqref{eq:bi}.)
\end{exercise}

\begin{exercise}
Show that $c_s^2<\Gamma-1$ for the ideal gas and that $c_s\to0$ as $p/\rho\to0$. For
$\Gamma=4/3$ what is the largest possible sound speed?
\end{exercise}
\begin{solution}
From \eqref{eq:cs}, $c_s^2=(\Gamma-1)\Gamma p/[(\Gamma-1)\rho+\Gamma p]$, increasing in
$p/\rho$ with limit $\Gamma-1$; for $\Gamma=4/3$, $c_s<1/\sqrt3$.
\end{solution}

\begin{exercise}
Derive \eqref{eq:c2p-v2} and \eqref{eq:c2p-f}. Show that for $\BB=0$ the problem
reduces to one equation for $Z$ with $v^2=S^2/(Z^2)$, i.e.\ $v=S/(E+p)$.
\end{exercise}

\begin{exercise}
Expand $\tau$ to fourth order in $\delta$ and identify the leading relativistic
correction to the kinetic energy.
\end{exercise}
\begin{solution}
With $W^2=1+v^2+v^4+\dots$ and $D=\rho W$, $\rho hW^2-p-D=\rho(1+\epsilon)W^2+p(W^2-1)-\rho W$;
the kinetic part $\rho(W^2-W)=\rho(\frac12v^2+\frac58v^4+\dots)$, so the first
correction is $\frac58\rho v^4$, plus $\rho\epsilon v^2+pv^2$ from the internal energy
and pressure.
\end{solution}

% ======================================================================== chapter 3
\chapter{Characteristics and hyperbolicity}\label{ch:characteristics}

\section{The one-dimensional system}

In planar symmetry the system \eqref{eq:UF} has seven components,
$\UU=(D,S_x,S_y,S_z,\tau,B^y,B^z)^{\mathsf T}$, and $B^x$ is a constant parameter. Written
in quasi-linear form,
\begin{equation}
\partial_t\UU+\mathsf{A}(\UU)\,\partial_x\UU=0,\qquad \mathsf{A}=\frac{\partial\FF^x}{\partial\UU},
\end{equation}
or, in primitive variables $\VV=(\rho,v_x,v_y,v_z,p,B^y,B^z)$,
\begin{equation}
\partial_t\VV+\tilde{\mathsf A}\,\partial_x\VV=0,\qquad
\tilde{\mathsf A}=\Bigl(\frac{\partial\UU}{\partial\VV}\Bigr)^{-1}\frac{\partial\FF^x}{\partial\VV}.
\end{equation}
The two matrices are similar and have the same eigenvalues, the \emph{characteristic speeds}
$\lambda$. A small-amplitude disturbance proportional to a right eigenvector $\bm r_k$
travels at $\lambda_k$. Writing $\mathsf A$ out explicitly is laborious; it is much more
economical to find the characteristic speeds covariantly and then specialise to the
$x$-direction, which is what we do now.

\section{Characteristic speeds from a covariant plane-wave analysis}\label{sec:covariant}

A characteristic hypersurface $\phi(x^\mu)=\text{const}$ is a surface across which
derivatives of the solution may jump. For a plane front moving along $x$ with
laboratory speed $\lambda$ we take $\phi=x-\lambda t$, i.e.
\begin{equation}
\phi_\mu\equiv\partial_\mu\phi=(-\lambda,1,0,0).
\end{equation}
Following \textcite{Anile1989}, \textcite{Komissarov1999} and \textcite{Anton2010} we
build three Lorentz scalars from $\phi_\mu$:
\begin{equation}\label{eq:aBG}
a\equiv u^\mu\phi_\mu=W(v_x-\lambda),\qquad
\calB\equiv b^\mu\phi_\mu=b^x-\lambda b^0,\qquad
G\equiv\phi^\mu\phi_\mu=1-\lambda^2 .
\end{equation}
(\textcite{MignoneBodo2006} use $a=W(\lambda-v_x)$; only $a^2$ enters the quartic below.)

\begin{lemma}[Rest-frame meaning of the scalars]\label{lem:restframe}
Let $V$ be the phase speed of the front measured in the fluid rest frame and $\theta$ the
angle between the rest-frame magnetic field $\bm b$ and the front normal. Then
\begin{equation}\label{eq:Vrest}
V^2=\frac{a^2}{G+a^2},\qquad b^2\cos^2\theta=\frac{\calB^2}{G+a^2}.
\end{equation}
\end{lemma}

\begin{proof}
In the rest frame $u^\mu=(1,\bm 0)$ and $b^\mu=(0,\bm b)$; write $\phi_\mu=(\phi_0',\bm\phi')$.
The front $\phi_0't'+\bm\phi'\cdot\bm x'=\text{const}$ moves along $\bm\phi'/|\bm\phi'|$
with phase speed $V=-\phi_0'/|\bm\phi'|$. Now $a=\phi_0'$, $G=-\phi_0'^2+|\bm\phi'|^2$, so
$|\bm\phi'|^2=G+a^2$ and $V^2=a^2/(G+a^2)$. Also $\calB=\bm b\cdot\bm\phi'
=b|\bm\phi'|\cos\theta$. Both relations involve only scalars and hold in every frame.
\end{proof}

\subsection{The dispersion relation in the fluid rest frame}

Linearise the equations about a uniform state at rest, $u^\mu=(1,\bm0)$, $b^\mu=(0,\bm b)$,
with perturbations proportional to $\exp[ik(x-Vt)]$, and choose axes so that
$\bm b=b(\cos\theta,\sin\theta,0)$. To first order $\delta u^\mu=(0,\delta\vv)$;
$u_\mu b^\mu=0$ gives $\delta b^0=\bm b\cdot\delta\vv$; $B^i=u^0b^i-u^ib^0$ gives
$\delta\BB=\delta\bm b$; and $\delta\Ptot=\delta p+\bm b\cdot\delta\bm b$. From
\eqref{eq:Tmunu}, $\delta T^{0i}=w\,\delta v^i-(\bm b\cdot\delta\vv)\,b^i$ and
$\delta T^{ij}=\delta\Ptot\,\delta^{ij}-b^i\delta b^j-\delta b^ib^j$. The divergence
constraint gives $\delta b_x=0$. The linearised equations are (after dividing by $ik$)
\begin{align}
&-V\delta\rho+\rho\,\delta v_x=0, \label{eq:lin-mass}\\
&-V\bigl[w\,\delta\vv-(\bm b\cdot\delta\vv)\bm b\bigr]+\hat{\bm x}\,\delta\Ptot-b\cos\theta\,\delta\bm b=0,
  \label{eq:lin-mom}\\
&-V\delta\bm b=b\cos\theta\,\delta\vv-\bm b\,\delta v_x, \label{eq:lin-ind}\\
&\delta p=hc_s^2\,\delta\rho . \label{eq:lin-ent}
\end{align}
The terms $b^2\delta\vv$ (inside $w\,\delta\vv$) and $-(\bm b\cdot\delta\vv)\bm b$ (from
$b^0b^i$) in \eqref{eq:lin-mom} together are the linearised electromagnetic momentum
$\EE\times\BB$, i.e.\ the displacement current; Newtonian MHD drops both.

\paragraph{Alfvén waves.} The $z$-components decouple: $-Vw\,\delta v_z-b\cos\theta\,
\delta b_z=0$ and $-V\delta b_z=b\cos\theta\,\delta v_z$ give
\begin{equation}\label{eq:alfven-rest}
V^2=c_a^2\cos^2\theta,\qquad c_a^2\equiv\frac{b^2}{w}=\frac{b^2}{\rho h+b^2}.
\end{equation}
$c_a<1$ always, in contrast with the Newtonian Alfvén speed $B/\sqrt\rho$.

\paragraph{Magnetosonic waves.} Eliminate $\delta\rho$, $\delta p$ and $\delta b_y=
(b/V)(\sin\theta\,\delta v_x-\cos\theta\,\delta v_y)$, so that
$\delta\Ptot=V^{-1}[(\rho hc_s^2+b^2\sin^2\theta)\delta v_x-b^2\sin\theta\cos\theta\,\delta v_y]$.
The $x$- and $y$-components of \eqref{eq:lin-mom}, multiplied by $V$, become
\begin{equation}
\begin{pmatrix}
\rho hc_s^2+b^2\sin^2\theta-V^2(\rho h+b^2\sin^2\theta) & -(1-V^2)b^2\sin\theta\cos\theta\\
-(1-V^2)b^2\sin\theta\cos\theta & b^2\cos^2\theta-V^2(\rho h+b^2\cos^2\theta)
\end{pmatrix}
\begin{pmatrix}\delta v_x\\ \delta v_y\end{pmatrix}=0 .
\end{equation}
Its determinant equals $\rho h\,[\,wV^4-(b^2+\rho hc_s^2+c_s^2b^2\cos^2\theta)V^2+
c_s^2b^2\cos^2\theta\,]$ (checked symbolically), so with $c_a^2=b^2/w$:
\begin{equation}\label{eq:ms-rest}
\boxed{\;V^4-\bigl(c_a^2+c_s^2-c_s^2c_a^2\sin^2\theta\bigr)V^2+c_s^2c_a^2\cos^2\theta=0\;}
\end{equation}
Without the electromagnetic momentum the relation is the Newtonian one, with
$v_A^2=b^2/(\rho h)$ and $v_A^2+c_s^2$ in the middle coefficient; keeping it and writing
$c_a^2=b^2/w$ produces the correction $-c_s^2c_a^2\sin^2\theta$. Together
with the entropy mode ($V=0$, $\delta p=\delta\vv=\delta\bm b=0$, $\delta\rho\ne0$),
\eqref{eq:alfven-rest} and \eqref{eq:ms-rest} account for all seven modes.

\subsection{Laboratory-frame speeds}

Substituting \eqref{eq:Vrest} into the rest-frame relations and multiplying by powers of
$G+a^2$ turns them into polynomials in $\lambda$:

\paragraph{Entropy (contact) wave.} $a=0$:
\begin{equation}
\lambda_{\rm e}=v_x .
\end{equation}

\paragraph{Alfvén waves.} $V^2=c_a^2\cos^2\theta$ becomes $w\,a^2=\calB^2$, i.e.
$\calB\pm\sqrt{w}\,a=0$, or $(b^x-\lambda b^0)\pm\sqrt{w}\,W(v_x-\lambda)=0$. Solving for $\lambda$,
\begin{equation}\label{eq:lamA}
\boxed{\;\lambda_A^{(\pm)}=\frac{b^x\pm\sqrt{w}\,u^x}{b^0\pm\sqrt{w}\,u^0}
=v_x+\frac{B^x(1-v^2)}{\vB\pm\sqrt{w}}\;}
\end{equation}
The second form follows by inserting \eqref{eq:b0}--\eqref{eq:bi} and dividing numerator
and denominator by $W$. It is the form used by \textcite{GiacomazzoRezzolla2006} and by the
code.

\begin{proposition}[Which Alfvén speed is which]\label{prop:alfvensign}
$\sqrt w>|\vB|$, hence $\sgn(\lambda_A^{(\pm)}-v_x)=\pm\sgn(B^x)$. For $B^x>0$,
$\lambda_A^{(-)}<v_x<\lambda_A^{(+)}$; for $B^x<0$ the labels are reversed.
\end{proposition}
\begin{proof}
$w=\rho h+b^2>b^2=B^2/W^2+\vB^2\ge\vB^2$, so $\vB+\sqrt w>0>\vB-\sqrt w$; the sign of
$\lambda_A^{(\pm)}-v_x$ is that of $B^x/(\vB\pm\sqrt w)$.
\end{proof}

\begin{codebox}[: the Alfvén pair for $B_n<0$]
The labels $(\pm)$ in \eqref{eq:lamA} refer to the sign in front of $\sqrt w$, not to the
direction of propagation; \textcite{Anton2010} make the same remark about their
eq.~(38). In the exact solver every helper that returned the pair as
\code{(la, ra)} (left, right) had used the $(-)$ root as ``left''. For $B_n<0$ this
built rotations of the wrong family. The callers now order the pair by value (e.g.\
\file{rmhd/shock_speed.py}: \code{if not (la <= ra): la, ra = ra, la}, commented
``labelled by the sign of sqrt(rho h+b\^{}2): swapped for B\_n < 0''); the low-level
helpers \code{xi} and \code{quartic_coeffs} still return it labelled by sign.
\end{codebox}

\paragraph{Magnetosonic waves.} Inserting \eqref{eq:Vrest} into \eqref{eq:ms-rest},
multiplying by $(G+a^2)^2$ and using $c_a^2/b^2=1/w$ gives
$(1-\zeta^2)a^4-G\,(\zeta^2a^2-c_s^2\calB^2/w)=0$ with $\zeta^2=c_s^2+c_a^2-c_s^2c_a^2$.
Multiplying by $w$ and using $w(1-\zeta^2)=\rho h(1-c_s^2)$ and $w\zeta^2=b^2+\rho hc_s^2$:
\begin{equation}\label{eq:quartic}
\boxed{\;\rho h(1-c_s^2)\,a^4-(1-\lambda^2)\bigl[(b^2+\rho hc_s^2)\,a^2-c_s^2\calB^2\bigr]=0,
\qquad a=W(v_x-\lambda),\ \ \calB=b^x-\lambda b^0\;}
\end{equation}
a quartic in $\lambda$ whose four roots are the slow and fast speeds. We checked
(\cref{app:verification}) that \eqref{eq:quartic} is identical to eq.~(56) of
\textcite{MignoneBodo2006}, to the normalised form eq.~(4.50) of
\textcite{GiacomazzoRezzolla2006} and eq.~(16) of \textcite{DelZanna2003} (which divide by
$w$ and use $\tilde b^\mu=b^\mu/\sqrt w$), to $c_s^2\calN_4=0$ with $\calN_4$ of
\textcite[eq.~35]{Anton2010}, and coefficient by coefficient to the quartic coded in
\file{rmhd/wave_speeds.py} and \file{rmhd/batched/wave_speeds_b.py}.\footnote{The same
quartic is attributed also to \textcite{Leismann2005}; that paper's full text could not
be accessed for this check, so the comparison with it was not made.}

\begin{remark}[The full characteristic determinant]
In the covariant variables $(u^\mu,b^\mu,p,s)$ of \textcite{Anile1989} the system has ten
unknowns, and the characteristic determinant factorises as
$\det(\mathcal A^\alpha\phi_\alpha)=w\,a^2\calA^2\calN_4$ with $\calA=wa^2-\calB^2$
\parencite[eq.~33]{Anton2010}. The squares reflect three unphysical roots of the enlarged
system: the extra root $a=0$ carries $\delta b^\mu\propto\phi^\mu$, i.e.\ $\delta B^x\neq0$, and
is excluded by $\nabla\cdot\BB=0$; the extra pair $\calA=0$ violates the linearised
constraints $u_\mu\delta u^\mu=0$ and $u_\mu\delta b^\mu+b_\mu\delta u^\mu=0$; the physical speeds are the seven
roots of $a\,\calA\,\calN_4=0$ found above. \Cref{app:eigen} gives the matrix and the
eigenvectors.
\end{remark}

\section{Rest-frame speeds and their ordering}\label{sec:ordering}

Write $X=V^2$ and $P(X)=X^2-QX+c_s^2c_a^2\cos^2\theta$ with
$Q=c_a^2+c_s^2-c_s^2c_a^2\sin^2\theta$. Let $X_s\le X_f$ be the roots (slow and fast).

\begin{proposition}[Ordering and bounds]\label{prop:ordering}
For $0<c_s<1$ and $0<c_a<1$,
\begin{equation}
0\le X_s\le c_a^2\cos^2\theta\le X_f\le\zeta^2<1,\qquad \zeta^2\le 2X_f .
\end{equation}
Equality in any of the first three inequalities requires $\cos\theta=0$ or
$\sin\theta=0$.
\end{proposition}

\begin{proof}
$P$ is an upward parabola in $X$.
\begin{enumerate}[label=(\roman*),nosep]
\item $P(0)=c_s^2c_a^2\cos^2\theta\ge0$ and
  $P(c_a^2\cos^2\theta)=-c_a^4\cos^2\theta\sin^2\theta\,(1-c_s^2)\le0$, so the Alfvén
  value lies between the roots and both roots are non-negative.
\item $Q-\zeta^2=c_s^2c_a^2\cos^2\theta\ge0$, hence
  $P(\zeta^2)=c_s^2c_a^2\cos^2\theta\,(1-\zeta^2)\ge0$; and the vertex lies to the left of
  $\zeta^2$, because $\zeta^2-Q/2=\tfrac12(\zeta^2-c_s^2c_a^2\cos^2\theta)\ge0$. Therefore
  $\zeta^2$ lies to the right of both roots, $X_f\le\zeta^2$.
\item $\zeta^2=1-(1-c_s^2)(1-c_a^2)<1$.
\item $X_f\ge Q/2\ge\zeta^2/2$.
\end{enumerate}
The polynomial identities used here were checked symbolically.
\end{proof}

So in the fluid frame the speeds are ordered exactly as in Newtonian MHD,
slow $\le$ Alfvén $\le$ fast, all are subluminal, and the fast speed is largest
perpendicular to the field, where it equals $\zeta$. \Cref{fig:friedrichs} shows the
phase-speed (Friedrichs) diagram.

\begin{figure}[t]
\centering
\begin{tikzpicture}
\begin{polaraxis}[width=7.6cm,ymax=0.85,ytick={0.2,0.4,0.6,0.8},xtick={0,30,...,330},
  legend style={at={(1.32,0.98)},anchor=north west,font=\small,draw=none},
  every axis plot/.append style={line width=0.9pt}]
  % c_s^2 = 0.2, c_a^2 = 0.45: Q = 0.65 - 0.09 sin^2
  \addplot[blue!70!black,domain=0:360,samples=361,data cs=polar]
    ({x},{sqrt((0.65-0.09*sin(x)^2+sqrt((0.65-0.09*sin(x)^2)^2-0.36*cos(x)^2))/2)});
  \addplot[red!70!black,domain=0:360,samples=361,data cs=polar]
    ({x},{sqrt(0.45)*abs(cos(x))});
  \addplot[green!45!black,domain=0:360,samples=361,data cs=polar]
    ({x},{sqrt(max(0,(0.65-0.09*sin(x)^2-sqrt((0.65-0.09*sin(x)^2)^2-0.36*cos(x)^2))/2))});
  \addplot[gray,dashed,domain=0:360,samples=181,data cs=polar] ({x},{sqrt(0.56)});
  \legend{fast,Alfvén,slow,$\zeta$}
\end{polaraxis}
\end{tikzpicture}
\caption{Rest-frame phase speeds $V(\theta)$ for $c_s^2=0.2$, $c_a^2=0.45$
($\theta$ measured from $\bm b$, which points to the right). The fast speed reaches
$\zeta=\sqrt{0.56}$ (dashed) at $\theta=90^\circ$; at $\theta=0$ the Alfvén speed
coincides with the fast speed (a Type~II degeneracy, \cref{sec:degeneracies}).}
\label{fig:friedrichs}
\end{figure}

\section{Properties of the laboratory-frame speeds}

\paragraph{Ordering.} The seven laboratory speeds are ordered as
\begin{equation}\label{eq:order}
\lambda_{\rm f}^-\le\lambda_{\rm A}^-\le\lambda_{\rm s}^-\le\lambda_{\rm e}\le
\lambda_{\rm s}^+\le\lambda_{\rm A}^+\le\lambda_{\rm f}^+ ,
\end{equation}
where now the superscripts denote the lower and upper value of each pair
\parencite[eq.~42]{Anton2010}, a result due to \textcite{Anile1989}. Aberration destroys
the symmetry of each pair about $v_x$. We do not reprove \eqref{eq:order}; the following
lemma is the key monotonicity behind it and is useful in its own right.

\begin{lemma}[Monotone map]\label{lem:monotone}
Along the $x$-direction, the rest-frame phase speed $V(\lambda)$ of \eqref{eq:Vrest}
increases monotonically from $0$ at $\lambda=v_x$ to $1$ at $\lambda=\pm1$ on either side
of $v_x$.
\end{lemma}
\begin{proof}
$V^2=1/(1+R)$ with $R=G/a^2=(1-\lambda^2)/[W^2(v_x-\lambda)^2]$, and
$\dd R/\dd\lambda=2(1-\lambda v_x)/[W^2(v_x-\lambda)^3]$, which is negative for
$\lambda>v_x$ and positive for $\lambda<v_x$ (since $1-\lambda v_x>0$).
\end{proof}

\begin{proposition}[An isotropic bound on the fast speeds]\label{prop:fastbound}
Let $\lambda_\zeta^\pm$ be the roots of $(1-\zeta^2)W^2(\lambda-v_x)^2=\zeta^2(1-\lambda^2)$,
\begin{equation}\label{eq:lamzeta}
\lambda_\zeta^\pm=\frac{v_x(1-\zeta^2)\pm\zeta\sqrt{(1-v^2)\,[1-v^2\zeta^2-v_x^2(1-\zeta^2)]}}
{1-v^2\zeta^2}.
\end{equation}
Then $\lambda_\zeta^-\le\lambda_{\rm f}^-$ and $\lambda_{\rm f}^+\le\lambda_\zeta^+$.
\end{proposition}
\begin{proof}
At $\lambda_\zeta^\pm$ the rest-frame phase speed is exactly $\zeta$; at $\lambda_{\rm
f}^\pm$ it is $\sqrt{X_f(\theta)}\le\zeta$ (\cref{prop:ordering}). By
\cref{lem:monotone} the fast root lies closer to $v_x$.
\end{proof}

Equation \eqref{eq:lamzeta} is the speed of a hypothetical isotropic wave of rest-frame
speed $\zeta$, boosted with the flow. \textcite{GammieMcKinneyToth2003} introduced this
replacement of the full dispersion relation (their eq.~28) and noted that it overestimates
the comoving fast speed by a bounded factor; \cref{prop:ordering} makes this precise:
$\zeta\le\sqrt2\,\sqrt{X_f}$. Because it needs only a square root it is a convenient
signal-speed estimate; HARM uses it, and so does the reference code.

\paragraph{Special cases.} For $B^x=0$ one finds $\calB=b^x-\lambda b^0=\vB\,a$ and
\eqref{eq:quartic} factorises as
\begin{equation}\label{eq:quad-bx0}
a^2\Bigl\{\rho h(1-c_s^2)a^2-(1-\lambda^2)\bigl[\rho hc_s^2+b^2-c_s^2\vB^2\bigr]\Bigr\}=0,
\end{equation}
so the slow speeds collapse onto $\lambda_e=v_x$ (as do both Alfvén speeds, by
\eqref{eq:lamA}) and the fast pair solves a quadratic; this is eq.~(58) of
\textcite{MignoneBodo2006}. For $\vv=0$, \eqref{eq:quartic} reduces to the biquadratic
$w\lambda^4-(b^2+\rho hc_s^2+c_s^2B_x^2)\lambda^2+c_s^2B_x^2=0$, their eq.~(57), which is
\eqref{eq:ms-rest} with $\cos^2\theta=B_x^2/B^2$.

\begin{codebox}[: wave speeds]
\file{rmhd/wave_speeds.py} (scalar) and \file{rmhd/batched/wave_speeds_b.py} (batched)
compute the coefficients of \eqref{eq:quartic} divided by $w$, with $\tilde b^\mu=
b^\mu/\sqrt w$ and \code{eps2} $=\zeta^2$, exactly as in eq.~(4.50) of
\textcite{GiacomazzoRezzolla2006}. The scalar code has a separate expression for
$|B_x|\le10^{-15}$. The batched code finds the roots by Ferrari's closed form (default) or
the eigenvalues of the companion matrix, polishes them with Newton iterations on the
quartic, and returns the Alfvén speeds from the same intermediates. The HLL/HLLD signal
speeds in \file{HLLD/src/physics/hlld.py} (\code{compute_cp_cm}) are the isotropic bound
\eqref{eq:lamzeta} with \code{v02} $=\zeta^2$; $c_{\max}$ and $c_{\min}$ are then
clamped to be non-negative.
\end{codebox}

\section{Degeneracies}\label{sec:degeneracies}

The inequalities in \eqref{eq:order} can be equalities. The classification follows
\textcite{Komissarov1999} and \textcite{Anton2010}, who characterise it by the component
$b_n$ of $b^\mu$ along the Alfvén front normal in the fluid frame,
$b_n^2=\calB^2/(G+a^2)$ evaluated at the Alfvén speed:
\begin{description}[nosep,leftmargin=1.2em]
\item[Type I] ($b_n=0$; for $x$-propagation, $B^x=0$): both Alfvén speeds and both slow
  speeds coincide with the entropy speed, $\lambda_A^\pm=\lambda_s^\pm=\lambda_e=v_x$.
  The Riemann fan has three waves: two fast waves and a tangential discontinuity.
\item[Type II] ($b_t=0$, i.e.\ $b_n^2=b^2$; the field is along the wave normal in the fluid
  frame): an Alfvén speed coincides with the fast or the slow speed of the same class.
  In Newtonian MHD both Alfvén waves are then degenerate. In RMHD, because of aberration,
  the condition can hold for one Alfvén wave and not the other; only when the tangential
  velocity vanishes is the classical behaviour recovered \parencite{Anton2010}.
\item[Type II$'$] (Type II with $c_s^2=c_a^2$): an Alfvén speed coincides with
  \emph{both} magnetosonic speeds of its class (a triple point, the relativistic
  counterpart of the triple umbilic point of Newtonian MHD; cf.\ \cref{fig:friedrichs}
  at $\theta=0$ when $c_s=c_a$).
\end{description}

At degenerate states the eigenvectors of \cref{app:eigen} vanish or become linearly
dependent, so they no longer form a basis; for a solver that projects on characteristic
fields this is fatal. \textcite{BrioWu1988} cured this in Newtonian MHD by renormalising
the eigenvectors; \textcite{Komissarov1999} gave relativistic eigenvectors valid through
Type~II degeneracies; \textcite{Anton2010} gave renormalised right and left eigenvectors
that are regular and complete at all physical states, including both types (their
eqs.~52 and 65--67 for the Alfvén waves).

\paragraph{Loss of strict hyperbolicity and of convexity.} At degenerate states the system
is hyperbolic but not strictly hyperbolic. Moreover, the entropy and Alfvén
characteristic fields are linearly degenerate ($\nabla_{\UU}\lambda_k\cdot\bm r_k\equiv0$),
while the magnetosonic fields are not genuinely nonlinear everywhere: their convexity
depends on the magnetic field \parencite{BrioWu1988,Ibanez2015}. Both properties are
hypotheses of the classical existence and uniqueness theory, and their failure is what
allows compound waves and non-unique solutions (\cref{ch:nonunique}).

\paragraph{Near-degeneracy.} For a numerical solver, nearly coincident speeds are worse
than exactly coincident ones: the system is nominally seven-wave but the Jacobian of any
formulation that treats the two waves separately is nearly singular. On the interface
states of the reference 2D rotor run the Alfvén and slow speeds differ by a median of
$5.5\times10^{-3}$, against $0.28$ between fast and Alfvén (the guide, \S1.2).

\begin{codebox}[: classifying an interface]
\file{rmhd/batched/classify.py} classifies each interface before any solve, with
thresholds relative to the width of the wave fan: \code{FULL7} (seven waves),
\code{COPLANAR}, \code{ALFVEN_DEGEN} (Type I), \code{SLOW_DEGEN} and \code{FAST_DEGEN}
(Type II with the slow or fast wave). The degenerate classes are handed to a reduced
solver (the guide, \S1.7).
\end{codebox}

\section*{Exercises}
\addcontentsline{toc}{section}{Exercises}

\begin{exercise}
Show that the Alfvén speeds \eqref{eq:lamA} are the roots of $\calA=wa^2-\calB^2=0$, and
that $\lambda_A^{(+)}\lambda_A^{(-)}$ and $\lambda_A^{(+)}+\lambda_A^{(-)}$ are rational
functions of the state. Evaluate both speeds for $\vv=0$.
\end{exercise}
\begin{solution}
For $\vv=0$: $\vB=0$, $W=1$, so $\lambda_A^{(\pm)}=\pm B^x/\sqrt w=\pm c_a\cos\theta$, the
rest-frame result.
\end{solution}

\begin{exercise}
Carry out the elimination leading to \eqref{eq:ms-rest} and show that for $\theta=\pi/2$
the fast speed is $\zeta$ and the slow speed vanishes, while for $\theta=0$ the roots are
$c_s^2$ and $c_a^2$.
\end{exercise}

\begin{exercise}
Show that $\zeta\le\sqrt2\max(c_s,c_a)$ and find the state for which the isotropic bound
overestimates the fast speed the most. Compare with the statement of
\textcite{GammieMcKinneyToth2003}.
\end{exercise}
\begin{solution}
The worst case is $\theta=0$, where $X_f=\max(c_s^2,c_a^2)$; for $c_s=c_a=c\ll1$,
$\zeta^2=2c^2-c^4\approx2X_f$.
\end{solution}

\begin{exercise}
Starting from \eqref{eq:quartic} with $B^x=0$, derive \eqref{eq:quad-bx0}. Show that if
moreover $\vv\cdot\BB=0$ the fast speeds are those of a gas with sound speed $\zeta$.
\end{exercise}

\begin{exercise}
For $B^x<0$, which of $\lambda_A^{(\pm)}$ is the left-going Alfvén speed? Construct a
state with $B^x=-1$ and check numerically that using the $(-)$ root as ``left'' gives a
speed larger than $v_x$.
\end{exercise}

\begin{exercise}
Using \cref{lem:monotone}, show that every characteristic speed satisfies
$|\lambda|<1$ whenever $c_s<1$.
\end{exercise}
\begin{solution}
All rest-frame speeds are $\le\zeta<1$; the monotone map sends $V<1$ to $|\lambda|<1$.
\end{solution}


% ======================================================================== chapter 4
\chapter{Discontinuities}\label{ch:discontinuities}

\section{Rankine--Hugoniot conditions}

Let a discontinuity move with constant speed $V_s$ and separate the state $a$ ahead of it
from the state $b$ behind it. Applying the integral form \eqref{eq:integralform} to a
small rectangle that the discontinuity crosses diagonally and letting it shrink gives the
Rankine--Hugoniot (RH) conditions
\begin{equation}\label{eq:RH}
\jump{\FF^x}=V_s\jump{\UU},\qquad\text{i.e.}\quad \FF^x_a-\FF^x_b=V_s(\UU_a-\UU_b).
\end{equation}
With \eqref{eq:UF}--\eqref{eq:Fx} these are seven conditions (mass, three momentum
components, energy, two tangential field components), plus $\jump{B^x}=0$.

\paragraph{Covariant form.} The front is the hypersurface $x-V_st=\text{const}$, with
unit spacelike normal
\begin{equation}
n_\mu=W_s(-V_s,1,0,0),\qquad n_\mu n^\mu=1,\qquad W_s=(1-V_s^2)^{-1/2}.
\end{equation}
For any tensor current, $X^{\mu\nu}n_\nu=W_s(X^{\mu x}-V_sX^{\mu0})$, so \eqref{eq:RH}
is equivalent to \parencite{Anile1989,GiacomazzoRezzolla2006}
\begin{equation}\label{eq:RHcov}
\jump{\rho u^\mu}n_\mu=0,\qquad \jump{T^{\mu\nu}}n_\nu=0,\qquad
\jump{{}^\ast\!F^{\mu\nu}}n_\nu=0 .
\end{equation}
For the last one, ${}^\ast\!F^{00}=0$ and ${}^\ast\!F^{0x}=B^x$ give $\jump{B^x}=0$,
while ${}^\ast\!F^{j0}=-B^j$ and ${}^\ast\!F^{jx}=u^jb^x-u^xb^j=-(v_xB^j-v^jB^x)$ give
the induction jump $V_s\jump{B^j}=\jump{v_xB^j-v^jB^x}$.

\section{Mass flux and the classification of discontinuities}

Define the rest-mass flux through the front and the normal component of the comoving
field,
\begin{equation}\label{eq:J}
J\equiv-\rho u^\mu n_\mu=\rho WW_s(V_s-v_x),\qquad
B_n\equiv b^\mu n_\mu=\frac{W_s}{W}B^x-\frac{\vB J}{\rho}.
\end{equation}
(Not to be confused with the lab normal field $B^x$, which the code and the later chapters
call $B_n$; the two coincide in general only when $J=0$ and $v_t=0$.)
These are eqs.~(4.5) and (4.28) of \textcite{GiacomazzoRezzolla2006}. The first condition
of \eqref{eq:RHcov} says $\jump{J}=0$. It also gives a useful identity: for any $Q$,
\begin{equation}\label{eq:jumpidentity}
\jump{Q\,(V_s-v_x)}=\frac{J}{W_s}\Bigl\llbracket\frac{Q}{D}\Bigr\rrbracket ,
\end{equation}
because $V_s-v_x=J/(W_sD)$ on both sides. Two classes follow:
\begin{itemize}[nosep]
\item $J=0$: no mass crosses the front; it moves with the fluid, $V_s=v_{x,a}=v_{x,b}$.
  These are \emph{contact} and \emph{tangential} discontinuities.
\item $J\neq0$: mass crosses the front. These are \emph{shocks} (compressive,
  $\jump{\rho}\ne0$) and \emph{rotational} (Alfvén) discontinuities.
\end{itemize}

It is convenient to write the projections of the fluxes on $n_\nu$ with \eqref{eq:J}.
With $u^\mu n_\mu=-J/\rho$ and $h_t\equiv w/\rho=h+b^2/\rho$ (the total specific
enthalpy),
\begin{align}
T^{\mu\nu}n_\nu&=-Jh_t\,u^\mu+\Ptot\,n^\mu-B_n\,b^\mu, \label{eq:Tn}\\
{}^\ast\!F^{\mu\nu}n_\nu&=u^\mu(b^\nu n_\nu)-b^\mu(u^\nu n_\nu)=B_n\,u^\mu+\frac{J}{\rho}\,b^\mu
\equiv C^\mu. \label{eq:Cmu}
\end{align}
The RH conditions are therefore $\jump{J}=0$, $\jump{-Jh_tu^\mu+\Ptot n^\mu-B_nb^\mu}=0$
and $\jump{C^\mu}=0$.

\section{Contact and tangential discontinuities}

\begin{proposition}\label{prop:contact}
Let $J=0$.
\begin{enumerate}[label=(\alph*),nosep]
\item If $B^x\ne0$ (contact discontinuity), then $\vv$, $\BB$, $\Ptot$ and $p$ are
  continuous; only $\rho$ (and with it $\epsilon$, $h$) can jump.
\item If $B^x=0$ (tangential discontinuity), then $v_x$ and $\Ptot$ are continuous, and
  $\rho$, $p$, the tangential velocity and the tangential field can all jump.
\end{enumerate}
\end{proposition}

\begin{proof}
With $J=0$, \eqref{eq:J} gives $B_n=W_sB^x/W$ and \eqref{eq:Cmu} gives
$C^\mu=(W_sB^x/W)u^\mu=W_sB^x(1,\vv)$.
(a) $\jump{C^\mu}=0$ and $B^x\neq0$ give $\jump{\vv}=0$, hence $\jump{W}=0$ and
$\jump{B_n}=0$. The momentum condition reduces to $\jump{\Ptot}n^\mu=B_n\jump{b^\mu}$;
contracting with $n_\mu$ gives $\jump{\Ptot}=B_n\jump{B_n}=0$, and then $\jump{b^\mu}=0$.
Since $B^i=u^0b^i-u^ib^0$ with $u^\mu$ continuous, $\jump{\BB}=0$; $b^2$ is continuous, so
$p=\Ptot-b^2/2$ is too.
(b) With $B^x=0$, $B_n=0$ and $C^\mu=0$ identically; the momentum condition is
$\jump{\Ptot}n^\mu=0$. The definition of $J$ gives $V_s=v_{x,a}=v_{x,b}$. Nothing else is
constrained.
\end{proof}

Across a contact the solution is thus fixed by $(\rho_a,\rho_b)$ alone. Across a
tangential discontinuity the gas pressure may jump as long as the total pressure does
not: magnetic pressure can balance gas pressure. This is the central wave of the
three-wave fan for $B^x=0$ (\cref{sec:bx0}).

\section{Three invariants of discontinuities with \texorpdfstring{$J\ne0$}{J != 0}}

\begin{proposition}[Shock invariants]\label{prop:invariants}
Across any discontinuity with $J\neq0$,
\begin{align}
&\jump{H}=0,\qquad H\equiv\frac{B_n^2}{J^2}-\frac{b^2}{\rho^2}, \label{eq:invH}\\
&\jump{h\,B_n}=0, \label{eq:invhB}\\
&\jump{\Ptot}+J^2\Bigl\llbracket\frac{h}{\rho}\Bigr\rrbracket=0, \label{eq:magtaub}
\end{align}
where $h$ is the specific enthalpy of the gas and $\Ptot$ the total pressure.
\end{proposition}

\begin{proof}
$C^\mu$ in \eqref{eq:Cmu} is continuous, and so is every scalar built from it and from
other continuous vectors. Using $u_\mu u^\mu=-1$, $u_\mu b^\mu=0$, $u_\mu n^\mu=-J/\rho$,
$b_\mu n^\mu=B_n$, $n_\mu n^\mu=1$:
\begin{itemize}[nosep]
\item $C_\mu C^\mu=J^2b^2/\rho^2-B_n^2=-J^2H$. Since $J$ is continuous and nonzero,
  $\jump{H}=0$.
\item Contract $C_\mu$ with the continuous vector $T^{\mu\nu}n_\nu$ of \eqref{eq:Tn}:
  $C_\mu u^\mu=-B_n$, $C_\mu n^\mu=B_nJ/\rho-B_nJ/\rho=0$, $C_\mu b^\mu=Jb^2/\rho$, so
  $C_\mu T^{\mu\nu}n_\nu=JB_n(h_t-b^2/\rho)=JB_nh$. Hence $\jump{hB_n}=0$.
\item Contract \eqref{eq:Tn} with $n_\mu$: $n_\mu T^{\mu\nu}n_\nu=J^2h_t/\rho+\Ptot-B_n^2$,
  which is continuous. Writing $h_t/\rho=h/\rho+b^2/\rho^2$,
  $\jump{\Ptot}+J^2\jump{h/\rho}-J^2\jump{B_n^2/J^2-b^2/\rho^2}=0$, and the last term
  vanishes by \eqref{eq:invH}.\qedhere
\end{itemize}
\end{proof}

Equation \eqref{eq:magtaub} is eq.~(4.26) of \textcite{GiacomazzoRezzolla2006}, who quote
the invariance of $H$ from \textcite{Anile1989}; they use \eqref{eq:invhB} in their
$B_t$-method. Remarkably, \eqref{eq:magtaub} has the same form as in hydrodynamics, with
the total pressure but the \emph{gas} enthalpy. A second relation, the
\emph{Lichnerowicz adiabat}, follows from the remaining jump conditions; in our notation it
reads \parencite[eq.~4.27]{GiacomazzoRezzolla2006}
\begin{equation}\label{eq:lichnerowicz}
\jump{h^2}-\Bigl(\!\Bigl(\frac{h}{\rho}\Bigr)\!\Bigr)\jump{\Ptot}-H\jump{b^2}
+2\Bigl\llbracket\frac{b^2h}{\rho}\Bigr\rrbracket-2J^2H\Bigl\llbracket\frac{h}{\rho}\Bigr\rrbracket=0 .
\end{equation}
We have not reproduced its derivation, but all of
\eqref{eq:invH}--\eqref{eq:lichnerowicz} were checked to 40 digits on a fast shock
constructed numerically from the full RH system (\cref{app:verification}). For $\bm b=0$,
$H=0$ and \eqref{eq:lichnerowicz} is the Taub adiabat of \cref{sec:taub}.

\section{Rotational (Alfvén) discontinuities}\label{sec:rotational}

\begin{proposition}[Alfvén discontinuities]\label{prop:alfven}
Suppose $J\neq0$ and $\rho$, $p$ and $b^2$ are continuous (hence $h$, $w$ and $\Ptot$ are).
Then either the discontinuity is trivial ($\jump{u^\mu}=0$) or
\begin{equation}\label{eq:JAlfven}
J^2=\frac{\rho^2B_n^2}{w}.
\end{equation}
Writing $J=s\rho B_n/\sqrt w$ with $s=\pm1$, the four-vector
\begin{equation}\label{eq:sigma}
\sigma^\mu\equiv b^\mu+s\sqrt{w}\,u^\mu
\end{equation}
is continuous across the front, and the front moves with the Alfvén speed of either state,
\begin{equation}\label{eq:VsAlfven}
V_s=\frac{\sigma^x}{\sigma^0}=\frac{b^x+s\sqrt w\,u^x}{b^0+s\sqrt w\,u^0}=\lambda_A^{(s)} .
\end{equation}
\end{proposition}

\begin{proof}
From \eqref{eq:invhB} with $h$ continuous, $B_n$ is continuous. The momentum and Maxwell
conditions become $-Jh_t\jump{u^\mu}-B_n\jump{b^\mu}=0$ and
$B_n\jump{u^\mu}+(J/\rho)\jump{b^\mu}=0$. Eliminating $\jump{b^\mu}=-(\rho B_n/J)\jump{u^\mu}$
gives $(J^2h_t-\rho B_n^2)\jump{u^\mu}=0$, i.e.\ \eqref{eq:JAlfven} since $\rho h_t=w$.
Then $\rho B_n/J=s\sqrt w$ and $\jump{b^\mu+s\sqrt wu^\mu}=0$. Finally,
$J=s\rho B_n/\sqrt w$ with \eqref{eq:J} reads $\sqrt w\,W(V_s-v_x)=s(b^x-V_sb^0)$, whose
solution is \eqref{eq:VsAlfven}; it can be evaluated on either side because $\sigma^\mu$ is
continuous.
\end{proof}

The continuity of $\sigma^\mu$ is the relativistic Walén relation. In the Newtonian
limit $\sigma^\mu\to(s\sqrt\rho,\ s\sqrt\rho\,\vv+\BB)$, so $\vv+s\BB/\sqrt\rho$ and the
normal velocity are continuous and the tangential field rotates rigidly. In SRMHD the
consequences are different \parencite{Komissarov1997}; in the words of later summaries
\parencite{MignoneBodo2006,MignoneUglianoBodo2009}, the thermodynamic quantities stay
continuous, but the tangential field components trace ellipses rather than circles and
the normal velocity $v_x$ jumps. The spatial vector
$\bm K\equiv\bm\sigma/\sigma^0$ is continuous, and its $x$-component is the Alfvén speed;
this is the invariant on which the HLLD solver of \cref{sec:hlld} is built.

\begin{example}[A relativistic rotational discontinuity]
\Cref{app:verification} constructs rotational discontinuities numerically: fix the state
ahead, compute $\sigma^\mu$, pick $u^z$ behind, and solve $u_\mu\sigma^\mu=-s\sqrt w$
(which is $u_\mu b^\mu=0$) together with $\jump{B^x}=0$ for $u^x,u^y$ behind. All seven RH
conditions then hold to $10^{-40}$ with $V_s=\lambda_A^{(s)}$. For the ahead state
$(\rho,\Ptot,\vv,\BB)=(1,\,2,\,(0.2,0.3,-0.1),\,(0.8,0.7,0.4))$, $\Gamma=5/3$, one
left-going example has $v_x$ changing by $8.9\times10^{-3}$ and $|B_t|$ by $-2.3\%$ across
the front, while $\rho$, $p$ and $b^2$ are unchanged.
\end{example}

\paragraph{Coplanar data.} If the tangential fields and velocities of the two states lie in
one plane containing the normal (as on every face of a two-dimensional flow with
$B_z=v_z=0$), the only rotational discontinuities that keep the solution in that plane are
the trivial one and a reversal of the direction of the tangential field, in general with a
change of $|B_t|$. In the solver record (\S2) this is the reason the
Alfvén unknowns degenerate on 2D runs: a continuous angle collapses to a binary choice.

\begin{codebox}[: Alfvén discontinuities]
Production uses \code{solve_alfven_pinned} (\file{rmhd/alfven_solver.py}) and its batched
form \code{solve_pinned} (\file{rmhd/batched/alfven_b.py}). These hold $\rho$, $\Ptot$ and
the direction $\psi=\operatorname{atan2}(B_z,B_y)$ of $B_t$ fixed, take $|B_t|$ from the
continuity of $b^2$, and solve for $(v_x,v_y,v_z)$ behind, at $V_s=\lambda_A$, by a damped
least-squares Newton iteration on the seven normalised RH residuals of
\file{rmhd/batched/jump.py} (scaled by the flux magnitudes, not by the jump, so that weak
discontinuities are not penalised). The iteration starts from a rigid rotation of $B_t$ and
$\bm v_t$ at fixed $W$, which is far from a strongly relativistic rotation: across the
$180^\circ$ rotation of Balsara's test~1 (\cref{ch:nonunique}) $|v_y|$ grows from $0.036$ to
$0.765$ and $W$ from $1.08$ to $1.69$, and from the rigid start the Newton does not
converge. Lanes on which it fails are therefore retried from the rotation built as in the
example above (\code{exact_rotation} in \file{rmhd/batched/fullcontact_b.py}, switch
\code{RMHD_ALFVEN_EXACT}, on by default). Because $B^i=\sigma^iu^0-\sigma^0u^i$ with
$\sigma^\mu$ continuous, the conditions $u_\mu\sigma^\mu=-s\sqrt w$, $\jump{B^x}=0$ and
``$B_t$ along $\psi$'' are linear in the $u^\mu$ behind: they define a line, which meets
$u_\mu u^\mu=-1$ in two points, and the physical one is the rotation. In two-dimensional runs
this retry raised the fraction of attempted faces solved exactly from 97.6\% to 98.5\% on
the rotor, from 97.7\% to 99.1\% on Orszag--Tang and from 93.1\% to 96.0\% on the
four-quadrant problem (solver record \S5e). The five-equation Newton on eqs.~(4.15)--(4.17),
(4.19)--(4.20) of \textcite{GiacomazzoRezzolla2006} (\code{solve_alfven}) is only a
fallback, for $B_t\approx0$ and in the scalar code's branch search; roots of those reduced
equations can violate conservation, so its results are gated by the same full RH
residual.
\end{codebox}

\section{Shocks}\label{sec:shocks}

\subsection{The relativistic hydrodynamic Taub adiabat}\label{sec:taub}

For $\bm b=0$ the momentum condition is $\jump{-Jhu^\mu+p\,n^\mu}=0$, i.e.
$J\jump{hu^\mu}=\jump{p}n^\mu$.

\begin{proposition}[Taub adiabat]\label{prop:taub}
Across a relativistic hydrodynamic shock \parencite{Taub1948}
\begin{equation}\label{eq:taub}
\jump{h^2}=\Bigl(\frac{h_a}{\rho_a}+\frac{h_b}{\rho_b}\Bigr)\jump{p},\qquad
J^2=-\frac{\jump{p}}{\jump{h/\rho}} .
\end{equation}
\end{proposition}
\begin{proof}
Contract $J(h_au_a^\mu-h_bu_b^\mu)=\jump{p}n^\mu$ with $h_au_{a\mu}+h_bu_{b\mu}$. The left
side is $J(-h_a^2+h_b^2)$, the cross terms cancelling. The right side is
$\jump{p}(h_au_a\cdot n+h_bu_b\cdot n)=-\jump{p}J(h_a/\rho_a+h_b/\rho_b)$. Divide by $-J$.
The second relation is the $n_\mu$-contraction, $\jump{J^2h/\rho+p}=0$, i.e.\
\eqref{eq:magtaub} with $\bm b=0$.
\end{proof}

For a given ahead state and post-shock pressure, \eqref{eq:taub} with the EOS determines
$h_b$ and $\rho_b$; for the ideal gas it is a quadratic for $h_b$ (\cref{ch:hydro}). In the
Newtonian limit $h\to1+\epsilon+p/\rho$ with $\epsilon,p/\rho\ll1$ and \eqref{eq:taub}
reduces to the Hugoniot relation $\jump{\epsilon}=-\tfrac12(p_a+p_b)\jump{1/\rho}$
(\cref{ex:hugoniot}). \textcite{Thorne1973} discusses the properties of the Taub adiabat.

\subsection{Magnetosonic shocks}

Fast and slow shocks have $J\ne0$, $\jump{\rho}\ne0$, and all components may jump. Given
the ahead state and the post-shock \emph{total} pressure $\Ptot_b$, the jump conditions
can be solved in closed form for the post-shock state as a function of $V_s$
\parencite[eqs.~4.13--4.25 and Appendix~A]{GiacomazzoRezzolla2006}; $V_s$ is then the
root of \eqref{eq:magtaub}, which must be found numerically. This is the ``$p$-method''.
The admissible windows for the root are set by the characteristic speeds of the ahead
state: a right-going fast shock has $V_s\in(\lambda_A^+,1)$, a right-going slow shock
$V_s\in(v_x,\lambda_A^+)$, and symmetrically on the left.

\begin{codebox}[: shock speeds]
\file{rmhd/shock_speed.py} evaluates the residual
$f(V_s)=\jump{\Ptot}+J^2\jump{h/\rho}$, building the post-shock state with
\file{rmhd/postshock.py} (\code{postshock_pressure}, a translation of the original
Fortran), and searches the windows $(-1,\lambda_A^-)$, $(\lambda_A^+,1)$,
$(\lambda_A^-,v_x)$, $(v_x,\lambda_A^+)$ for the left fast, right fast, left slow and
right slow shock. Because $f$ can have several roots and NaN bands (where the post-shock
state is unphysical), the window is scanned on fixed nodes, refined near its edges, and
the bracketed root is polished by a NaN-tolerant bisection. The solver record (\S5a)
documents a case where a pair of roots closer than one scan step was missed and the fix
(a pair-splitting scan), and that on 119 of 120 such pairs exactly one root was an
evolutionary fast shock.
\end{codebox}

\begin{example}[A relativistic fast shock]\label{ex:fastshock}
Take the ahead state with $\rho=1$, $\Ptot=2$ (gas pressure $p=0.842$),
$\vv=(0.1,\,0.2,\,0.05)$, $\BB=(1.2,\,0.9,\,0.3)$ and $\Gamma=5/3$, and impose
$\Ptot_b=3$ behind a right-going shock. Solving the seven RH conditions for the state behind and $V_s$
(\cref{app:verification}) gives $V_s=0.86778$ and
\[
(\rho,v_x,v_y,v_z,B^y,B^z)_b=(1.2966,\,0.2911,\,0.1262,\,0.0291,\,1.3518,\,0.4430).
\]
The right-going fast speeds are $\lambda_{\rm f}^+=0.8275$ ahead and $0.8943$ behind, and the
right-going Alfvén speed behind is $0.6128$, so
$\lambda_{\rm f}^+(\UU_a)<V_s<\lambda_{\rm f}^+(\UU_b)$ and $V_s>\lambda_{\rm A}^+(\UU_b)$: an
evolutionary $1\to2$ fast shock (\cref{sec:admissibility}). The tangential field grows from
$|B_t|=0.949$ to $1.422$, and the invariants of \cref{prop:invariants} and the Lichnerowicz
adiabat hold to $10^{-39}$.
\end{example}

\paragraph{What a shock does to the tangential field.} In Newtonian MHD a fast shock
increases $|B_t|$ and a slow shock decreases it; the limiting cases are the
\emph{switch-on} shock (fast, with $B_t=0$ ahead) and the \emph{switch-off} shock (slow,
with $B_t=0$ behind). \textcite{Komissarov2003} showed that relativistic switch-on and
switch-off shocks exist as limits of fast and slow shocks, as in the Newtonian
case.\unverified{the content of this statement was checked against the abstract of the
preprint arXiv:astro-ph/0206076 by the same author on limit shocks of relativistic MHD;
the text of the journal paper \citet{Komissarov2003} was not accessed.} A
lab-frame rule ``slow shocks weaken $|B_t|$'' is, however, not reliable in SRMHD when the
flow has a tangential velocity: the solver record (\S5a) documents a slow shock that
\emph{strengthens} the lab-frame tangential field (tangential velocity $0.52$) and a fast
shock that weakens it. A solver that chooses between shock and rarefaction by the sign of
the change of $|B_t|$ inherits this problem (\cref{sec:failure-classes}).

\section{Admissibility}\label{sec:admissibility}

The RH conditions admit discontinuities that are not physical, for instance rarefaction
shocks. Two classical criteria select the admissible ones.

\paragraph{Lax conditions.} A discontinuity is a Lax $k$-shock if the characteristics of
family $k$ run into it from both sides, while the characteristics of the other families
cross it \parencite{Lax1957}. For a right-going shock of family $k$ this means
$\lambda_k(\UU_a)<V_s<\lambda_k(\UU_b)$, $\lambda_{k-1}(\UU_b)<V_s$ and
$V_s<\lambda_{k+1}(\UU_a)$.

\paragraph{Evolutionary condition.} Perturb a discontinuity slightly. The linearised jump
conditions are equations for the amplitudes of the \emph{outgoing} small-amplitude waves
(those carried away from the front on either side) and for the perturbation of the front
speed; the incoming waves are given. The problem has a unique solution only if the number
of unknowns equals the number of equations \parencite{JeffreyTaniuti1964,FalleKomissarov2001}.
For a shock whose two states are coplanar---always the case in Newtonian MHD, where
shocks do not rotate the tangential field---the reflection symmetry $z\to-z$ splits the
linearised problem into two independent parts: the out-of-plane (Alfvénic) part, two
equations ($z$-momentum and $B^z$) with no front-speed perturbation, which needs exactly
\emph{two} outgoing Alfvén waves; and the in-plane part, five equations for the outgoing
magnetosonic and entropy waves plus the front speed, which needs exactly \emph{four}.
Classify each side of a right-going front by where $V_s$ lies among that state's
right-going speeds:
\begin{equation}
\text{1: }V_s>\lambda_{\rm f}^+,\quad
\text{2: }\lambda_{\rm A}^+<V_s<\lambda_{\rm f}^+,\quad
\text{3: }\lambda_{\rm s}^+<V_s<\lambda_{\rm A}^+,\quad
\text{4: }\lambda_{\rm e}<V_s<\lambda_{\rm s}^+ .
\end{equation}
An outgoing characteristic is one faster than $V_s$ on the ahead side or slower than $V_s$
on the behind side. Counting them (\cref{tab:shocks}) shows that only the transitions
$1\to2$ (fast shock) and $3\to4$ (slow shock), ahead state first, have the right numbers.
The \emph{intermediate} shocks $1\to3$, $1\to4$, $2\to3$ and $2\to4$ all have one outgoing
Alfvén wave instead of two and are non-evolutionary in ideal MHD (for $2\to3$ the total
is six, but split as $1+5$). Switch-on and switch-off shocks are limiting cases of fast and
slow shocks.

\begin{table}[t]
\centering\small
\begin{tabular}{@{}llcl@{}}
\toprule
transition (ahead$\to$behind) & name & outgoing: Alfvén + in-plane & ideal-MHD status\\
\midrule
$1\to2$ & fast shock & $2+4$ & evolutionary\\
$3\to4$ & slow shock & $2+4$ & evolutionary\\
$1\to3$ & intermediate & $1+4$ & non-evolutionary\\
$1\to4$ & intermediate & $1+3$ & non-evolutionary\\
$2\to3$ & intermediate & $1+5$ & non-evolutionary\\
$2\to4$ & intermediate & $1+4$ & non-evolutionary\\
\midrule
required & & $2+4$ &\\
\bottomrule
\end{tabular}
\caption{Counting of outgoing characteristics for a right-going coplanar shock. The
rotational discontinuity ($V_s=\lambda_A$, linearly degenerate field) is not in this
table: it is neither compressive nor of Lax type.}
\label{tab:shocks}
\end{table}

\paragraph{The status of intermediate shocks.} Whether intermediate shocks are physical
has been debated for decades. \textcite{Wu1987,Wu1988,Wu1990} showed in dissipative
simulations that they form from smooth waves and possess stable structures.
\textcite{FalleKomissarov2001} argued from the relation between evolutionary conditions and
dissipative shock structure that non-evolutionary shocks are inadmissible, explained the
intermediate shocks seen in planar MHD simulations, and suggested modifications of numerical
schemes to avoid them in symmetric problems. \textcite{InoueInutsuka2007} found that in
dissipative MHD all shock types, including intermediate shocks, are evolutionary. Exact
solvers take different positions: \textcite{GiacomazzoRezzolla2006} exclude compound waves
by construction (each fast or slow wave is a single shock or rarefaction) and do not discuss
intermediate shocks, although the left slow shock of their stored Balsara~1 solution is a
$2\to3$ intermediate shock (solver record \S5), whereas the Newtonian solver of
\textcite{TakahashiYamada2014} includes all types of intermediate shocks and switch-on/off
waves and can find all solutions of a given problem. We return to the consequences in
\cref{ch:nonunique}.

\begin{codebox}[: which shock roots are accepted]
The solver record (\S5a) traced 120 fast waves on which the standard scan had found no
shock: each had a pair of roots, and on 119 of them exactly one root was evolutionary
(two-sided Lax on the fast family \emph{and} super-Alfvénic behind, i.e.\ $1\to2$), the
other leaving the flow sub-Alfvénic behind (an intermediate shock). The Lax condition on
the fast family alone, which is what the scan checks, admits both. The pair-splitting
scan (\code{RMHD_FAST_EDGE_SCAN}, on by default) runs on the lanes the standard scan gave
up on; it splits the pair where the shock speed equals the Alfvén speed behind and takes
the first root outward from there. An evolutionary
selection of slow-shock roots exists but is parked on a branch (\code{RMHD_EVOLUTIONARY},
branch \code{evolutionary-selection}). Every
accepted answer must in addition pass the full RH residual and the wave-order gate
(\cref{sec:student-impl}).
\end{codebox}

\section*{Exercises}
\addcontentsline{toc}{section}{Exercises}

\begin{exercise}
Verify that $T^{\mu\nu}n_\nu=W_s(T^{\mu x}-V_sT^{\mu0})$ and hence that
$\jump{T^{\mu\nu}}n_\nu=0$ is equivalent to the momentum and energy rows of
\eqref{eq:RH}. Derive \eqref{eq:jumpidentity}.
\end{exercise}

\begin{exercise}\label{ex:hugoniot}
Show that \eqref{eq:taub} reduces to the Newtonian Hugoniot relation
$\epsilon_b-\epsilon_a=\frac12(p_a+p_b)(1/\rho_a-1/\rho_b)$ when $\epsilon$ and
$p/\rho$ are small.
\end{exercise}
\begin{solution}
Write $h=1+\epsilon+p/\rho$; to first order in small quantities
$h_a^2-h_b^2\approx2(\epsilon_a-\epsilon_b)+2(p_a/\rho_a-p_b/\rho_b)$ and
$(h_a/\rho_a+h_b/\rho_b)(p_a-p_b)\approx(1/\rho_a+1/\rho_b)(p_a-p_b)$. Rearranging gives the
Hugoniot relation.
\end{solution}

\begin{exercise}
Show that the Lichnerowicz adiabat \eqref{eq:lichnerowicz} reduces to \eqref{eq:taub} for
$\bm b=0$, and that for an Alfvén discontinuity ($\jump{\rho}=\jump{p}=\jump{b^2}=0$) it
is satisfied identically.
\end{exercise}

\begin{exercise}
Starting from \eqref{eq:sigma}, show that $u_\mu\sigma^\mu=-s\sqrt w$ and
$\sigma_\mu\sigma^\mu=b^2-w=-\rho h$ on both sides of a rotational discontinuity, and
explain why these are consistent with the continuity of $\sigma^\mu$.
\end{exercise}
\begin{solution}
$u_\mu\sigma^\mu=u_\mu b^\mu+s\sqrt w\,u_\mu u^\mu=-s\sqrt w$, and
$\sigma_\mu\sigma^\mu=b^2+2s\sqrt w\,u_\mu b^\mu+w\,u_\mu u^\mu=b^2-w$. Both sides have
the same $w$ and $b^2$, as required.
\end{solution}

\begin{exercise}
Reproduce \cref{tab:shocks}. Why is the total count of outgoing characteristics not a
sufficient test for coplanar shocks?
\end{exercise}

\begin{exercise}
In the Newtonian limit, show that \eqref{eq:sigma} gives the Walén relation
$\jump{\vv+s\BB/\sqrt{\rho}}=0$ and $\jump{|B_t|}=0$ for a rotational discontinuity.
\end{exercise}


% ======================================================================== chapter 5
\chapter{Rarefaction waves}\label{ch:rarefactions}

\section{Self-similar simple waves}

A rarefaction is a continuous self-similar solution $\UU=\UU(\xi)$, $\xi=x/t$. Substituting
into $\partial_t\UU+\mathsf A\,\partial_x\UU=0$ with $\partial_t=-(\xi/t)\dd/\dd\xi$ and
$\partial_x=t^{-1}\dd/\dd\xi$ gives
\begin{equation}\label{eq:simplewave}
\bigl(\mathsf A(\UU)-\xi\,\mathsf I\bigr)\,\frac{\dd\UU}{\dd\xi}=0 .
\end{equation}
A non-constant solution therefore needs $\xi=\lambda_k(\UU(\xi))$ for some family $k$, with
$\dd\UU/\dd\xi$ parallel to the right eigenvector $\bm r_k$. Differentiating
$\xi=\lambda_k(\UU(\xi))$ gives $1=\nabla_{\UU}\lambda_k\cdot\dd\UU/\dd\xi$, hence
\begin{equation}\label{eq:fanode}
\boxed{\;\frac{\dd\UU}{\dd\xi}=\frac{\bm r_k(\UU)}{\nabla_{\UU}\lambda_k\cdot\bm r_k(\UU)}\;}
\end{equation}
This requires $\nabla_{\UU}\lambda_k\cdot\bm r_k\ne0$ along the wave: the $k$-th field must
be genuinely nonlinear there. Linearly degenerate fields (entropy, Alfvén) have no
rarefactions; their simple waves are discontinuities moving at the characteristic
speed. A $k$-rarefaction joining a left state $\UU_l$ to a right state $\UU_r$ occupies
$\lambda_k(\UU_l)\le\xi\le\lambda_k(\UU_r)$ and is admissible only if $\lambda_k$
increases monotonically through the fan. The edge facing the undisturbed state is the
\emph{head}, the other the \emph{tail}.

Along the integral curve any function $w(\UU)$ with $\nabla_{\UU}w\cdot\bm r_k=0$ is
constant: a $k$-\emph{Riemann invariant}. An $m$-component system has $m-1$ independent
ones for each family. In practice it is easier to parametrise the curve by a physical
variable, e.g.\ the pressure: $\dd\UU/\dd p=\bm r_k/r_k^{(p)}$, where $r_k^{(p)}$ is the
pressure component of $\bm r_k$ in primitive variables.

\section{Relativistic hydrodynamics}\label{sec:hydro-raref}

For normal flow ($v_y=v_z=0$) take $\VV=(\rho,v,p)$ with $v\equiv v_x$. The eigenvalues of
$\tilde{\mathsf A}$ are
\begin{equation}
\lambda_{\rm e}=v,\qquad \lambda_\pm=\frac{v\pm c_s}{1\pm vc_s},
\end{equation}
the relativistic composition of the flow speed with $\pm c_s$, and along the $\pm$ acoustic
eigenvectors
\begin{equation}\label{eq:dvdp}
\frac{\dd v}{\dd p}=\pm\frac{1}{W^2\rho hc_s},\qquad \dd p=hc_s^2\,\dd\rho
\end{equation}
(both checked from the Jacobian in \cref{app:verification}). With the rapidity
$\vartheta=\operatorname{artanh}v$, $\dd\vartheta=W^2\dd v$, so \eqref{eq:dvdp} becomes
$\dd\vartheta=\pm\dd p/(\rho hc_s)=\pm c_s\,\dd\rho/\rho$. The Riemann invariants are
therefore
\begin{equation}\label{eq:riemanninv}
\mathcal J_\pm=\vartheta\pm\int\frac{c_s}{\rho}\,\dd\rho ,
\end{equation}
with $\mathcal J_+$ constant across a left-going ($\lambda_-$) rarefaction and
$\mathcal J_-$ constant across a right-going ($\lambda_+$) one.

\paragraph{Ideal gas.} On the isentrope $p=K\rho^\Gamma$ one finds
$c_s\,\dd\rho/\rho=2\,\dd c_s/(\Gamma-1-c_s^2)$ (checked symbolically), so
\begin{equation}
\int\frac{c_s}{\rho}\dd\rho=\mathcal G(c_s)\equiv\frac{1}{\sqrt{\Gamma-1}}
\ln\frac{\sqrt{\Gamma-1}+c_s}{\sqrt{\Gamma-1}-c_s} ,
\end{equation}
and across a rarefaction facing the state $a$ (sign $s=-1$ for a left-going wave, $+1$ for a
right-going one)
\begin{equation}\label{eq:raref-hydro}
\operatorname{artanh}v_b=\operatorname{artanh}v_a+s\,\bigl[\mathcal G(c_{s,b})-\mathcal G(c_{s,a})\bigr],
\qquad \rho_b=\rho_a\Bigl(\frac{p_b}{p_a}\Bigr)^{1/\Gamma}.
\end{equation}
Inside the fan the state at $\xi$ follows from $\xi=\lambda_\pm(v,c_s)$, with the sign of $s$, together with
\eqref{eq:raref-hydro}.

\paragraph{Tangential velocities.} With $v_t\ne0$ the tangential momentum equation gives one
more invariant. For the self-similar flow, $\dd(S_t/D)=0$ across the wave, i.e.
\begin{equation}\label{eq:Ainv}
A\equiv hWv_t=\text{const},\qquad\text{hence}\qquad
v_t^2=\frac{A^2(1-v_x^2)}{h^2+A^2},
\end{equation}
and the acoustic ODE becomes \parencite{Pons2000}
\begin{equation}\label{eq:pmm-ode}
\frac{\dd v_x}{\dd p}=s\,\frac{(1-v_x^2)\sqrt{h^2+A^2(1-c_s^2)}}{(h^2+A^2)\,\rho\,c_s},
\end{equation}
which reduces to \eqref{eq:dvdp} for $A=0$. We checked \eqref{eq:pmm-ode} and the
invariance of $A$ against the eigenvectors of the four-variable Jacobian
(\cref{app:verification}). The tangential velocity enters through $h^2+A^2$: unlike in
Newtonian hydrodynamics, the normal and tangential motions are coupled through the
Lorentz factor.

\section{Relativistic MHD}\label{sec:mhd-raref}

For SRMHD the self-similar equations \eqref{eq:simplewave} are seven ODEs for
$(\rho,\Ptot,v_x,v_y,v_z,B^y,B^z)$, written out by \textcite[eqs.~4.43--4.49]{GiacomazzoRezzolla2006}.
Non-trivial solutions need $\xi$ to be a root of the magnetosonic quartic \eqref{eq:quartic}
(their eq.~4.50): fans exist only in the fast and slow families. There is no closed form
for the integral curves, and \textcite{GiacomazzoRezzolla2006} integrate them numerically
with the total pressure as the independent variable (six ODEs; their Appendix~B). Two
structural facts are worth deriving.

\paragraph{The Alfvén singularity.} The right-hand sides of Appendix~B of
\textcite{GiacomazzoRezzolla2006} contain the factor
\begin{equation}
\mathcal R=\frac{1}{\rho hW^4(\vB^2-w)\,(\lambda_A^{(+)}-\xi)(\lambda_A^{(-)}-\xi)}
\end{equation}
and, in the equations for $\dd B^y/\dd\Ptot$ and $\dd B^z/\dd\Ptot$, the denominator
\begin{equation}\label{eq:Bydenom}
(B^x)^2+2B^x\vB W^2(v_x-\xi)+W^4(\vB^2-w)(v_x-\xi)^2
=W^4(\vB^2-w)\,(\lambda_A^{(+)}-\xi)(\lambda_A^{(-)}-\xi).
\end{equation}
The factorisation follows by viewing the left side as a quadratic in $W^2(v_x-\xi)$ with
roots $-B^x/(\vB\pm\sqrt w)$, i.e.\ $\xi=\lambda_A^{(\pm)}$ by \eqref{eq:lamA}; it was
checked symbolically. So the fan equations are singular exactly where the fan speed
$\xi$ (a fast or slow speed) equals an Alfvén speed: at a Type~II degeneracy.
\textcite{GiacomazzoRezzolla2006} note that such points cannot be integrated without a
regularisation, which their code does not include. When the field becomes nearly normal
to the front in the fluid frame ($\theta\to0$), the rest-frame roots of \eqref{eq:ms-rest}
tend to $c_s^2$ and $c_a^2$, so the slow fan approaches this singularity if $c_a<c_s$ and the
fast fan if $c_a>c_s$ (the case of \cref{fig:friedrichs}).

\paragraph{Other parametrisations.} Between the two slow waves GR06 use the tangential field
as the unknown. Their slow fans are integrated with $|B_t|$ as the independent variable at
fixed angle $\psi=\arctan(B^z/B^y)$, and, when the field rotates inside the slow fan (as in
Balsara's test~5), with $\psi$ itself \parencite[\S5.2]{GiacomazzoRezzolla2006}. Note that
in SRMHD $\psi$ need not be constant across fast and slow waves, unlike in Newtonian MHD.

\paragraph{End points and admissibility.} A fan is integrated from the head (the state
it faces, at $\xi=\lambda_k(\UU_a)$) until the independent variable reaches its target; the
tail speed is $\lambda_k$ of the end state. The fan is admissible only if $\lambda_k$ is
monotone along it. If it is not---because the field is not genuinely nonlinear along the
path---the wave must be replaced by a compound structure (\cref{ch:nonunique}).

\begin{codebox}[: fans]
\file{rmhd/rarefaction.py} integrates the fan ODEs with fixed-step fourth-order
Runge--Kutta, with three independent variables: the total pressure (\code{'pressure'}, 800
steps), $|B_t|$ at fixed $\psi$ (\code{'normB'}, 500 steps) and $\psi$ (\code{'psi'}, 1000
steps). The right-hand sides contain the factors $(\lambda_A^{(+)}-\xi)(\lambda_A^{(-)}-\xi)$
discussed above. The batched versions are in \file{rmhd/batched/rarefaction_b.py} and
\file{fan_njit.py}.

Which variable to use matters. The solver record (\S5a) measured, along the slow family at
the input states, $\dd\ln\Ptot/\dd\ln|B_t|\approx-0.005$ (median) on interfaces production
solves, and $-0.54$ and $-0.69$ on unsolved interfaces where the field is
nearly normal to the face: there $|B_t|$ hardly moves along the slow family, so a slow wave
steered by $|B_t|$ is steered by a coordinate that does not change. The four-unknown solver
(\cref{sec:student-impl}) therefore chooses, for each slow wave, between $B_t$ and $\Ptot$
from the tangent of the slow family.
\end{codebox}

\section{The state on the ray \texorpdfstring{$\xi=0$}{xi = 0}}

Godunov's flux needs the state at $\xi=0$. If $\xi=0$ falls inside a fan, the integration
is continued from the head only until $\lambda_k(\UU)=0$. In the code this is
\code{state_at_xi} in \file{rmhd/batched/ray_b.py}, which also handles the case of a
constant zone.

\section*{Exercises}
\addcontentsline{toc}{section}{Exercises}

\begin{exercise}
Show that \eqref{eq:fanode} implies $\dd\lambda_k/\dd\xi=1$ along the fan, and that a
rarefaction across which $\lambda_k$ would have to decrease is not a solution.
\end{exercise}

\begin{exercise}
Derive \eqref{eq:dvdp} directly from the self-similar form of the mass, momentum and
energy equations of relativistic hydrodynamics (normal flow). \emph{Hint:} the mass
equation gives $(v-\xi)\rho'+\rho\,W^2(1-v\xi)\,v'=0$.
\end{exercise}

\begin{exercise}
Gas at rest with sound speed $c_{s,0}$ fills $x>0$; $x<0$ is vacuum. The gas expands to the
left through a right-going rarefaction. Show that the gas--vacuum interface moves with
$v_{\rm vac}=-\tanh\mathcal G(c_{s,0})$. Evaluate it for $\Gamma=5/3$, $p/\rho=1$, and compare
with the Newtonian escape speed $2c_{s,0}/(\Gamma-1)$.
\end{exercise}
\begin{solution}
At the vacuum end $\rho\to0$, so $c_s\to0$ and $\mathcal G\to0$; \eqref{eq:raref-hydro}
gives $|\operatorname{artanh}v_{\rm vac}|=\mathcal G(c_{s,0})$. For $\Gamma=5/3$, $p/\rho=1$:
$h=3.5$, $c_{s,0}^2=5/10.5$, $c_{s,0}\approx0.690$, $\sqrt{\Gamma-1}\approx0.816$, so
$\mathcal G\approx1.225\ln(1.506/0.126)\approx3.04$ and $|v_{\rm vac}|\approx0.995$. The
Newtonian formula would give $2.07$, faster than light.
\end{solution}

\begin{exercise}
Show that $A=hWv_t$ is continuous across a relativistic hydrodynamic shock as well.
\emph{Hint:} use \eqref{eq:jumpidentity} on the tangential momentum.
\end{exercise}

\begin{exercise}
Prove the factorisation \eqref{eq:Bydenom}.
\end{exercise}


% ======================================================================== chapter 6
\chapter{The relativistic hydrodynamic Riemann problem}\label{ch:hydro}

Before the seven-wave problem we solve the three-wave one completely. The method---express
the state behind each outer wave as a function of the pressure behind it, then match at
the contact---is exactly the one generalised to SRMHD in \cref{ch:srmhd}. It is due to
\textcite{MartiMuller1994} for normal flow and to \textcite{Pons2000} for arbitrary
tangential velocities; \textcite{RezzollaZanotti2001,RezzollaZanottiPons2003} added an
a~priori determination of the wave pattern.

\section{Structure of the solution}

The solution consists of a left-going wave (shock or rarefaction), a contact
discontinuity, and a right-going wave:
\begin{equation}
L\;\longrightarrow\;\mathcal W_\leftarrow\;\longrightarrow\;L_\ast\;\longrightarrow\;
\mathcal C\;\longrightarrow\;R_\ast\;\longrightarrow\;\mathcal W_\rightarrow\;
\longrightarrow\;R .
\end{equation}
Across the contact $p$ and $v_x$ are continuous (\cref{prop:contact} with $\bm b=0$), so
the unknowns are $p_\ast$ and $v_\ast$; $\rho$ and the tangential velocity jump. Each
outer wave is a shock if the pressure rises across it ($p_\ast$ larger than the pressure
of the state it faces) and a rarefaction otherwise.

\section{The shock branch}\label{sec:hydro-shock}

Let $a$ be the state the shock faces and $p_b$ the pressure behind it. Using the mass flux
$j\equiv J=\rho_aW_aW_s(V_s-v_{x,a})$ of \eqref{eq:J} (positive for a right-going shock):
\begin{enumerate}[label=(\roman*)]
\item \emph{Post-shock enthalpy.} The Taub adiabat \eqref{eq:taub} with the ideal-gas
  relation $\rho_b=\Gamma p_b/[(\Gamma-1)(h_b-1)]$ is a quadratic for $h_b$,
  \begin{equation}\label{eq:taubquad}
  \Bigl[1+\frac{(\Gamma-1)(p_a-p_b)}{\Gamma p_b}\Bigr]h_b^2
  -\frac{(\Gamma-1)(p_a-p_b)}{\Gamma p_b}\,h_b
  +\frac{h_a(p_a-p_b)}{\rho_a}-h_a^2=0,
  \end{equation}
  whose positive root is the physical one.
\item \emph{Mass flux.} $j^2=(p_b-p_a)/(h_a/\rho_a-h_b/\rho_b)$, from \eqref{eq:taub}.
\item \emph{Shock speed.} Solving $j=\rho_aW_aW_s(V_s-v_{x,a})$ for $V_s$,
  \begin{equation}\label{eq:Vsj}
  V_s=\frac{\rho_a^2W_a^2v_{x,a}+j\sqrt{j^2+\rho_a^2W_a^2(1-v_{x,a}^2)}}{\rho_a^2W_a^2+j^2},
  \end{equation}
  with $j>0$ for a right-going and $j<0$ for a left-going shock.
\item \emph{Post-shock normal velocity.} The momentum condition with
  \eqref{eq:jumpidentity} gives $h_bW_bv_{x,b}=h_aW_av_{x,a}+W_s(p_b-p_a)/j$, and the energy
  and mass conditions give $h_bW_b=h_aW_a+(p_b-p_a)\bigl[W_sv_{x,a}/j+1/(\rho_aW_a)\bigr]$.
  Dividing,
  \begin{equation}\label{eq:vbshock}
  v_{x,b}=\frac{h_aW_av_{x,a}+W_s(p_b-p_a)/j}
  {h_aW_a+(p_b-p_a)\bigl[W_sv_{x,a}/j+1/(\rho_aW_a)\bigr]} .
  \end{equation}
\end{enumerate}
These are the relations of \textcite{Pons2000} (see also eq.~4.13 of \textcite{Romero2005},
who write them for the magnetised case of \cref{sec:romero}). They hold with tangential
velocities too, if $W_a$ includes $v_{t,a}$; the tangential velocity behind follows from
the invariance of $A=hWv_t$ (\cref{ex:Ashock}) and \eqref{eq:Ainv}. A shock built from
(i)--(iv) satisfies all jump conditions to 38 digits (\cref{app:verification}).

\paragraph{Derivation of (iv).} With $S/D=hWv_x$ and $E/D=hW-p/D$, the momentum and energy
jumps read $(j/W_s)\jump{hWv_x}=\jump{p}$ and $(j/W_s)\jump{hW-p/D}=\jump{pv_x}$, and mass
conservation gives $1/D_b=1/D_a+(W_s/j)(v_{x,a}-v_{x,b})$. Substituting the last relation
into the energy jump, the terms in $v_{x,b}$ cancel and $h_bW_b$ follows; the momentum jump
gives $h_bW_bv_{x,b}$.

\section{The rarefaction branch}

Across a rarefaction facing $a$, \eqref{eq:raref-hydro} gives $v_{x,b}$ in closed form for
normal flow. With tangential velocities one integrates \eqref{eq:pmm-ode} from $p_a$ to $p_b$
along the isentrope.

\section{The equation for the central pressure}

Define, for the left and right waves,
\begin{equation}
v_{L\ast}(p)=\begin{cases}\mathcal S_L(p),&p>p_L\\ \mathcal R_L(p),&p\le p_L\end{cases},
\qquad
v_{R\ast}(p)=\begin{cases}\mathcal S_R(p),&p>p_R\\ \mathcal R_R(p),&p\le p_R\end{cases},
\end{equation}
where $\mathcal S$ and $\mathcal R$ denote the shock and rarefaction branches of
\cref{sec:hydro-shock} and \eqref{eq:raref-hydro}. The solution is the root of
\begin{equation}\label{eq:pstar}
\boxed{\;v_{L\ast}(p_\ast)=v_{R\ast}(p_\ast)\;}
\end{equation}
$v_{L\ast}$ decreases and $v_{R\ast}$ increases monotonically with $p$, so for normal flow
\eqref{eq:pstar} has at most one root; it exists unless $v_{L\ast}(0)\le v_{R\ast}(0)$, in which
case the gas separates and a vacuum forms between the two rarefactions (last exercise of this
chapter). When it exists, bisection or a safeguarded Newton iteration finds it reliably
\parencite{MartiMuller1994}. \Cref{fig:hydro-pstar} shows the two curves
for a strong blast wave.

\begin{figure}[t]
\centering
\begin{tikzpicture}
\begin{semilogxaxis}[width=11.5cm,height=6.6cm,xmin=0.01,xmax=1e4,ymin=-0.85,ymax=1.05,
  xlabel={$p$},ylabel={$v_x$},grid=major,grid style={gray!20},
  legend style={at={(0.02,0.04)},anchor=south west,font=\small},
  every axis plot/.append style={line width=1pt}]
\addplot[blue!70!black] table[x index=0,y index=1,row sep=\\]{
0.01 0.999974 0.000000\\
0.0125893 0.999967 0.018018\\
0.0158489 0.999959 0.036874\\
0.0199526 0.999949 0.056646\\
0.0251189 0.999936 0.077448\\
0.0316228 0.999920 0.099422\\
0.0398107 0.999900 0.122729\\
0.0501187 0.999874 0.147549\\
0.0630957 0.999843 0.174065\\
0.0794328 0.999804 0.202456\\
0.1 0.999755 0.232883\\
0.125893 0.999693 0.265476\\
0.158489 0.999616 0.300316\\
0.199526 0.999519 0.337413\\
0.251189 0.999398 0.376690\\
0.316228 0.999247 0.417957\\
0.398107 0.999058 0.460900\\
0.501187 0.998821 0.505075\\
0.630957 0.998524 0.549913\\
0.794328 0.998152 0.594747\\
1 0.997686 0.638849\\
1.25893 0.997104 0.681483\\
1.58489 0.996374 0.721970\\
1.99526 0.995461 0.759732\\
2.51189 0.994319 0.794339\\
3.16228 0.992889 0.825528\\
3.98107 0.991102 0.853197\\
5.01187 0.988867 0.877393\\
6.30957 0.986073 0.898278\\
7.94328 0.982585 0.916099\\
10 0.978232 0.931152\\
12.5893 0.972805 0.943756\\
15.8489 0.966048 0.954229\\
19.9526 0.957647 0.962875\\
25.1189 0.947221 0.969974\\
31.6228 0.934315 0.975774\\
39.8107 0.918383 0.980495\\
50.1187 0.898788 0.984324\\
63.0957 0.874793 0.987421\\
79.4328 0.845571 0.989919\\
100 0.810217 0.991930\\
125.893 0.767786 0.993547\\
158.489 0.717347 0.994844\\
199.526 0.658069 0.995883\\
251.189 0.589330 0.996715\\
316.228 0.510845 0.997380\\
398.107 0.422806 0.997912\\
501.187 0.325987 0.998336\\
630.957 0.221808 0.998675\\
794.328 0.112304 0.998945\\
1000 0.000000 0.999160\\
1258.93 -0.112316 0.999332\\
1584.89 -0.221890 0.999468\\
1995.26 -0.326234 0.999577\\
2511.89 -0.423318 0.999663\\
3162.28 -0.511702 0.999732\\
3981.07 -0.590568 0.999787\\
5011.87 -0.659680 0.999831\\
6309.57 -0.719282 0.999866\\
7943.28 -0.769968 0.999893\\
10000 -0.812557 0.999915\\
};
\addplot[red!70!black] table[x index=0,y index=2,row sep=\\]{
0.01 0.999974 0.000000\\
0.0125893 0.999967 0.018018\\
0.0158489 0.999959 0.036874\\
0.0199526 0.999949 0.056646\\
0.0251189 0.999936 0.077448\\
0.0316228 0.999920 0.099422\\
0.0398107 0.999900 0.122729\\
0.0501187 0.999874 0.147549\\
0.0630957 0.999843 0.174065\\
0.0794328 0.999804 0.202456\\
0.1 0.999755 0.232883\\
0.125893 0.999693 0.265476\\
0.158489 0.999616 0.300316\\
0.199526 0.999519 0.337413\\
0.251189 0.999398 0.376690\\
0.316228 0.999247 0.417957\\
0.398107 0.999058 0.460900\\
0.501187 0.998821 0.505075\\
0.630957 0.998524 0.549913\\
0.794328 0.998152 0.594747\\
1 0.997686 0.638849\\
1.25893 0.997104 0.681483\\
1.58489 0.996374 0.721970\\
1.99526 0.995461 0.759732\\
2.51189 0.994319 0.794339\\
3.16228 0.992889 0.825528\\
3.98107 0.991102 0.853197\\
5.01187 0.988867 0.877393\\
6.30957 0.986073 0.898278\\
7.94328 0.982585 0.916099\\
10 0.978232 0.931152\\
12.5893 0.972805 0.943756\\
15.8489 0.966048 0.954229\\
19.9526 0.957647 0.962875\\
25.1189 0.947221 0.969974\\
31.6228 0.934315 0.975774\\
39.8107 0.918383 0.980495\\
50.1187 0.898788 0.984324\\
63.0957 0.874793 0.987421\\
79.4328 0.845571 0.989919\\
100 0.810217 0.991930\\
125.893 0.767786 0.993547\\
158.489 0.717347 0.994844\\
199.526 0.658069 0.995883\\
251.189 0.589330 0.996715\\
316.228 0.510845 0.997380\\
398.107 0.422806 0.997912\\
501.187 0.325987 0.998336\\
630.957 0.221808 0.998675\\
794.328 0.112304 0.998945\\
1000 0.000000 0.999160\\
1258.93 -0.112316 0.999332\\
1584.89 -0.221890 0.999468\\
1995.26 -0.326234 0.999577\\
2511.89 -0.423318 0.999663\\
3162.28 -0.511702 0.999732\\
3981.07 -0.590568 0.999787\\
5011.87 -0.659680 0.999831\\
6309.57 -0.719282 0.999866\\
7943.28 -0.769968 0.999893\\
10000 -0.812557 0.999915\\
};
\addplot[only marks,mark=*,mark size=2pt] coordinates {(18.597,0.96041)};
\node[font=\small,anchor=south east] at (axis cs:18,0.7) {$(p_\ast,v_\ast)$};
\addplot[only marks,mark=square*,mark size=1.6pt,gray] coordinates {(1000,0) (0.01,0)};
\legend{$v_{L\ast}(p)$ (left wave),$v_{R\ast}(p)$ (right wave)}
\end{semilogxaxis}
\end{tikzpicture}
\caption{The functions $v_{L\ast}(p)$ and $v_{R\ast}(p)$ for the blast wave
$(\rho,p,v_x)_L=(1,1000,0)$, $(\rho,p,v_x)_R=(1,0.01,0)$, $\Gamma=5/3$. Each curve is a
rarefaction branch below the pressure of its state (grey squares) and a shock branch above
it. They cross at $p_\ast=18.597$, $v_\ast=0.96041$.}
\label{fig:hydro-pstar}
\end{figure}

\begin{example}[Blast wave]\label{ex:blast}
For the data of \cref{fig:hydro-pstar} the left wave is a rarefaction and the right wave a
shock. Solving \eqref{eq:pstar} in 30-digit arithmetic with the formulas above gives
$p_\ast=18.5970787$, $v_\ast=0.96040961$, $\rho_{L\ast}=0.0915518$,
$\rho_{R\ast}=10.41558$ and a shock speed $V_s=0.98680$. The exact solver of the reference
implementation (built-in case~3 of \file{rmhd/initial_data.py}, run with $B=0$) returns
$p_\ast=18.5970777$ and $v_\ast=0.96040961$; the relative difference in $p_\ast$,
$5\times10^{-8}$, is consistent with its fixed-step integration of the fan and its
$10^{-8}$ tolerance. The thin dense shell behind the shock ($\rho_{R\ast}/\rho_R\approx10.4$,
more than the Newtonian limit $(\Gamma+1)/(\Gamma-1)=4$) is the hallmark of relativistic
blast waves.
\end{example}

\section{Wave patterns and the relative velocity}

Which of the four patterns (shock--shock, shock--rarefaction, rarefaction--shock,
rarefaction--rarefaction) occurs is decided by where $p_\ast$ lies relative to $p_L$ and
$p_R$, and hence by the velocity jump between the initial states. Because $v_{L\ast}$ and
$v_{R\ast}$ are monotone, the limiting cases are obtained by evaluating them at $p=p_L$ and
$p=p_R$. \textcite{RezzollaZanotti2001} showed that the pattern is determined by the
relativistic relative velocity of the initial states compared with such limiting values,
which removes the initial guesswork, and \textcite{RezzollaZanottiPons2003} extended this to
flows with tangential velocities.

\section{Tangential velocities}

With $v_t\ne0$ the problem is still a two-unknown problem in $(p_\ast,v_\ast)$, but the
tangential velocity couples to the normal one through $W$: $A=hWv_t$ is invariant across
shocks and rarefactions (not across the contact), the shock relations
\eqref{eq:Vsj}--\eqref{eq:vbshock} use the full $W_a$, and the rarefaction branch uses
\eqref{eq:pmm-ode}. Large tangential velocities can change the wave pattern for fixed
thermodynamic states, and the pattern changes continuously as $v_t$ varies
\parencite{Pons2000,RezzollaZanottiPons2003}. The tangential velocity changes its
magnitude across a wave but, since $\jump{S_y/D}=\jump{S_z/D}=0$, not its direction
\parencite[\S5.2]{GiacomazzoRezzolla2006}.

\section*{Exercises}
\addcontentsline{toc}{section}{Exercises}

\begin{exercise}
Derive \eqref{eq:taubquad} and show that its positive root satisfies $h_b>h_a$ when
$p_b>p_a$.
\end{exercise}

\begin{exercise}
Derive \eqref{eq:Vsj}. Show that $|V_s|\to1$ as $|j|\to\infty$ and that $V_s\to v_{x,a}$ as
$j\to0$.
\end{exercise}

\begin{exercise}\label{ex:Ashock}
Using \eqref{eq:jumpidentity}, show that across a hydrodynamic shock
$\jump{hWv_t}=0$, and that the tangential velocity behind is
$v_t^2=A^2(1-v_x^2)/(h^2+A^2)$.
\end{exercise}
\begin{solution}
The tangential momentum jump is $V_s\jump{S_t}=\jump{S_tv_x}$, i.e.\
$\jump{S_t(V_s-v_x)}=0=(j/W_s)\jump{S_t/D}$ and $S_t/D=hWv_t$. Then $A^2=h^2W^2v_t^2$ with
$W^{-2}=1-v_x^2-v_t^2$ gives the formula.
\end{solution}

\begin{exercise}[Programming]
Implement the solver of this chapter (normal flow) and reproduce \cref{ex:blast}. Then
give the right state a tangential velocity $v_t=0.9$ and observe how $p_\ast$ and $v_\ast$
change. Does the wave pattern change?
\end{exercise}

\begin{exercise}
Show that for $p_L=p_R$ and $v_{x,L}>v_{x,R}$ the solution has two shocks, and for
$v_{x,L}<v_{x,R}$ two rarefactions. What happens if $v_{x,R}-v_{x,L}$ is large enough that
$p_\ast\to0$?
\end{exercise}


% ======================================================================== chapter 7
\chapter{The SRMHD Riemann problem}\label{ch:srmhd}

\section{Seven waves, eight states}

For $B^x\ne0$ and generic data the solution of the SRMHD Riemann problem consists of seven
waves: a fast wave, an Alfvén discontinuity and a slow wave moving left, a contact
discontinuity, and a slow wave, an Alfvén discontinuity and a fast wave moving right
(\cref{fig:sevenwave}). They separate eight constant states, $R_1=L$, $R_2,\dots,R_7$,
$R_8=R$, in the numbering of \textcite{GiacomazzoRezzolla2006}. Fast and slow waves are
shocks or rarefactions; the Alfvén waves and the contact are discontinuities.

\begin{figure}[t]
\centering
\begin{tikzpicture}[x=1.0cm,y=1.0cm,>=Stealth]
  \draw[->] (-6,0) -- (6,0) node[right] {$x$};
  \draw[->] (0,-0.2) -- (0,4.6) node[above] {$t$};
  \fill[red!15] (0,0) -- (155:4.4) -- (146:4.4) -- cycle;
  \fill[red!15] (0,0) -- (30:4.4) -- (39:4.4) -- cycle;
  \foreach \ang/\lab/\sty in {155/LF/thick, 130/LA/dashed, 112/LS/thick, 95/CD/dotted,
                              78/RS/thick, 60/RA/dashed, 30/RF/thick}{
    \draw[\sty,red!70!black] (0,0) -- (\ang:4.4);
    \node[font=\small] at (\ang:4.75) {\lab};
  }
  \foreach \ang/\lab in {168/{$R_1=L$}, 139/$R_2$, 121/$R_3$, 103.5/$R_4$, 86.5/$R_5$,
                         69/$R_6$, 46/$R_7$, 14/{$R_8=R$}}{
    \node[font=\small] at (\ang:3.3) {\lab};
  }
\end{tikzpicture}
\caption{The seven-wave Riemann fan of SRMHD (fast: LF, RF; Alfvén: LA, RA; slow: LS, RS;
contact: CD). The shaded wedges indicate that fast and slow waves may be rarefaction fans.
For $B^x=0$ the Alfvén and slow waves collapse onto the contact (three waves); for
coplanar data the Alfvén waves are either absent or reverse the field.}
\label{fig:sevenwave}
\end{figure}

\paragraph{Counting unknowns.} Each wave belongs to a one-parameter family of states
reachable from the state ahead of it (a wave curve). Fixing one parameter per non-contact
wave---six parameters---fixes $R_2,\dots,R_4$ from the left and $R_7,\dots,R_5$ from the
right. The contact then imposes six conditions (\cref{prop:contact}): continuity of $v_x$,
$v_y$, $v_z$, $B^y$, $B^z$ and $\Ptot$. The exact Riemann problem is therefore a system of
six nonlinear equations in six unknowns. Every exact method is a choice of these six
parameters, and the difficulty lies in that choice: it decides how well conditioned the
system is and where it is even defined.

\section{Special cases}\label{sec:bx0}

\subsection{Purely tangential field orthogonal to the flow}\label{sec:romero}

\textcite{Romero2005} solved the case $B^x=0$ with $\vv=(v_x,0,v_z)$ and $\BB=(0,B^y,0)$,
so that $\vv\cdot\BB=0$. Then $b^0=0$ and $b^\mu=(0,0,b,0)$ with $b=B^y/W$, and
\eqref{eq:Tmunu} becomes
\begin{equation}
T^{\mu\nu}=\rho\hat h\,u^\mu u^\nu+\hat p\,g^{\mu\nu}-b^\mu b^\nu,\qquad
\hat h=h+\frac{b^2}{\rho},\qquad \hat p=p+\frac{b^2}{2}.
\end{equation}
The last term has only a $yy$ component, and $\partial_\mu(b^\mu b^\nu)=\partial_y(b\,b^\nu)=0$
in a flow that depends on $(t,x)$ only. The energy--momentum equations are therefore those
of relativistic hydrodynamics with pressure $\hat p$ and enthalpy $\hat h$. The induction
equation is $\partial_tB^y+\partial_x(v_xB^y)=0$, and with mass conservation it says that
$B^y/D=b/\rho$ is carried with the fluid; by \eqref{eq:jumpidentity} it is also continuous
across shocks. So on each side of the contact $\hat B\equiv b/\rho$ is a constant, and the
problem is relativistic hydrodynamics with the effective equation of state
\parencite[Appendix~A]{Romero2005}
\begin{equation}\label{eq:romeroeos}
\hat p=p+\tfrac12\hat B^2\rho^2,\qquad
\hat\epsilon=\epsilon+\tfrac12\hat B^2\rho,\qquad
\hat c_s^2=\frac{hc_s^2+\hat B^2\rho}{\hat h}=c_s^2+v_A^2-c_s^2v_A^2,\quad
v_A^2=\frac{b^2}{\rho\hat h}.
\end{equation}
The last identity (checked symbolically) shows that the effective sound speed is the fast
speed $\zeta$ for propagation perpendicular to the field. The general relations of
\cref{ch:hydro} (\eqref{eq:taub}, \eqref{eq:Vsj}, \eqref{eq:vbshock}, the invariance of
$A=\hat hWv_t$ and \eqref{eq:pmm-ode}) apply with hats, but not the ideal-gas forms
\eqref{eq:taubquad} and \eqref{eq:raref-hydro}: the three waves are fast shocks or rarefactions and a tangential
discontinuity across which $\hat p$ and $v_x$ are continuous. Magnetic pressure raises
$\hat p$ in the initial states and can change the wave pattern for fixed gas states
\parencite{Romero2005}.

\subsection{Normal field zero: three waves}

For $B^x=0$ but $\vv\cdot\BB\ne0$ the reduction above fails, but the fan still has three
waves (Type~I degeneracy): two fast waves and a tangential discontinuity at which only
$v_x$ and $\Ptot$ are continuous (\cref{prop:contact}). \textcite{GiacomazzoRezzolla2006}
solve it like the hydrodynamic problem: the states behind the fast waves are functions of
the total pressure behind them (the $p$-method of \cref{sec:shocks} and
\cref{sec:mhd-raref}), and the single unknown $\Ptot_\ast$ is the root of
$v_{x,L\ast}(\Ptot_\ast)=v_{x,R\ast}(\Ptot_\ast)$.

\subsection{Coplanar data: five waves}\label{sec:coplanar}

If the tangential fields and velocities of $L$ and $R$ lie in one plane containing the
normal---as on every face of a two-dimensional flow---the solution stays in that plane.
A rotational discontinuity can then only leave the field unchanged or reverse it
(\cref{sec:rotational}). The fan has five waves (fast, slow, contact, slow, fast), plus
possibly a reversal at one or both Alfvén speeds, which is a discrete choice rather than a
continuous parameter. With $B_t$ a \emph{signed} scalar, the unknowns are the two fast-wave
strengths and the tangential field at the contact, and the contact conditions are
$\jump{v_x}=\jump{v_t}=\jump{\Ptot}=0$: three equations in three unknowns. This is the
formulation of the planar solver below. Coplanar Riemann problems are also where the
non-uniqueness of \cref{ch:nonunique} lives.

\section{The exact solution of Giacomazzo and Rezzolla}\label{sec:gr06}

\textcite{GiacomazzoRezzolla2006} gave the first exact solution of the general problem.
Its ingredients are:
\begin{itemize}[nosep]
\item \emph{the $p$-method} for the fast waves and the Alfvén discontinuities: the total
  pressure behind the fast wave is the unknown; the post-shock state follows from the jump
  conditions and the shock speed from \eqref{eq:magtaub}; a fast rarefaction is integrated
  in $\Ptot$; across the Alfvén discontinuity $\rho$ and $\Ptot$ are continuous and the
  remaining five jump conditions are solved by Newton's method at $V_s=\lambda_A$;
\item \emph{the $B_t$-method} for the slow waves: the tangential field between the slow
  waves is the unknown; across a slow shock the post-shock state follows from the jump
  conditions with $(B^y,B^z)$ prescribed, the pressure from the invariance of $hB_n$
  \eqref{eq:invhB}, and the speed again from \eqref{eq:magtaub}; a slow rarefaction is
  integrated in $|B_t|$ (or in $\psi$ when the field rotates);
\item \emph{a $4\times4$ Newton--Raphson iteration} on the unknowns $(\Ptot_L,B^y_c,B^z_c,
  \Ptot_R)$---the total pressures behind the two fast waves and the tangential field at the
  contact---with the residuals $\jump{v_x},\jump{v_y},\jump{v_z},\jump{\Ptot}$ at the
  contact. Continuity of $B^y,B^z$ is built in by the $B_t$-method. Using the total pressure
  as the unknown in all three pairs of regions instead would give three unknowns but five
  conditions at the contact; GR06 found that round-off alone then prevents the five
  conditions from being satisfied simultaneously.
\end{itemize}
They report that the residuals change very rapidly near the root, so the iteration
converges only from a sufficiently close initial guess; they take it from an HLLE solution
on about 800 cells. Their test suite---Komissarov's shock tubes, Balsara's five tests
(\cref{tab:balsara}) and a generic Alfvén test---is solved to an accuracy of about
$10^{-10}$ in most cases ($10^{-11}$--$10^{-12}$ for $B^x=0$), but only to $10^{-6}$ for
Komissarov's collision, $10^{-7}$ for Balsara's test~4 and $3\times10^{-4}$ for test~5. Two
cases are instructive: Balsara's test~1 develops a compound wave in numerical solutions,
which the method excludes by construction, although the left slow shock of its own exact
solution is a $2\to3$ intermediate shock, and the test also has a regular solution, with a
$180^\circ$ rotation (\cref{ch:nonunique}); and in test~5 the left
Alfvén and slow waves are so close that the Alfvén jump is resolved less accurately and the
field rotates inside the slow rarefaction.

\begin{example}[Balsara's test~2]\label{ex:balsara2}
\Cref{tab:balsara2} reproduces the exact solution of Balsara's test~2 given by
\textcite[Table~8]{GiacomazzoRezzolla2006}. Both initial states have $v_y=v_z$ (zero) and
$B^y=B^z$, so the data are coplanar in the plane spanned by $\hat{\bm x}$ and
$(\hat{\bm y}+\hat{\bm z})/\sqrt2$; accordingly $R_2=R_3$ and $R_6=R_7$: the Alfvén waves are
silent and the fan has five waves. The left fast and slow waves are rarefactions, the right
slow and fast waves are shocks, and across the contact only $\rho$ jumps. The built-in initial
guess for this case in \file{rmhd/initial_data.py} (case~9),
$(\Ptot_{\rm LF},B^y_c,B^z_c,\Ptot_{\rm RF})=(23.214,\,3.2052,\,3.2052,\,20.721)$, agrees with
the table to the digits printed.
\end{example}

\begin{table}[t]
\centering\small
\begin{tabular}{@{}lrrrrrrr@{}}
\toprule
 & $\rho$ & $\Ptot$ & $v_x$ & $v_y$ & $v_z$ & $B^y$ & $B^z$\\
\midrule
$R_1=L$ & 1.000 & 78.50 & 0.0000 & 0.0000 & 0.0000 & 6.000 & 6.000\\
$R_2$ & 0.4300 & 23.21 & 0.6344 & $-0.09981$ & $-0.09981$ & 3.045 & 3.045\\
$R_3$ & 0.4300 & 23.21 & 0.6344 & $-0.09981$ & $-0.09981$ & 3.045 & 3.045\\
$R_4$ & 0.3830 & 22.84 & 0.6770 & $-0.05566$ & $-0.05566$ & 3.205 & 3.205\\
$R_5$ & 2.828 & 22.84 & 0.6770 & $-0.05566$ & $-0.05566$ & 3.205 & 3.205\\
$R_6$ & 1.582 & 20.72 & 0.4688 & $-0.2538$ & $-0.2538$ & 3.971 & 3.971\\
$R_7$ & 1.582 & 20.72 & 0.4688 & $-0.2538$ & $-0.2538$ & 3.971 & 3.971\\
$R_8=R$ & 1.000 & 13.99 & 0.0000 & 0.0000 & 0.0000 & 0.7000 & 0.7000\\
\bottomrule
\end{tabular}
\caption{Exact solution of Balsara's test~2 ($B^x=5$, $\Gamma=5/3$), from
\textcite[Table~8]{GiacomazzoRezzolla2006}; $\Ptot$ is the total pressure.}
\label{tab:balsara2}
\end{table}

\section{An exact solver used as a flux function}\label{sec:student-impl}

The reference implementation (\file{rmhd_final}) ports the method above to Python, batches
it over many interfaces, and uses it as the flux of a two-dimensional constrained-transport
code (\file{HLLD}). What follows is a summary of its own documentation (the guide, Parts
1, 3 and 4; the solver record, \S\S1--3, 5--5a); the numbers quoted are measurements
recorded there, mostly on a $64^2$ relativistic rotor.

\paragraph{State and frame.} Every face is rotated so that its normal is $x$; a state is the
vector $(\rho,\Ptot,v_x,v_y,v_z,B^y,B^z)$, with $B_n=B^x$ held separately. Faces are first
classified (\cref{sec:degeneracies}).

\paragraph{Seven-wave Newton.} The production formulation has six unknowns and six residuals,
\begin{align}
\bm u_6&=\bigl[\ln\Ptot_{\rm LF},\ \ln|B_{t,\rm CD}|,\ \psi_{\rm CD},\ \ln\Ptot_{\rm RF},\
\varphi_L,\ \varphi_R\bigr],\\
\bm f_6&=\bigl[\jump{v_x},\ \jump{v_y},\ \jump{v_z},\ \jump{\Ptot},\ s_L,\ s_R\bigr],
\end{align}
where $\Ptot_{\rm LF},\Ptot_{\rm RF}$ are the total pressures behind the fast waves,
$(|B_{t,\rm CD}|,\psi_{\rm CD})$ the tangential field at the contact, and
$\varphi_{L,R}$ the rotation angles of the Alfvén waves. Each slow wave is told both
components of the contact field but has one free parameter; the \emph{slacks}
$s_{L,R}$ measure how far it misses. The Jacobian is a finite difference whose six probes
are stacked into the batch; the step comes from a singular value decomposition with small
singular values truncated, followed by a line search. Converged lanes are frozen, and every
lane's arithmetic is bit-for-bit what it would be alone (batch independence).

\paragraph{Seeds.} Because the basin of convergence is narrow, as GR06 already found, the
Newton is seeded by a small neural network that predicts $\bm u_6$ from the two states,
written in a canonical form that removes the symmetries of the problem. Failed lanes are
retried from other networks' seeds and from jittered seeds. The recorded experience is that
convergence is threshold-shaped in the quality of the seed and that more \emph{diverse}
starts help where better single seeds do not (the guide, \S1.9). Part of what looked like a
narrow basin was the Alfvén solve failing inside the iterations: with the closed-form
retry of \cref{sec:rotational} the first seed alone converges on 89\% of the attempted
faces of the $64^2$ rotor, against 50\% without it (solver record \S5e).

\paragraph{The five-wave planar solver.} On a 2D flow every face is coplanar to about
$10^{-10}$ (median of the cross product of the normalised tangential fields, $7.6\times10^{-11}$
on the rotor). There the seven-wave formulation is ill-posed in the directions it does not
need: $\psi_{\rm CD}$ is fixed by geometry, the slow waves are over-determined, and the
Alfvén and slow speeds nearly coincide. The planar solver (\file{rmhd/batched/planar5_b.py})
rotates the frame so that the left tangential field points along $+y$ and solves
\begin{equation}
\bm u_3=[\ln\Ptot_{\rm LF},\ B_{t,\rm CD},\ \ln\Ptot_{\rm RF}],\qquad
\bm f_3=[\jump{v_x},\ \jump{v_t},\ \jump{\Ptot}],
\end{equation}
with $B_{t,\rm CD}$ signed, so that a field reversal through the contact is a sign change
and a reversal at an Alfvén speed is a discrete branch (\code{flip}). Only a small part of
the natural box of unknowns is constructible (about 6--8\%), so the Jacobian falls back to
backward differences where a forward probe leaves the feasible set.

\paragraph{The four-unknown planar solver.} Where the field is nearly normal to the face and
the Alfvén speed exceeds the sound speed, the fast waves carry the tangential field and the
slow waves carry the pressure, so steering a slow wave by $B_t$ fails (\cref{ch:rarefactions}).
The four-unknown solver (\file{rmhd/batched/planar4_b.py}) gives each wave its own strength,
\begin{equation}
\bm u_4=[\ln\Ptot_{\rm LF},\ s_L,\ s_R,\ \ln\Ptot_{\rm RF}],\qquad
\bm f_4=\bigl[\jump{v_x},\ \jump{v_t},\ \jump{B_t}/\sqrt{2\Ptot},\ \jump{\Ptot}/\Ptot\bigr],
\end{equation}
where $s$ is the signed tangential field (kind B) or $\ln\Ptot$ (kind P) behind the slow
wave, chosen per wave from the tangent of the slow family. It starts from slow waves of zero
strength and carries its own certificate: every shock satisfies the full RH conditions to
$10^{-8}$, every fan is an integral of the fan ODE, the contact closes to $10^{-8}$, the
input is planar and the fan is ordered.

\paragraph{What ``exact'' means.} An answer is accepted as exact if and only if the state
behind it satisfies the full seven-wave jump conditions to $10^{-8}$, verified per
interface, \emph{and} the fan is admissible: its waves are ordered in space. Two lessons are
recorded with this definition. First, a zero residual is not a solution: the jump
conditions do not order the waves, and a fan whose Alfvén wave lies on the wrong side of its
slow wave can have a perfect residual, hence the wave-order gate. Second, a residual is not
a state tolerance: the conversion factor is the condition number of the formulation, so a
residual of $10^{-8}$ pins the state to about $10^{-7}$ in the planar form but, on about 2\%
of coplanar faces, only to about $10^{-4}$ in the seven-wave form, whose Jacobian has a
condition number of order $10^4$ on those faces against about $30$ elsewhere. Re-solving
accepted faces at a tolerance of $10^{-12}$ moved the flux by a
median of $10^{-12}$--$10^{-11}$ and at most $9.4\times10^{-7}$ (solver record, \S5b).

\paragraph{The flux pipeline.} A sweep computes the HLLD flux on every face first, then
overwrites the faces the exact solver earns (the guide, Part~4):
\begin{enumerate}[nosep]
\item \emph{gates}: superluminal or non-positive states keep HLLD; faces whose relative jump
  is below $10^{-2}$ are not attempted and take a linearised flux, which agrees with the
  exact flux to second order in the jump (\cref{sec:roe}); faces whose whole fan lies on one
  side keep the HLLD flux, which is then the exact upwind flux;
\item \emph{canonical orientation}: each face is replaced by the canonical member of its
  symmetry orbit (reflection, signed exchanges of the tangential axes, $\BB\to-\BB$, with
  $B_n\ge0$), solved, and mapped back;
\item \emph{seven-wave Newton} from the network seeds, up to three rounds of 40
  iterations, every Alfvén solve inside it retrying a failed rotation from the closed form
  (\cref{sec:rotational}); an answer is accepted if it converged, the state on the ray is
  resolved and physical, and the fan is ordered;
\item \emph{rescues} on the faces it lost: the five-wave planar solver with extra starts,
  then the four-unknown solver (both only on faces planar to $10^{-6}$);
\item \emph{fallback}: a face every rung lost takes the HLLD flux (if the weak-jump gate is
  lowered below $10^{-2}$, a lost face whose jump is below that takes the linearised flux
  instead);
\item the flux is the flux of the state at $\xi=0$.
\end{enumerate}

\begin{table}[t]
\centering\small
\begin{tabular}{@{}lrr@{}}
\toprule
what happens to a face (all $8.0\times10^6$ face evaluations) & without & with\\
\midrule
below the weak-jump gate (linearised flux) & 85.93\% & 85.94\%\\
fan entirely on one side (upwind) & 2.01\% & 2.01\%\\
attempted by the exact solver & 12.06\% & 12.05\%\\
\midrule
of the attempted: seven-wave Newton & 55.5\% & 92.3\%\\
\phantom{of the attempted: }five-wave planar rescue & 38.9\% & 3.4\%\\
\phantom{of the attempted: }four-unknown planar solver & 2.1\% & 1.8\%\\
\phantom{of the attempted: }degenerate class & 1.1\% & 1.1\%\\
\phantom{of the attempted: }HLLD fallback & 2.38\% & 1.46\%\\
\phantom{of the attempted: }\textbf{solved exactly} & \textbf{97.62\%} & \textbf{98.54\%}\\
\bottomrule
\end{tabular}
\caption{Provenance of the fluxes in two complete $64^2$ relativistic-rotor runs with the
production defaults, without and with the closed-form Alfvén retry of
\cref{sec:rotational}, which is now a default (solver record \S5e). The run of \S5a,
``Every face of one run'', made with the defaults of 2026-10-03, had 43.8\% and 50.5\%.}
\label{tab:provenance}
\end{table}

\Cref{tab:provenance} shows where the fluxes of one complete run came from. Two things
stand out. With the closed-form Alfvén retry the seven-wave Newton answers nine faces in
ten. Without it the planar rescue answered between a quarter and a half of the attempted
faces on the three test problems (\cref{sec:tests}); that looked like a consequence of the
coplanarity of 2D flows, but it was mostly the Alfvén solve failing inside the Newton
iterations, since with the retry the same faces are answered by the seven-wave Newton. And
some faces are not answered at all: 1.5\% of the attempted ones on the rotor, 0.9\% on
Orszag--Tang and 4.0\% on the four-quadrant problem. What is known about them is the
subject of the next chapter.

\section*{Exercises}
\addcontentsline{toc}{section}{Exercises}

\begin{exercise}
Verify that the effective enthalpy and internal energy of \eqref{eq:romeroeos} are
consistent, $\hat h=1+\hat\epsilon+\hat p/\rho$, and derive $\hat c_s$.
\end{exercise}
\begin{solution}
$1+\epsilon+\tfrac12\hat B^2\rho+p/\rho+\tfrac12\hat B^2\rho=h+\hat B^2\rho=h+b^2/\rho$.
Along an adiabat of the effective EOS $\hat B$ is constant, so
$\partial\hat p/\partial\rho=hc_s^2+\hat B^2\rho$ and $\hat c_s^2=(hc_s^2+\hat B^2\rho)/\hat h$.
\end{solution}

\begin{exercise}
Count unknowns and conditions in GR06's method and in the six-unknown formulation. Explain
why each slow wave needs a ``slack'' in the latter but not in the former.
\end{exercise}

\begin{exercise}
Let $L$ and $R$ be coplanar, with $v_z=B^z=0$ in both. Show that the reflection $z\to-z$
maps the Riemann problem to itself, and conclude that if the solution is unique it has
$v_z=B^z=0$ everywhere. Which components of $\bm u_6$ are then fixed, and to which values?
\end{exercise}
\begin{solution}
Under $z\to-z$, $v_z\to-v_z$ and $B^z\to-B^z$ while the data are unchanged; a unique
solution must be invariant, so $v_z=B^z=0$. Then $\psi_{\rm CD}\in\{0,\pi\}$ and
$\varphi_L,\varphi_R\in\{0,\pi\}$: three of the six unknowns are discrete, and the
remaining continuous problem is the three-unknown planar one.
\end{solution}

\begin{exercise}[Programming]
Implement the three-wave solver for $B^x=0$, $\vv\cdot\BB=0$ by reusing your hydrodynamic
solver (generalised to this EOS and to tangential velocity) with the effective EOS
\eqref{eq:romeroeos}, and compare with the $p$-method applied
directly to the magnetised states.
\end{exercise}


% ======================================================================== chapter 8
\chapter{Non-regular solutions and non-uniqueness}\label{ch:nonunique}

The construction of \cref{ch:srmhd} assumes that the solution consists of
\emph{regular} waves: Lax shocks of the fast and slow families, rarefactions, rotational
discontinuities and a contact. Following \textcite{Torrilhon2003jcp} we call such a
solution an \emph{r-solution}, and a solution that contains any other wave (a compound wave
or an intermediate shock) a \emph{c-solution}. Because MHD is neither strictly hyperbolic
nor convex (\cref{sec:degeneracies}), c-solutions are not exotic.

\section{Non-convexity and compound waves}

Where $\nabla_{\UU}\lambda_k\cdot\bm r_k$ changes sign along a wave curve, the classical
construction breaks down: a single shock or a single rarefaction cannot connect the two
states admissibly. The admissible structure can then be a \emph{compound wave}, a shock
with a rarefaction of the same family attached to it, the shock moving at the characteristic
speed of the state it is attached to. \textcite{BrioWu1988} found that the MHD equations are
not convex and observed such a slow compound wave in the numerical solution of their
coplanar shock-tube problem, which is now a standard test. Balsara's test~1
(\cref{tab:balsara}) is its relativistic version; \textcite{GiacomazzoRezzolla2006} report
that its numerical solution develops a left-going slow compound wave (a slow shock adjacent
to a slow rarefaction), which their exact solver excludes by construction. In the coplanar
geometry a compound wave reverses the tangential field through zero while compressing the
gas, which in Newtonian MHD no regular wave can do: a slow shock at most switches the field
off, and a reversal at the Alfvén speed is not compressive. For SRMHD, \textcite{Ibanez2015}
showed that the convexity of the magnetosonic fields depends on the magnetic field through a
generalised fundamental derivative.

\section{Coplanar non-uniqueness}

\textcite{Barmin1996} showed that shock-capturing schemes applied to coplanar MHD data
produce a non-evolutionary slow compound wave, and that a small out-of-plane perturbation
of the field makes a rotational discontinuity split off, after which the compound wave
turns into an evolutionary slow shock. \textcite{Torrilhon2003jpp,Torrilhon2003jcp} made
this precise for Newtonian MHD with an exact solver that includes intermediate waves:
\begin{itemize}[nosep]
\item over a wide range of initial data the r-solution exists and is unique;
\item for certain critical data, among them coplanar fields with vanishing transverse
  velocities, a c-solution exists in addition; the critical data form a hypersurface of
  codimension one, so arbitrarily close to them the solution is unique again;
\item in his coplanar test (states with $\rho,p=1$ and $0.2$, $B_n=1$, tangential fields at
  $180^\circ$) the r-solution has a single $180^\circ$ rotation and the c-solution a compound
  wave instead, and the densities behind the contact differ by about 2\%;
\item finite-volume schemes (Roe, HLLE, central) show \emph{pseudo-convergence}: on coarse
  grids they converge towards the c-solution, and only after heavy refinement towards the
  true one; the effect is strongest when the solution is unique but close to a non-unique
  one, and is expected for any diffusive scheme.
\end{itemize}
Balsara's test~1 behaves the same way in SRMHD (\cref{tab:b1family}): besides its compound
wave and the intermediate shock of \textcite{GiacomazzoRezzolla2006} it has an r-solution
with a single $180^\circ$ rotation, whose star-plateau density is 1.1\% above that of the
intermediate-shock answer (solver record \S5).
\textcite{TakahashiYamada2013} built a solver for regular and non-regular solutions and
found, starting from the Brio--Wu data, an initial condition with \emph{no} regular solution,
only a non-regular one on the boundary of the regular ones, and showed that non-regular
solutions are not unique. Their later exact solver handles all types of intermediate shocks
and switch-on/off waves and finds all solutions of a given problem
\parencite{TakahashiYamada2014}.

\section{Selection by dissipation}

If ideal MHD admits two solutions, something outside ideal MHD must select one. In
dissipative MHD the structure of intermediate shocks depends on the dissipation; in Wu's
simulations they form from smooth waves and their structures carry free parameters related
to their stability \parencite{Wu1990}, and \textcite{InoueInutsuka2007} found all MHD shocks
evolutionary once dissipation is included. A numerical scheme supplies its own
dissipation, and therefore its own selection.

\begin{keybox}[Measured: which exact answer a finite-volume code picks]
The solver record (\S5) finds that Balsara's test~1 has not two but a one-parameter family
of exact solutions in SRMHD (\cref{tab:b1family}). Its intermediate shock is not
evolutionary, so the jump conditions leave the state behind it free; with $b$ the
tangential field behind the left field reversal, the family runs from the compound wave,
through the intermediate shock alone (the stored solution of
\textcite{GiacomazzoRezzolla2006}) and an intermediate shock followed by a Lax slow shock,
to the end where the shock has no density or pressure jump and moves at the Alfvén speed on
both sides: a $180^\circ$ rotational discontinuity followed by a Lax slow shock, the
regular solution. Every member satisfies the jump conditions to $10^{-13}$, and the
entropy rises across every shock. One-dimensional runs up to 16384 cells (32768 for HLLE
with the piecewise-linear MC limiter), measured where the members differ, land in different
places of the family. With WENO5-Z reconstruction both HLLD and HLLE approach the compound
wave. With MC neither settles on the intermediate shock: a compressive slow wave grows
behind it, and the tangential field behind the run's shock drifts toward the regular end as
the grid is refined ($b=-0.448$, $-0.458$, $-0.472$ at 8192, 16384 and 32768 cells with
HLLE). Earlier runs at up to 4096 cells, judged by the plateau alone, had seemed to agree on
the intermediate shock. The record's interpretation: the selection happens inside the
discontinuity, where every scheme is first order and the profile is set by its dissipation,
reconstruction included, so a code that reports one member reports a property of its own
solver.
\end{keybox}

\begin{table}[t]
\centering\small
\begin{tabularx}{\textwidth}{@{}lXl@{}}
\toprule
$b$ & the left slow structure & star-plateau $\rho$, relative\\
\midrule
$-0.310$ & intermediate shock, sonic behind, its rarefaction attached (compound wave) & $-1.1\times10^{-3}$\\
$(-0.4285,\,-0.310)$ & intermediate shock and a detached slow rarefaction & \\
$-0.4285$ & the intermediate shock alone \parencite{GiacomazzoRezzolla2006} & $0$\\
$(-0.6414,\,-0.4285)$ & intermediate shock and a Lax slow shock & \\
$-0.6414$ & $180^\circ$ rotational discontinuity and a Lax slow shock (regular) & $+1.10\times10^{-2}$\\
\bottomrule
\end{tabularx}
\caption{The family of exact solutions of Balsara's test~1 (solver record \S5), labelled by
the tangential field $b$ behind the left field reversal. The star-plateau density is
relative to that of the intermediate-shock solution. Beyond $b=-0.6414$ the shock is no
longer intermediate and the entropy falls across it.}
\label{tab:b1family}
\end{table}

\section{What this means for an exact solver}

An exact solver built on the regular seven-wave family cannot represent compound waves or
intermediate shocks. On data whose admissible solution is non-regular it must fail, or it
returns a regular answer if one exists. Two further consequences matter in practice.
\begin{itemize}[nosep]
\item \emph{Inadmissible roots exist.} The jump conditions alone admit fans whose waves are
  in the wrong order; the record (\S5, ``A zero residual is not a solution'') reports that
  a residual-minimising search converged to such inadmissible roots on 11--12\% of
  interfaces where an admissible root was known. An admissibility (wave-order) check is part
  of the definition of a solution.
\item \emph{Uniqueness within the regular family} is an empirical question for the
  populations a code meets. On 7620 lanes compared by two independent routes the record
  found no second admissible regular solution (\S5b). Counting intermediate shocks, answers
  are not unique: when the closed-form Alfvén retry went in, 35 of 35\,000 paired faces
  changed root, 7 of them between a regular and an intermediate-shock answer (\S5e).
\item \emph{The reference solver is not purely regular.} Its slow-shock search extends $0.03$
  past the Alfvén speed (\code{MARGIN}), which admits intermediate shocks: they are 2--4\%
  of the exact answers in its two-dimensional runs, and no diagnostic flags them in
  production. Until the closed-form Alfvén retry it returned the $2\to3$ intermediate shock
  on Balsara~1, as does the stored solution of \textcite{GiacomazzoRezzolla2006}; with it,
  the seven-wave Newton returns the regular solution (solver record \S5e). Which root wins
  when both are reachable is decided implicitly, by the order of the solver's stages.
\end{itemize}

\section{The failure classes of the reference implementation}\label{sec:failure-classes}

The reference implementation keeps a ledger of every interface its exact solver loses.
Three classes were documented (solver record, \S5--5a and the four-quadrant section of
\S7). The closed-form Alfvén retry of \cref{sec:rotational} has since answered many of the
faces in the first and the third, so each is given with what is known now:

\begin{enumerate}
\item \textbf{Compound waves.} Running each lost interface as a one-dimensional shock tube
  shows a reversal of the tangential field through zero fused to a density and magnitude
  jump---a structure outside the seven-wave family. On the $64^2$ rotor, compound
  interfaces were 2.5\% of the attempted ones and 75.7\% of them went unsolved (2.3\% and
  68.1\% at $128^2$). A two-condition rule computed from the states alone,
  \begin{equation}
  \min(|B_{t,L}|,|B_{t,R}|)<0.3\,|B_n|\quad\text{and}\quad|\Delta\psi|>\pi/2,
  \end{equation}
  (the field is nearly normal to the face and its tangential part turns by more than a right
  angle) flags 2.5\% of the attempted interfaces with 81\% precision and 79\% recall. This
  is the regime near $B_t=0$ where the slow, Alfvén and fast speeds approach one another.
  \emph{Revised:} at the resolution of these tube tests a $180^\circ$ rotation riding on a
  slow shock looks the same, and the searches that found no such root built their rotations
  from the rigid seed of \cref{sec:rotational}. With the closed-form retry about half of the
  faces the runs had lost are answered, most of them by regular solutions with a
  $180^\circ$ rotation: 72\% of the 573 rotor faces gained, 90--93\% on Orszag--Tang
  (solver record \S5e).
\item \textbf{Field nearly normal to the face.} About half of what production lost on the
  rotor (695 of 1371 failing interfaces in the ledger) is elementary according to the
  shock-tube test. These interfaces have $|B_n|/|B_t|$ with median 4.4 (0.85 on solved
  ones). Where, in addition, the Alfvén speed exceeds the sound speed, the fast waves carry the
  tangential field and the slow waves carry the pressure, $|B_t|$ is the wrong coordinate for
  the slow wave, and the planar Newton's residual jumps between branches. The four-unknown
  solver, which steers such slow waves by the pressure, certifies 135 of 200 of them; with
  all rescue stages, the exact fraction on the ledger windows rose from 94.39\% to 97.86\%.
\item \textbf{Slow waves riding their Alfvén waves in fast tangential flow.} In a
  four-quadrant problem with a $45^\circ$ field (\cref{sec:tests}) 93.0\% of the attempted
  faces were solved. The lost faces sit just inside the two quadrants moving at $0.99c$,
  in flow with $W\approx6.4$ and magnetisation $0.045$, where the slow waves move with their
  Alfvén waves (gap $\sim2\times10^{-4}$) and the whole inner fan is about $0.01$ wide.
  Nothing in the solver chain closed 95\% of them, and the record concluded that a
  formulation for a slow wave travelling with its Alfvén wave was needed. \emph{Revised:} a
  large part of it was the Alfvén solve. With the closed-form retry the run solves 96.0\% of the
  attempted faces, and of 4000 faces it had lost, 1506 are solved and none lost: 643 by
  regular solutions with a $180^\circ$ rotation and 849 by intermediate shocks, almost all
  from the planar solver (solver record \S5e).
\end{enumerate}
The remaining losses, 1.5\%, 0.9\% and 4.0\% of the attempted faces on the rotor,
Orszag--Tang and the four-quadrant problem at $64^2$, have not been classified since.

\section*{Exercises}
\addcontentsline{toc}{section}{Exercises}

\begin{exercise}
Explain why critical data of codimension one can make a convergent scheme look as if it
converged to the wrong solution (pseudo-convergence), and why refining the grid eventually
removes the effect.
\end{exercise}

\begin{exercise}
Show that the compound-wave rule above is invariant under the symmetries used to
canonicalise a face: rotation of the tangential plane about $x$, $\BB\to-\BB$, and the
reflection $x\to-x$ (which exchanges $L$ and $R$ and flips the normal components).
\end{exercise}
\begin{solution}
$|B_t|$ and $|B_n|$ are invariant under all three; $\Delta\psi$, the angle between the two
tangential fields, is invariant under rotation and under $\BB\to-\BB$, and the exchange
$L\leftrightarrow R$ only changes its sign.
\end{solution}

\begin{exercise}
Balsara's test~1 has a one-parameter family of exact solutions (\cref{tab:b1family}), and
finite-volume codes land on different members depending on their reconstruction. Which
member should an exact Riemann solver used as their flux return: the regular solution, the
intermediate shock of \textcite{GiacomazzoRezzolla2006}, the compound wave, or none? Discuss
in terms of evolutionarity and of consistency between the flux and the dissipation of the
scheme.
\end{exercise}


% ======================================================================== chapter 9
\chapter{Approximate Riemann solvers}\label{ch:approx}

An approximate Riemann solver replaces the exact fan by a simpler model that still
satisfies the integral form of the conservation law. We derive the HLL solver, then the
two relativistic refinements that restore the contact (HLLC) and the Alfvén waves (HLLD),
and finally describe linearised solvers.

\begin{figure}[t]
\centering
\begin{tikzpicture}[x=0.85cm,y=0.85cm]
  \foreach \ox/\title in {0/HLL, 6/HLLC, 12/HLLD}{
    \draw[->] (\ox-2.6,0) -- (\ox+2.6,0);
    \node[font=\small\bfseries] at (\ox,3.2) {\title};
  }
  % HLL
  \draw[thick] (0,0) -- (-2.2,2.6) node[above,font=\scriptsize] {$S_L$};
  \draw[thick] (0,0) -- (2.2,2.6) node[above,font=\scriptsize] {$S_R$};
  \node[font=\scriptsize] at (0,1.6) {$\UU^{\rm hll}$};
  \node[font=\scriptsize] at (-2.1,0.8) {$\UU_L$};
  \node[font=\scriptsize] at (2.1,0.8) {$\UU_R$};
  % HLLC
  \draw[thick] (6,0) -- (3.8,2.6) node[above,font=\scriptsize] {$\lambda_L$};
  \draw[thick] (6,0) -- (8.2,2.6) node[above,font=\scriptsize] {$\lambda_R$};
  \draw[dotted,thick] (6,0) -- (6.5,2.6) node[above,font=\scriptsize] {$\lambda^\ast$};
  \node[font=\scriptsize] at (5.3,1.8) {$\UU_L^\ast$};
  \node[font=\scriptsize] at (7.0,1.8) {$\UU_R^\ast$};
  % HLLD
  \draw[thick] (12,0) -- (9.8,2.6) node[above,font=\scriptsize] {$\lambda_L$};
  \draw[dashed] (12,0) -- (11.0,2.6) node[above,font=\scriptsize] {$\lambda_{aL}$};
  \draw[dotted,thick] (12,0) -- (12.4,2.6) node[above,font=\scriptsize] {$\lambda_c$};
  \draw[dashed] (12,0) -- (13.4,2.6) node[above,font=\scriptsize] {$\lambda_{aR}$};
  \draw[thick] (12,0) -- (14.2,2.6) node[above,font=\scriptsize] {$\lambda_R$};
  \node[font=\tiny] at (10.75,1.25) {$\UU_{aL}$};
  \node[font=\tiny] at (11.75,1.95) {$\UU_{cL}$};
  \node[font=\tiny] at (12.55,1.95) {$\UU_{cR}$};
  \node[font=\tiny] at (13.45,1.25) {$\UU_{aR}$};
\end{tikzpicture}
\caption{The wave models of the HLL, HLLC and HLLD solvers. HLL keeps the two outer waves
and one averaged state; HLLC adds the contact; HLLD adds the two Alfvén (rotational) waves,
with a single total pressure across the whole fan. None of them models the slow waves.}
\label{fig:hllmodels}
\end{figure}

\section{HLL}\label{sec:hll}

Let $S_L\le S_R$ bound all signal speeds of the Riemann fan, i.e.\ the fan lies inside
$S_Lt\le x\le S_Rt$. Apply the integral form \eqref{eq:integralform} to the rectangle
$[S_LT,S_RT]\times[0,T]$: the states on its left and right edges are $\UU_L$ and $\UU_R$,
so
\begin{equation}
\int_{S_LT}^{S_RT}\UU(x,T)\,\dd x=T\bigl(S_R\UU_R-S_L\UU_L+\FF_L-\FF_R\bigr).
\end{equation}
\textcite{HartenLaxvanLeer1983} replace the fan by a single constant state equal to this
average,
\begin{equation}\label{eq:Uhll}
\UU^{\rm hll}=\frac{S_R\UU_R-S_L\UU_L+\FF_L-\FF_R}{S_R-S_L}.
\end{equation}
The flux on the ray $\xi=0$ follows from the integral form on $[S_LT,0]\times[0,T]$, which
gives $\FF^{\rm hll}=\FF_L+S_L(\UU^{\rm hll}-\UU_L)$, or equivalently (using
$[0,S_RT]$) $\FF^{\rm hll}=\FF_R+S_R(\UU^{\rm hll}-\UU_R)$. Eliminating $\UU^{\rm hll}$,
\begin{equation}\label{eq:Fhll}
\hat\FF=\begin{cases}
\FF_L, & S_L\ge0,\\[2pt]
\dfrac{S_R\FF_L-S_L\FF_R+S_LS_R(\UU_R-\UU_L)}{S_R-S_L}, & S_L<0<S_R,\\[8pt]
\FF_R, & S_R\le0 .
\end{cases}
\end{equation}
With $c_{\min}\equiv\max(0,-S_L)$ and $c_{\max}\equiv\max(0,S_R)$ the three cases merge into
the single formula used in \file{hlld.py},
\begin{equation}
\hat\FF=\frac{c_{\max}\FF_L+c_{\min}\FF_R-c_{\max}c_{\min}(\UU_R-\UU_L)}{c_{\max}+c_{\min}}
\end{equation}
(identities checked symbolically). The HLL flux is consistent ($\hat\FF(\UU,\UU)=\FF(\UU)$),
conservative, and robust: with suitable signal speeds \textcite{Einfeldt1988} showed it
preserves positivity (the ``HLLE'' solver). Its weakness is equally clear: by averaging the
fan into one state it smears every intermediate wave---contact, Alfvén and slow---over the
numerical diffusion of the fastest waves.

\paragraph{Signal speeds.} Common choices are the minimum and maximum of the outer
characteristic speeds of the two states \parencite[as in][]{MignoneBodo2006}, from the quartic
\eqref{eq:quartic}, or the isotropic bound \eqref{eq:lamzeta} \parencite{GammieMcKinneyToth2003},
which by \cref{prop:fastbound} can only overestimate them. The code uses the latter.

\section{HLLC for relativistic MHD}\label{sec:hllc}

\textcite{MignoneBodo2006} restore the contact. The fan is modelled by two outer waves
$\lambda_L,\lambda_R$ and a contact $\lambda^\ast$ separating two star states
$\UU^\ast_L,\UU^\ast_R$, related by the RH conditions across each wave and constrained by
the two consistency conditions that the star states average to $\UU^{\rm hll}$ and the star
fluxes to $\FF^{\rm hll}$. The contact moves with the fluid, $\lambda^\ast=v_x^\ast$, and the
total pressure and normal velocity are continuous across it.

\paragraph{The case $B^x=0$.} (Following \textcite{MignoneBodo2006} we write
$\bm m\equiv\SSS$, $E=\tau+D$, and $\gamma$ for the Lorentz factor $W$.) The star states have
$m_x^\ast=(E^\ast+p^\ast)v^\ast$ and
star fluxes $F^\ast_{m_x}=m_x^\ast v^\ast+p^\ast$, $F^\ast_E=m_x^\ast$, with common
$p^\ast,v^\ast$ ($p$ is the total pressure here). Hence on both sides
\begin{equation}
F^\ast_{m_x}-v^\ast F^\ast_E=p^\ast,\qquad m_x^\ast-v^\ast E^\ast=p^\ast v^\ast .
\end{equation}
Both relations are linear in the star quantities. The consistency conditions express
$\UU^{\rm hll}$ and $\FF^{\rm hll}$ as averages of the two star values with weights that sum to
one, so the relations hold for the HLL averages too:
$p^\ast=F^{\rm hll}_{m_x}-v^\ast F^{\rm hll}_E$ and $m_x^{\rm hll}=(E^{\rm hll}+p^\ast)v^\ast$.
Eliminating $p^\ast$:
\begin{equation}\label{eq:hllc-quad}
F^{\rm hll}_E\,(v^\ast)^2-\bigl(E^{\rm hll}+F^{\rm hll}_{m_x}\bigr)v^\ast+m_x^{\rm hll}=0,
\end{equation}
the root with the minus sign being the physical one \parencite[eq.~47]{MignoneBodo2006}.

\paragraph{The case $B^x\ne0$.} Now $\vv$, $\BB$ and the total pressure are continuous at the
contact. The tangential components of the consistency conditions give
\begin{equation}
\BB_t^\ast=\BB_t^{\rm hll},\qquad B^xv_t^\ast=B_t^\ast v_x^\ast-F^{\rm hll}_{B_t},
\end{equation}
and the normal-momentum and energy components give \parencite[eqs.~39--40]{MignoneBodo2006}
\begin{align}
&E^{\rm hll}v_x^\ast+p^\ast v_x^\ast-B^x(\vv^\ast\cdot\BB^\ast)=m_x^{\rm hll},\\
&\bigl[F^{\rm hll}_E-B^x(\vv^\ast\cdot\BB^\ast)\bigr]v_x^\ast-\Bigl(\frac{B^x}{\gamma^\ast}\Bigr)^2
+p^\ast-F^{\rm hll}_{m_x}=0, \label{eq:hllc-p}
\end{align}
with $\gamma^\ast$ the star Lorentz factor. Multiplying the second by $v_x^\ast$ and
subtracting the first gives a quadratic $a(v_x^\ast)^2+bv_x^\ast+c=0$ with
\begin{equation}
a=F^{\rm hll}_E-\BB^{\rm hll}_\perp\cdot\FF^{\rm hll}_{B_\perp},\quad
b=-F^{\rm hll}_{m_x}-E^{\rm hll}+|\BB^{\rm hll}_\perp|^2+|\FF^{\rm hll}_{B_\perp}|^2,\quad
c=m_x^{\rm hll}-\BB^{\rm hll}_\perp\cdot\FF^{\rm hll}_{B_\perp},
\end{equation}
again taking the minus root; $v_t^\ast$ and $p^\ast$ follow, and the conserved star states
from the RH conditions across the outer waves. Two remarks of \textcite{MignoneBodo2006}
deserve emphasis: the formulas for $B^x\ne0$ do not reduce continuously to those for $B^x=0$,
because the contact becomes a tangential discontinuity; and the transverse velocities are
ill-defined as $B^x\to0$, with a possible $O(1/B^x)$ term in the transverse momenta that
vanishes for strictly two-dimensional flows ($B^z=v^z=0$) but not in general.

\section{HLLD for relativistic MHD}\label{sec:hlld}

\textcite{MignoneUglianoBodo2009} extend the multi-state approach of
\textcite{MiyoshiKusano2005} to SRMHD. The model fan has five waves---fast $\lambda_L$,
Alfvén $\lambda_{aL}$, contact $\lambda_c$, Alfvén $\lambda_{aR}$, fast $\lambda_R$---and
four intermediate states $\UU_{aL},\UU_{cL},\UU_{cR},\UU_{aR}$; slow waves are not modelled,
so the total pressure $p$ is the same everywhere in the fan. The single unknown is $p$.

\paragraph{Across the fast waves.} With $\bm R\equiv\lambda\UU-\FF$ evaluated on the outer
state (known; $\lambda=\lambda_L$ or $\lambda_R$), the jump conditions for the state behind
the fast wave read \parencite[eqs.~17--20]{MignoneUglianoBodo2009}
\begin{align}
(\lambda-v^x)D&=R_D,\\
(\lambda-v^x)wW^2v^k+b^k(b^x-\lambda b^0)-p\,\delta_{kx}&=R_{m_k},\\
(\lambda-v^x)wW^2-\lambda p+b^0(b^x-\lambda b^0)&=R_E,\\
(\lambda-v^x)B^k+B^xv^k&=R_{B^k},
\end{align}
with $w$ the total enthalpy and $p$ the total pressure. The last gives
$B^k=(R_{B^k}-B^xv^k)/(\lambda-v^x)$, and combining momentum and energy gives the velocity as
an explicit function of $p$ (their eqs.~23--30; the auxiliary quantities $\tilde A,\dots,\tilde X$
are unrelated to those of \cref{ch:characteristics}):
\begin{align}
v^x&=\frac{B^x(\tilde AB^x+\lambda\tilde C)-(\tilde A+\tilde G)(p+R_{m_x})}{\tilde X},\qquad
v^{t}=\frac{\tilde QR_{m_t}+R_{B^t}\bigl[\tilde C+B^x(\lambda R_{m_x}-R_E)\bigr]}{\tilde X},\\
\tilde A&=R_{m_x}-\lambda R_E+p(1-\lambda^2),\quad \tilde G=|\bm R_{B_t}|^2,\quad
\tilde C=\bm R_{m_t}\cdot\bm R_{B_t},\\
\tilde Q&=-\tilde A-\tilde G+(B^x)^2(1-\lambda^2),\quad
\tilde X=B^x(\tilde A\lambda B^x+\tilde C)-(\tilde A+\tilde G)(\lambda p+R_E),
\end{align}
and the total enthalpy is $w=p+(R_E-\vv\cdot\bm R_m)/(\lambda-v^x)$ (their eq.~31).

\paragraph{Across the Alfvén waves.} The invariant of \cref{sec:rotational} is the key:
$\bm K\equiv\bm\sigma/\sigma^0$, with $\sigma^\mu=\eta u^\mu+b^\mu$ and
$\eta=\pm\sgn(B^x)\sqrt w$ ($+$ on the right, $-$ on the left), is continuous across the
Alfvén wave, and $K^x=\lambda_a$. Expressed through $p$ and the outer state
\parencite[eq.~43]{MignoneUglianoBodo2009},
\begin{equation}
K^k=\frac{R_{m_k}+p\,\delta_{kx}+R_{B^k}\,\eta}{\lambda p+R_E+B^x\eta}.
\end{equation}

\paragraph{At the contact.} $\BB$ is continuous; the consistency condition between the two
Alfvén waves gives
\begin{equation}
\BB_c=\frac{\bigl[\BB(\lambda_a-v^x)+B^x\vv\bigr]_{aR}-\bigl[\BB(\lambda_a-v^x)+B^x\vv\bigr]_{aL}}
{\lambda_{aR}-\lambda_{aL}},
\end{equation}
and inverting the definition of $\bm K$ on each side gives the velocity at the contact,
$\vv_c=\bm K-\BB_c\,(1-K^2)/(\eta-\bm K\cdot\BB_c)$ (their eq.~47). Continuity of the
normal velocity is the equation for $p$:
\begin{equation}\label{eq:hlld-f}
\boxed{\;f(p)\equiv v^x_{cR}(p)-v^x_{cL}(p)
=\Delta K^x-B^x\Bigl[\frac{1-K_R^2}{\eta_R-\bm K_R\cdot\BB_c}-\frac{1-K_L^2}{\eta_L-\bm K_L\cdot\BB_c}\Bigr]=0,\;}
\end{equation}
with $\Delta K^x=K^x_R-K^x_L$. Multiplying out, this is eq.~(48) of
\textcite{MignoneUglianoBodo2009}, $\Delta K^x[1-B^x(Y_R-Y_L)]=0$ with
$Y_S=(1-K_S^2)/(\eta_S\Delta K^x-\bm K_S\cdot\Delta K^x\BB_c)$.

\paragraph{Solution and safeguards.} MUB09 solve \eqref{eq:hlld-f} by the secant method from
the HLL total pressure, or from the solution of the $B^x=0$ limit when
$(B^x)^2/p^{\rm hll}<0.1$. They warn that \eqref{eq:hlld-f} may have several roots and impose
$w_L>p$, $w_R>p$ and the ordering $v^x_{aL}>\lambda_L$, $v^x_{cL}>\lambda_{aL}$,
$v^x_{aR}<\lambda_R$, $v^x_{cR}<\lambda_{aR}$; where these fail, they switch to HLL, which in
their tests happened in usually less than 0.1\% of the zones. In the limit $B^x\to0$ the
equation for $p$ becomes a quadratic and the same algorithm applies.

\begin{codebox}[: HLLD]
\file{HLLD/src/physics/hlld.py} follows MUB09: \code{HLLDComputation} evaluates the
quantities above for a trial $p$ and returns $-f(p)$, i.e.\ the same equation with the
opposite sign. The start is MUB09's rule (from the total pressure of the HLL state, or from
the $B_x=0$ limit), the two starting points are moved apart until $f$ changes sign (at most
five times), and a clamped secant iterates until $|f(p)|<10^{-10}$ (about 3.5 iterations
according to the guide, \S2.3; Newton variants were measured to be slower and less reliable
because the residual has several roots, and a comment in the code records an interface with
seven sign changes of $f$ on $10^{-3}\le p\le10^2$, where the starting point decides the root). Faces fall back to HLLE when the root-find fails, when the wave order is violated, or
when $B_n=0$ on both sides; the last is a difference from MUB09, whose formulation covers
$B^x\to0$, and it matters in 2D because a field initially along $x$ makes every $y$-sweep
degenerate.
\end{codebox}

\section{Roe-type and linearised solvers}\label{sec:roe}

A linearised solver decomposes the jump into the characteristic fields of an averaged
state,
\begin{equation}
\UU_R-\UU_L=\sum_k\alpha_k\bm r_k,\qquad
\hat\FF=\tfrac12(\FF_L+\FF_R)-\tfrac12\sum_k|\lambda_k|\,\alpha_k\bm r_k ,
\end{equation}
which resolves every wave family, including contacts and Alfvén waves, as a discontinuity.
It needs the full eigenstructure, which in SRMHD has no closed-form Roe average, and the
eigenvectors must be renormalised at degeneracies (\cref{sec:degeneracies}).
\textcite{Komissarov1999} built a Godunov-type scheme for RMHD on linear Riemann solvers;
\textcite{Balsara2001} a total-variation-diminishing scheme; and \textcite{Anton2010} a
full-wave-decomposition solver based on their renormalised eigenvectors, which they found
less diffusive than HLL and HLLC and comparable in accuracy to HLLD, but more expensive in one
dimension, with the relative cost improving in multidimensional runs.

\begin{codebox}[: the linearised flux below the weak-jump gate]
\file{HLLD/src/physics/linearised.py} uses a linearisation for faces whose relative jump is
below $10^{-2}$: the Jacobian $(\partial\UU/\partial\VV)^{-1}\partial\FF/\partial\VV$ at the
arithmetic mean state (central differences), its eigen-decomposition, and the state on the
ray $\xi=0$, $\VV^\ast=\VV_L+\sum_{\lambda_k<0}\alpha_k\bm r_k$; the flux is $\FF(\VV^\ast)$.
For a weak jump this agrees with the exact solution to second order in the jump, whereas
HLLD's error is first order there. It is never counted as exact.
\end{codebox}

\section{Trade-offs}

\textcite{MattiaMignone2022} compared five solvers that avoid the full eigenstructure (HLL,
HLLC, HLLD, HLLEM and GFORCE) on standard one- and multidimensional tests. For weak to
moderate magnetisation HLLD was the most accurate, followed by HLLC; GFORCE was less
dissipative than HLL and equally robust at strong magnetisation, and HLLEM did not justify its
cost. MUB09 found that HLL needed up to about twice the resolution of HLLD to match its results
in their multidimensional tests. The exact solver is the end point of this sequence. In the
reference implementation's $64^2$ rotor runs, $L_1(\rho)$ against a $512^2$ HLLD reference at
$t=0.4$ is $0.2881$ with HLLD and $0.2874$ with the exact flux on the attempted faces
(solver record, \S5a), and changing the reconstruction from piecewise-linear to fifth order
moved the rotor some 24 times more than switching on the exact flux (\S7). On the magnetised
four-quadrant problem at $64^2$, by contrast, the exact flux reduced $L_1(\rho)$ against the
$512^2$ reference at $t=0.4$ from $0.1100$ (HLLD) to $0.1019$, about 7\% of the gap
($0.1023$ with the closed-form Alfvén retry, which solves three more per cent of the
attempted faces; solver record \S5e). Which solver is ``best'' is therefore problem- and
resolution-dependent; what the exact solver provides unambiguously is a reference.

\section*{Exercises}
\addcontentsline{toc}{section}{Exercises}

\begin{exercise}
Derive \eqref{eq:Fhll} and show that $\hat\FF$ is continuous when $S_L$ or $S_R$ passes
through zero.
\end{exercise}

\begin{exercise}
Show that, for an isolated stationary contact ($\vv_L=\vv_R=0$, $p_L=p_R$, $\BB=0$,
$\rho_L\ne\rho_R$), the HLL flux has a nonzero mass component while the mass component of the HLLC flux vanishes.
\end{exercise}
\begin{solution}
$F_D=0$ on both sides but $\UU_R-\UU_L\ne0$, so the HLL mass flux is
$S_LS_R(D_R-D_L)/(S_R-S_L)\neq0$ (negative for $D_R>D_L$, since $S_L<0<S_R$). In HLLC, $v^\ast=0$ (the quadratic has $c=m_x^{\rm hll}=0$
and the minus root vanishes), so the star states carry no mass flux.
\end{solution}

\begin{exercise}
Derive \eqref{eq:hllc-quad} from the two consistency conditions, filling in the steps above.
\end{exercise}

\begin{exercise}
Show that \eqref{eq:hlld-f} and MUB09's form $\Delta K^x[1-B^x(Y_R-Y_L)]=0$ are the same
equation, and check that the code's residual is $-f(p)$.
\end{exercise}


% ======================================================================== chapter 10
\chapter{Riemann solvers in multidimensional codes}\label{ch:multid}

\section{The finite-volume update}

In two dimensions the semi-discrete finite-volume scheme for cell averages $\UU_{ij}$ is
\begin{equation}\label{eq:fv2d}
\frac{\dd\UU_{ij}}{\dd t}=-\frac{\hat\FF_{i+1/2,j}-\hat\FF_{i-1/2,j}}{\Delta x}
-\frac{\hat{\bm G}_{i,j+1/2}-\hat{\bm G}_{i,j-1/2}}{\Delta y}\equiv\mathcal L(\UU)_{ij},
\end{equation}
where each face flux is obtained from a one-dimensional Riemann problem in the direction
normal to the face, with the left and right states reconstructed at the face. For a
$y$-face the components are permuted so that $y$ plays the role of $x$; the normal field
$B^y$ is then the constant parameter of that Riemann problem. The scheme is \emph{unsplit}
when both flux differences are evaluated from the same state, as in \eqref{eq:fv2d}.

\begin{codebox}[: sweeps]
In \file{HLLD/src/physics/} a sweep computes all face fluxes in one direction
(\file{sweep.py}). The flux functions take the direction as a parameter and cycle the
tangential components, $t_1=(\code{idir}+1)\bmod3$, $t_2=(\code{idir}+2)\bmod3$; the exact
solver works in its own frame, into which \code{to_solver_frame} rotates every face.
\end{codebox}

\section{Reconstruction}

Godunov's scheme is first order because the data are piecewise constant. Second order in
smooth regions is obtained by piecewise-linear reconstruction with a limited slope. With the
monotonised-central (MC) limiter of \textcite{vanLeer1977},
\begin{gather}
\Delta_i=\operatorname{minmod}\Bigl(2(q_i-q_{i-1}),\ \tfrac12(q_{i+1}-q_{i-1}),\ 2(q_{i+1}-q_i)\Bigr),\\
q^L_{i+1/2}=q_i+\tfrac12\Delta_i,\qquad q^R_{i-1/2}=q_i-\tfrac12\Delta_i,
\end{gather}
where $\operatorname{minmod}$ returns the argument of smallest modulus if all have the same
sign and zero otherwise. Higher orders use weighted essentially non-oscillatory (WENO)
reconstruction, which combines three third-order candidate stencils with nonlinear weights
based on smoothness indicators \parencite{JiangShu1996}; the WENO-Z variant modifies the
weights with a global smoothness indicator to reduce dissipation near critical points
\parencite{Borges2008}. Reconstruction acts on primitive variables. In relativity it is
safer to reconstruct the spatial four-velocity $W\vv$ than $\vv$, since any reconstructed
$W\vv$ gives a subluminal velocity.

\begin{codebox}[: reconstruction]
\file{reconstruction.py} offers \code{pcm}, \code{minmod}, \code{mc} (production), \code{weno5z},
\code{mp5} and \code{mp7}; it reconstructs $W\vv$, and a face whose reconstructed state
violates the density or pressure floor falls back to piecewise-constant data (the guide,
\S2.2).
\end{codebox}

\section{Time integration}

The method of lines integrates $\dd\UU/\dd t=\mathcal L(\UU)$ with a strong-stability-preserving
Runge--Kutta scheme; the third-order scheme of \textcite{ShuOsher1988} is
\begin{align}
\UU^{(1)}&=\UU^n+\Delta t\,\mathcal L(\UU^n),\notag\\
\UU^{(2)}&=\tfrac34\UU^n+\tfrac14\bigl[\UU^{(1)}+\Delta t\,\mathcal L(\UU^{(1)})\bigr],\\
\UU^{n+1}&=\tfrac13\UU^n+\tfrac23\bigl[\UU^{(2)}+\Delta t\,\mathcal L(\UU^{(2)})\bigr].\notag
\end{align}
Each stage needs the primitive variables of every cell (\cref{sec:c2p}), a reconstruction,
and a Riemann solve on every face. The time step obeys a CFL condition with the largest
signal speed; the reference code uses this scheme with Courant number $0.25$ (the guide,
\S2.5).

\section{Constrained transport}\label{sec:ct}

The constraint $\nabla\cdot\BB=0$ is preserved by the continuous equations but not, in
general, by a cell-centred finite-volume update. Constrained transport
\parencite{EvansHawley1988} stores the normal field components on the faces,
$B^x_{i-1/2,j}$ and $B^y_{i,j-1/2}$, and updates them by Stokes' theorem from the electric
field $E_z=v_yB_x-v_xB_y$ (from $\EE=-\vv\times\BB$) at the cell corners:
\begin{align}
\partial_tB^x_{i-1/2,j}&=-\frac{E_{z,i-1/2,j+1/2}-E_{z,i-1/2,j-1/2}}{\Delta y},\\
\partial_tB^y_{i,j-1/2}&=+\frac{E_{z,i+1/2,j-1/2}-E_{z,i-1/2,j-1/2}}{\Delta x}.
\end{align}
The discrete divergence of cell $(i,j)$,
$(B^x_{i+1/2,j}-B^x_{i-1/2,j})/\Delta x+(B^y_{i,j+1/2}-B^y_{i,j-1/2})/\Delta y$, has a time
derivative in which each of the four corner values appears twice with opposite signs, so it
is conserved to round-off whatever the corner values are (\cref{ex:ct}).

\paragraph{The Riemann solver supplies the EMF.} The induction fluxes of \eqref{eq:Fx} are
\begin{equation}
F^x(B^y)=v_xB^y-v_yB^x=-E_z,\qquad F^y(B^x)=v_yB^x-v_xB^y=+E_z ,
\end{equation}
so every face flux carries a face-centred value of $E_z$, upwinded by the Riemann solver in
the direction normal to the face. Constrained transport needs it at the corners, where it
must be upwinded in both directions. The arithmetic average of the four face values
\parencite{BalsaraSpicer1999} lacks the upwinding along the face.
\textcite{GardinerStone2005} derived corner EMFs that add corrections built from the
cell-centred $E_z$, chosen by the sign of the mass flux through the adjacent faces, and that
reduce to the one-dimensional solver's flux for grid-aligned flow.
\textcite{MignoneDelZanna2021} give a systematic construction of upwind constrained-transport
schemes from the information of the Riemann solver and compare the alternatives.

\begin{codebox}[: the corner EMF]
\file{ct.py} implements the Gardiner--Stone corner EMF. For the corner $(i-\frac12,j-\frac12)$,
with face EMFs $E^{y}_{E},E^{y}_{W}$ on the $y$-faces to its east and west and
$E^{x}_{N},E^{x}_{S}$ on the $x$-faces to its north and south, cell-centred values
$E_{NE},E_{NW},E_{SE},E_{SW}$, and upwind signs $s$ of the mass flux on those faces,
\begin{align*}
E_z={}&\tfrac14\bigl(E^y_E+E^y_W+E^x_N+E^x_S\bigr)\\
&+\tfrac18\Bigl[(1-s_S)(E^y_E-E_{SE})+(1+s_S)(E^y_W-E_{SW})\\
&\qquad-(1-s_N)(E_{NE}-E^y_E)-(1+s_N)(E_{NW}-E^y_W)\Bigr]\\
&+\tfrac18\Bigl[(1-s_W)(E^x_N-E_{NW})+(1+s_W)(E^x_S-E_{SW})\\
&\qquad-(1-s_E)(E_{NE}-E^x_N)-(1+s_E)(E_{SE}-E^x_S)\Bigr],
\end{align*}
where $s_N,s_S$ are the signs on the north and south $x$-faces and $s_E,s_W$ on the east and
west $y$-faces. Dropping the bracketed corrections (option \code{upwind=False}) gives the
four-face average; with $s=0$ the corrections are the centred average of the two one-sided
estimates. The code sets $s=0$ in a deadband around zero mass flux: in a static medium the
computed mass flux is round-off with no definite sign, and taking its sign broke the discrete
symmetry of the rotor. The measured
$|\nabla\cdot\BB|$ stays at $\sim10^{-13}$ (the guide, \S2.4).
\end{codebox}

\section{Standard tests}\label{sec:tests}

\paragraph{One-dimensional shock tubes.} \Cref{tab:balsara} lists the five Riemann problems of
\textcite{Balsara2001}, which, with those of \textcite{Komissarov1999} (some of whose
parameters were corrected in \textcite{Komissarov2002tests}), are the standard SRMHD shock
tubes. \textcite{GiacomazzoRezzolla2006} give exact solutions for all five.

\begin{table}[t]
\centering\small
\setlength{\tabcolsep}{4.5pt}
\begin{tabular}{@{}ll rrrrrrrr r@{}}
\toprule
test & state & $\rho$ & $p$ & $v_x$ & $v_y$ & $v_z$ & $B_x$ & $B_y$ & $B_z$ & $t$\\
\midrule
1 ($\Gamma=2$) & L & 1.0 & 1.0 & 0.0 & 0.0 & 0.0 & 0.5 & 1.0 & 0.0 & 0.4\\
              & R & 0.125 & 0.1 & 0.0 & 0.0 & 0.0 & 0.5 & $-1.0$ & 0.0 & \\
2 & L & 1.0 & 30.0 & 0.0 & 0.0 & 0.0 & 5.0 & 6.0 & 6.0 & 0.4\\
  & R & 1.0 & 1.0 & 0.0 & 0.0 & 0.0 & 5.0 & 0.7 & 0.7 & \\
3 & L & 1.0 & 1000.0 & 0.0 & 0.0 & 0.0 & 10.0 & 7.0 & 7.0 & 0.4\\
  & R & 1.0 & 0.1 & 0.0 & 0.0 & 0.0 & 10.0 & 0.7 & 0.7 & \\
4 & L & 1.0 & 0.1 & 0.999 & 0.0 & 0.0 & 10.0 & 7.0 & 7.0 & 0.4\\
  & R & 1.0 & 0.1 & $-0.999$ & 0.0 & 0.0 & 10.0 & $-7.0$ & $-7.0$ & \\
5 & L & 1.08 & 0.95 & 0.4 & 0.3 & 0.2 & 2.0 & 0.3 & 0.3 & 0.55\\
  & R & 1.0 & 1.0 & $-0.45$ & $-0.2$ & 0.2 & 2.0 & $-0.7$ & 0.5 & \\
\bottomrule
\end{tabular}
\caption{The relativistic MHD Riemann problems of \textcite[Table~1]{Balsara2001} ($p$ is the
gas pressure, $t$ the final time). Test~1 uses $\Gamma=2$; the others $\Gamma=1.6666$ (as
printed; $5/3$ in \textcite{GiacomazzoRezzolla2006}). The original runs use a unit domain of
1600 zones with the discontinuity at its centre. Verified against Balsara's table; the same
values are in \textcite[Table~4]{GiacomazzoRezzolla2006} and in
\file{rmhd/initial_data.py} (cases 8--12).}
\label{tab:balsara}
\end{table}

\paragraph{The relativistic rotor.} \textcite{DelZanna2003} adapted the classical MHD rotor to
SRMHD: a disk of radius $0.1$ with $\rho=10$ rotating rigidly with $\omega=9.95$ (rim speed
$0.995$, $W_{\max}\approx10$) inside a static medium with $\rho=1$; uniform $p=1$ and
$\BB=(1,0,0)$; $\Gamma=5/3$; output at $t=0.4$. The field winds up and brakes the rotor,
launching fast and Alfvén waves; the central density is swept into a shell. The problem is
invariant under a rotation by $\pi$, which gives a free measure of numerical error.

\paragraph{The relativistic Orszag--Tang vortex.} The Newtonian vortex of
\textcite{OrszagTang1979} has several relativistic versions.
\textcite{vanderHolst2008} use the doubly periodic domain $[0,2\pi]^2$ with $\rho=1$,
$p=10$, $\Gamma=4/3$, $\vv=(-A\sin y,\,A\sin x)$ with $A=0.99/\sqrt2$ (maximum $W\approx7$)
and $\BB=(-\sin y,\,\sin2x)$. \textcite{BeckwithStone2011} use the unit square with
$\rho=25/(36\pi)$, $p=5/(12\pi)$, $\Gamma=5/3$, velocity amplitude $0.5$ per component
(printed as $V^i=[0.5\sin(2\pi y),0.5\sin(2\pi x)]$; the usual orientation has
$v_x=-0.5\sin2\pi y$), and $A_z=(B_0/4\pi)\cos4\pi x+(B_0/2\pi)\cos2\pi y$ with
$B_0=1/\sqrt{4\pi}$, shown at $t=1$. The reference code's problem combines the flow of the
former (peak speed $0.99$, $\Gamma=4/3$) with the density, pressure and field of the latter,
as its documentation states.

\paragraph{Four-quadrant problems.} \textcite{DelZannaBucciantini2002} proposed a relativistic
two-dimensional Riemann problem on a square divided into quadrants with
$(\rho,v_x,v_y,p)=(0.1,0,0,0.01)$ (NE), $(0.1,0.99,0,1)$ (NW), $(0.5,0,0,1)$ (SW) and
$(0.1,0,0.99,1)$ (SE), evolved to $t=0.4$; it produces two shocks and two shear layers.
The reference implementation adds a uniform in-plane field
$\BB=B_0(\cos45^\circ,\sin45^\circ,0)$, $B_0=0.5$, with $\Gamma=5/3$ on $[-0.5,0.5]^2$; its
documentation records that no published magnetised version was found among the codes it
checked.

\section*{Exercises}
\addcontentsline{toc}{section}{Exercises}

\begin{exercise}\label{ex:ct}
Show that the constrained-transport update preserves the discrete divergence of every cell
exactly, for arbitrary corner values of $E_z$.
\end{exercise}
\begin{solution}
Label the corners of cell $(i,j)$ $NE,NW,SE,SW$. The update gives
$\partial_tB^x_{i+1/2,j}=-(E_{NE}-E_{SE})/\Delta y$, $\partial_tB^x_{i-1/2,j}=-(E_{NW}-E_{SW})/\Delta y$,
$\partial_tB^y_{i,j+1/2}=(E_{NE}-E_{NW})/\Delta x$, $\partial_tB^y_{i,j-1/2}=(E_{SE}-E_{SW})/\Delta x$.
Hence
$\Delta x\Delta y\,\partial_t(\nabla\cdot\BB)_{ij}=-(E_{NE}-E_{SE})+(E_{NW}-E_{SW})
+(E_{NE}-E_{NW})-(E_{SE}-E_{SW})=0$.
\end{solution}

\begin{exercise}
Show that $F^x(B^y)=-E_z$ and $F^y(B^x)=+E_z$ with $\EE=-\vv\times\BB$.
\end{exercise}

\begin{exercise}
For Balsara's test~1, compute the magnetic and gas pressures and the magnetisation
$\sigma=b^2/(\rho h)$ in both initial states, and identify the magnetically dominated side.
\end{exercise}
\begin{solution}
With $\vv=0$, $b^2=B^2=1.25$ on both sides, so $b^2/2=0.625$. For $\Gamma=2$,
$h=1+2p/\rho$: $h_L=3$, $\sigma_L=1.25/3\approx0.42$; $h_R=2.6$, $\rho_Rh_R=0.325$,
$\sigma_R\approx3.8$. The right state is magnetically dominated (gas pressure $0.1$ against
magnetic pressure $0.625$).
\end{solution}

\begin{exercise}
Show that the MC slope equals the centred slope $\frac12(q_{i+1}-q_{i-1})$ where the data are
smooth and monotone and the two one-sided slopes do not differ by more than a factor of
three, and that it vanishes at local extrema.
\end{exercise}


% ======================================================================== appendices
\appendix

\chapter{Algebra of \texorpdfstring{$b^\mu$}{b} and the components of \texorpdfstring{$T^{\mu\nu}$}{T}}\label{app:bmu}

With $u^\mu=W(1,\vv)$, $\vB=\vv\cdot\BB$ and $w=\rho h+b^2$:
\begin{align*}
&b^0=W\vB=\vv\cdot\bm b, &&b^i=\frac{B^i}{W}+W\vB\,v^i, &&b_0=-W\vB,\\
&b^2=\frac{B^2}{W^2}+\vB^2, &&W^2b^2-(b^0)^2=B^2, &&B^i=u^0b^i-u^ib^0=W(b^i-b^0v^i),\\
&u_\mu b^\mu=0, &&{}^\ast\!F^{\mu\nu}=u^\mu b^\nu-u^\nu b^\mu, &&\EE=-\vv\times\BB .
\end{align*}
Components of $T^{\mu\nu}=w\,u^\mu u^\nu+\Ptot g^{\mu\nu}-b^\mu b^\nu$:
\begin{align*}
T^{00}&=wW^2-\Ptot-(b^0)^2=\rho hW^2+B^2-p-\tfrac12b^2=\rho hW^2-p+\tfrac12(|\EE|^2+B^2),\\
T^{0i}&=wW^2v^i-b^0b^i=(\rho hW^2+B^2)v^i-\vB B^i=\rho hW^2v^i+(\EE\times\BB)^i,\\
T^{ij}&=wW^2v^iv^j+\Ptot\delta^{ij}-b^ib^j
=(\rho hW^2+B^2)v^iv^j-\vB(B^iv^j+B^jv^i)-\frac{B^iB^j}{W^2}+\Ptot\delta^{ij}\\
&=S^jv^i+\Ptot\delta^{ij}-\frac{B^ib^j}{W},
\end{align*}
with $|\EE|^2=v^2B^2-\vB^2$. In the fluid rest frame $b^\mu=(0,\BB')$, $T^{00}=e+\frac12b^2$,
$T^{0i}=0$ and $T^{ij}=(p+\frac12b^2)\delta^{ij}-b^ib^j$. Useful combinations:
$h_t=h+b^2/\rho$ (total specific enthalpy), $w=\rho h_t$, $c_a^2=b^2/w$,
$\zeta^2=c_s^2+c_a^2-c_s^2c_a^2$, $\sigma=b^2/(\rho h)=c_a^2/(1-c_a^2)$. All identities in
this appendix were checked symbolically (\cref{app:verification}).

\chapter{Eigenstructure}\label{app:eigen}

\section{Characteristic matrix in covariant variables}

Following \textcite{Anile1989} and \textcite{Anton2010}, take the ten covariant variables
$\tilde\UU=(u^\mu,b^\mu,p,s)$, $s$ the specific entropy. For a characteristic covector
$\phi_\mu$, with $a=u^\mu\phi_\mu$, $\calB=b^\mu\phi_\mu$, $G=\phi^\mu\phi_\mu$ and
$\mathcal E\equiv w$, the characteristic matrix acting on
$(\delta u^\nu,\delta b^\nu,\delta p,\delta s)$ is
\begin{equation}\label{eq:charmat}
\mathcal A^\alpha\phi_\alpha=\begin{pmatrix}
\mathcal E a\,\delta^\mu_\nu & m^\mu{}_\nu & l^\mu & 0^\mu\\
\calB\,\delta^\mu_\nu & -a\,\delta^\mu_\nu & f^\mu & 0^\mu\\
\rho h\,\phi_\nu & 0_\nu & a/c_s^2 & 0\\
0_\nu & 0_\nu & 0 & a\end{pmatrix},
\end{equation}
with
\begin{gather*}
m^\mu{}_\nu=(\phi^\mu+2au^\mu)b_\nu-\calB\,\delta^\mu_\nu,\qquad
l^\mu=\phi^\mu+\frac{(\rho h-b^2/c_s^2)\,a\,u^\mu+\calB\,b^\mu}{\rho h},\\
f^\mu=\frac{1}{\rho h}\Bigl(\frac{a}{c_s^2}b^\mu-\calB u^\mu\Bigr).
\end{gather*}
Its determinant is
\begin{equation}
\det(\mathcal A^\alpha\phi_\alpha)=\mathcal E\,a^2\calA^2\calN_4,\qquad
\calA=\mathcal Ea^2-\calB^2,\qquad
\calN_4=\rho h\Bigl(\frac{1}{c_s^2}-1\Bigr)a^4-\Bigl(\rho h+\frac{b^2}{c_s^2}\Bigr)a^2G+\calB^2G .
\end{equation}

\begin{remark}[A sign in the printed matrix]
In \textcite{Anton2010} the $(b,b)$ block is $-u^\mu\delta^\alpha_\beta$ in the general form
(their eq.~22), i.e.\ $-a\,\delta^\mu_\nu$ after contraction with $\phi$, but $+a\,\delta^\mu_\nu$
in the contracted form (their eq.~32); and $l^\mu$ includes the term $\calB b^\mu/(\rho h)$ in
the text below eq.~32 but not in their eq.~26. We checked numerically, at an arbitrary
state and several values of $\lambda$, that with $-a\,\delta^\mu_\nu$ and with the
$\calB b^\mu/(\rho h)$ term the determinant identity above holds exactly and the eigenvectors
below are exact null vectors; with $+a\,\delta^\mu_\nu$, or without the $\calB b^\mu/(\rho h)$
term, the determinant identity fails and the magnetosonic eigenvectors are not null vectors
(\cref{app:verification}).
\end{remark}

\section{Right eigenvectors}

From \textcite[eqs.~46--50]{Anton2010}, who take them from \textcite{Anile1989}:
\begin{align}
\bm r_e&=(0^\alpha,0^\alpha,0,1)^{\mathsf T}, \\
\bm r_{a,\pm}&=\bigl(a\,\varepsilon^\alpha,\ \calB\,\varepsilon^\alpha,\ 0,\ 0\bigr)^{\mathsf T},
\qquad \varepsilon^\alpha=\varepsilon^\alpha{}_{\beta\gamma\delta}u^\beta\phi^\gamma b^\delta, \\
\bm r_{m,\pm}&=\bigl(a\,d^\alpha,\ \calB\,d^\alpha+a\calA f^\alpha,\ a^2\calA,\ 0\bigr)^{\mathsf T},
\end{align}
for the entropy, Alfvén and magnetosonic ($m=$ s, f) waves, with
\begin{equation}
d^\alpha=\frac{c_s^2-1}{c_s^2}\Bigl(\frac{a^4}{G}(\phi^\alpha+au^\alpha)-\frac{a^2\calB}{\rho h}b^\alpha\Bigr),
\qquad
f^\alpha=\frac{1}{\rho h}\Bigl(\frac{a}{c_s^2}b^\alpha-\calB u^\alpha\Bigr),
\end{equation}
each evaluated at the corresponding root ($a=0$; $\calA=0$; $\calN_4=0$). At the Alfvén roots
$\calB/a=\mp\sqrt{\mathcal E}$.

\paragraph{Renormalisation.} The Alfvén vectors vanish at Type~I degeneracies ($a=\calB=0$)
and at Type~II degeneracies ($\varepsilon^\alpha=0$ because $b^\alpha$ becomes a combination
of $u^\alpha$ and $\phi^\alpha$). Dividing by $a$ and using $\calB/a=\mp\sqrt{\mathcal E}$ removes
the first problem; writing $\varepsilon^\alpha=g_1\alpha^\alpha_1+g_2\alpha^\alpha_2$ in a
basis that stays regular, replacing $g_{1,2}$ by $\hat g_{1,2}\equiv g_{1,2}/\sqrt{g_1^2+g_2^2}$
and setting $\hat g_1=\hat g_2=1/\sqrt2$ at the degenerate point removes the second
\parencite[eqs.~52--67]{Anton2010}.
This is the relativistic version of the procedure of \textcite{BrioWu1988}. The magnetosonic
vectors are renormalised similarly \parencite[\S3.3.2]{Anton2010}; their expressions, the left
eigenvectors, and the transformation to conserved variables are long and are not reproduced
here.

\chapter{Notation}\label{app:notation}

\begin{center}
\small
\begin{tabular}{@{}ll@{}}
\toprule
symbol & meaning\\
\midrule
$\rho,\ \epsilon,\ p,\ h$ & rest-mass density, specific internal energy, gas pressure, gas specific enthalpy\\
$\Gamma$ & adiabatic index of the ideal gas\\
$\vv,\ W,\ u^\mu$ & three-velocity, Lorentz factor, four-velocity\\
$\BB,\ \EE$ & laboratory magnetic and electric fields ($\sqrt{4\pi}$ absorbed)\\
$b^\mu,\ b^2$ & comoving magnetic field four-vector and its square\\
$\vB$ & $\vv\cdot\BB$\\
$w,\ h_t$ & total enthalpy density $\rho h+b^2$, total specific enthalpy $w/\rho$\\
$\Ptot$ & total pressure $p+b^2/2$\\
$\sigma$ & magnetisation $b^2/(\rho h)$\\
$D,\ \SSS,\ \tau,\ E$ & conserved density, momentum, energy minus rest mass, total energy\\
$\bm m$ & the momentum $\SSS$ in the notation of \textcite{MignoneBodo2006} (\cref{ch:approx})\\
$c_s,\ c_a,\ \zeta$ & sound speed, rest-frame Alfvén speed $b/\sqrt w$, $\sqrt{c_s^2+c_a^2-c_s^2c_a^2}$\\
$\lambda_{\rm e},\lambda_{\rm A}^\pm,\lambda_{\rm s}^\pm,\lambda_{\rm f}^\pm$ & entropy, Alfvén, slow, fast speeds\\
$a,\ \calB,\ G$ & $u^\mu\phi_\mu$, $b^\mu\phi_\mu$, $\phi^\mu\phi_\mu$ for $\phi_\mu=(-\lambda,1,0,0)$\\
$V_s,\ W_s,\ n_\mu$ & front speed, its Lorentz factor, unit normal $W_s(-V_s,1,0,0)$\\
$J,\ B_n$ & mass flux $-\rho u^\mu n_\mu$, normal comoving field $b^\mu n_\mu$\\
$H$ & shock invariant $B_n^2/J^2-b^2/\rho^2$\\
$\sigma^\mu,\ \bm K$ & Alfvén invariant $b^\mu+s\sqrt w\,u^\mu$, its spatial ratio $\bm\sigma/\sigma^0$\\
$\jump{Q},\ \avg{Q}$ & $Q_a-Q_b$ and $Q_a+Q_b$ (ahead, behind)\\
$\xi$ & similarity variable $x/t$\\
$R_1,\dots,R_8$ & the states of the seven-wave fan\\
$\UU^{\rm hll},\ \FF^{\rm hll}$ & HLL average state and flux\\
\bottomrule
\end{tabular}
\end{center}

\chapter{What was verified, and how}\label{app:verification}

The script \file{sympy_checks.py} in the directory of these notes runs every check below and
prints one line per check (about ten seconds on a laptop). Symbolic checks use \textsc{SymPy};
numerical ones use \textsc{mpmath} at 40 significant digits.

\paragraph{Checked symbolically.}
\begin{itemize}[nosep]
\item \Cref{prop:bmu}: $b^\mu={}^\ast\!F^{\mu\nu}u_\nu$ with $\EE=-\vv\times\BB$ gives
  \eqref{eq:b0}--\eqref{eq:b2}, $u_\mu b^\mu=0$ and ${}^\ast\!F^{\mu\nu}=u^\mu b^\nu-u^\nu b^\mu$.
\item The identities and the components of $T^{\mu\nu}$ in \cref{app:bmu}, including the
  general and rest-frame forms and $\sigma=c_a^2/(1-c_a^2)$.
\item The conserved variables and fluxes \eqref{eq:UF}--\eqref{eq:Fx} from the components of
  \eqref{eq:Tmunu}, both forms of $\tau$ used in \file{hlld.py}, and the momentum-flux form of
  \textcite{MignoneBodo2006}.
\item The Newtonian limits of \cref{sec:newtonian}.
\item The rest-frame dispersion relation: the determinant of the linearised system and of
  the $2\times2$ in-plane system; the covariantisation leading to \eqref{eq:quartic}.
\item The identity of \eqref{eq:quartic} with the coefficients coded in
  \file{rmhd/wave_speeds.py}, with eq.~(16) of \textcite{DelZanna2003}, and its reductions for
  $B^x=0$ and $\vv=0$.
\item The Alfvén speeds \eqref{eq:lamA} and $\calA=0$; the polynomial identities behind
  \cref{prop:ordering}; the isotropic bound \eqref{eq:lamzeta} and \cref{lem:monotone}.
\item The HLL identities; the effective sound speed \eqref{eq:romeroeos}; the relativistic
  Riemann invariant for the ideal gas; the primitive-recovery relations
  \eqref{eq:c2p-v2}--\eqref{eq:c2p-f}, checked exactly, in rational arithmetic, at one
  state; the factorisation \eqref{eq:Bydenom}.
\end{itemize}

\paragraph{Checked numerically.}
\begin{itemize}[nosep]
\item The eigenvalues and acoustic eigenvectors of one-dimensional relativistic hydrodynamics,
  \eqref{eq:dvdp}, and with tangential velocity \eqref{eq:pmm-ode} and the invariance of
  $A=hWv_t$.
\item Hydrodynamic shocks built from \eqref{eq:taubquad}--\eqref{eq:vbshock} satisfy all
  jump conditions to $10^{-38}$; conversely, a shock solved from the jump conditions alone
  satisfies the Taub adiabat \eqref{eq:taub}.
\item A fast MHD shock solved from the full RH system satisfies \eqref{eq:invH}--\eqref{eq:magtaub}
  and the Lichnerowicz adiabat \eqref{eq:lichnerowicz} to $10^{-39}$, and is an evolutionary
  $1\to2$ transition.
\item Rotational discontinuities built from \cref{prop:alfven} satisfy all RH conditions to
  $10^{-40}$ at $V_s=\lambda_A$, with $\rho$, $p$ and $b^2$ continuous and $v_x$, $|B_t|$ not.
\item The characteristic determinant and the eigenvectors of \cref{app:eigen}, and that both
  variants of the matrix printed by \textcite{Anton2010} fail these tests.
\end{itemize}

\paragraph{Checked against the sources.} Every bibliography entry was checked on its arXiv
page or through the Crossref record of its DOI (or a publisher or library record for books).
Formulas quoted from papers were compared with their text: the quartic and HLLC relations
with \textcite{MignoneBodo2006}; the HLLD relations with \textcite{MignoneUglianoBodo2009};
the shock relations, the Lichnerowicz adiabat, the rarefaction ODEs and the test problems
with \textcite{GiacomazzoRezzolla2006}; the effective EOS with \textcite{Romero2005}; the
eigenstructure with \textcite{Anton2010}; \cref{tab:balsara} with Table~1 of
\textcite{Balsara2001}; the rotor with \textcite{DelZanna2003}; the Orszag--Tang set-ups with
\textcite{vanderHolst2008} and \textcite{BeckwithStone2011}; the four-quadrant data with
\textcite{DelZannaBucciantini2002}. Statements about the reference implementation were
compared with its source and with \file{HLLD/docs/GUIDE.md} and
\file{HLLD/docs/what_we_solve.md}.

\paragraph{Not verified.} (i) The form of the quartic in \textcite{Leismann2005}, whose full
text could not be accessed. (ii) The content of \textcite{Komissarov2003} beyond its title,
taken from the abstract of the author's preprint on the same subject. (iii) The derivation of
the Lichnerowicz adiabat, which was only checked numerically. (iv) The page ranges of
\textcite{Thorne1973} and \textcite{RyuJones1995}, for which only the first page is given.
(v) Claims attributed to \textcite{Komissarov1997} are quoted from the summaries in
\textcite{MignoneBodo2006} and \textcite{MignoneUglianoBodo2009}; the paper itself was not
accessed.

\paragraph{Checked separately (not by the script).} \Cref{ex:blast} was compared against
the exact solver of the reference implementation (its built-in case~3).

%%NEXT-CHAPTER%%

\backmatter
\bibliographystyle{plainnat}
\bibliography{refs}

\end{document}
